\documentclass{article}
\usepackage{graphicx} % Required for inserting images
\usepackage[a4paper, left=2cm, right=2cm, top=3cm, bottom=3cm]{geometry}
\usepackage{amsmath}
\usepackage{amssymb}
\usepackage{bm}
\usepackage{hyperref}
\allowdisplaybreaks

\title{Mathematical Foundations Notes}
\author{Siheon Lee}
\date{October 2025}

\begin{document}

\maketitle
\tableofcontents
\newpage

\section{Lecture 1 (07/10/2025)}
\subsection{Slide 1 (Foundations of Mathematics: Set Theory and Logic)}
\begin{itemize}
    \item Ordering types: partial ordering and total ordering
    \item $a \preceq b$ means "a is ranked less than or equal to b"
    \item $a \prec b$ means "a is ranked strictly below b"
    \item Partial Ordering
    \begin{itemize}
        \item Relation $\preceq$ on a set $A$ is called a partial ordering if it satisfies the following properties $\forall a,b,c \in A$:
        \begin{itemize}
            \item Reflexivity: $a \preceq a$
            \begin{itemize}
                \item $a=a$ is always true $\forall a \in A$ regardless of whether or not any order relation is defined on $A$
                \item $a \preceq a$ is putting itself in a relation with itself
                \item Let's define $\preceq \equiv <$. Then reflexivity is not satisfied since $1 \nless 1$ (strict poset).
            \end{itemize}
            \item Antisymmetry: If $a \preceq b$ and $b \preceq a$, then $a=b$
            \item Transitivity: If $a \preceq b$ and $b \preceq c$, then $a \preceq c$
        \end{itemize}
        \item Partially ordered set is often denoted as a poset
    \end{itemize}
    \item Example Poset I
    \begin{itemize}
        \item $S=\{1,2\}$
        \item The power set (set of all subsets) of $S$ is $P(S) = \{ \emptyset, \{1\},\{2\},\{1,2\}\}$
        \item The partial order denoted by $\subseteq$ is the subset relation
        \item For subsets of $A$ and $B$ of $S$, $A \subseteq B$ if every element of $A$ is also an element of $B$
        \item Poset example:
        \begin{itemize}
            \item $\emptyset \subseteq \emptyset$
            \item $\emptyset \subseteq \{1\}$
            \item $\emptyset \subseteq \{2\}$
            \item $\emptyset \subseteq \{1,2\}$
            \item $\{1\} \subseteq \{1\}$
            \item $\{1\} \subseteq \{1,2\}$
            \item $\{2\} \subseteq \{2\}$
            \item $\{2\} \subseteq \{1,2\}$
            \item $\{1,2\} \subseteq \{1,2\}$
            \item The relation is reflexive, antisymmetric, and transitive
        \end{itemize}
    \end{itemize}
    \item Example: Divisibility
    \begin{itemize}
        \item Define the set $S = \{1,2,3,6\}$. For all $a,b \in S$ we say $a \preceq b$ if $a$ divides $b$. Note that $a \preceq b \Longleftrightarrow (\exists k \in \mathbb{Z} - \{0\}) \ b = ka$
        \item $1 \preceq 1$
        \item $1 \preceq 2$
        \item $1 \preceq 3$
        \item $1 \preceq 6$
        \item $2 \preceq 2$
        \item $2 \preceq 6$
        \item $3 \preceq 3$
        \item $3 \preceq 6$
        \item $6 \preceq 6$
        \item We can see that for all $a \in S$, $a \preceq a$ (reflexivity)
        \item Also, $a \preceq b \Longleftrightarrow (\exists k \in \mathbb{Z} - \{0\}) \ b = ka$/ If, in addition, $b \preceq a$, then $k, \frac{1}{k} \in \mathbb{Z} - \{0\} \Rightarrow k=1$. Therefore, $a=b$ (antisymmetry).
        \item Suppose $a \preceq b$ and $b \preceq c$. Then $(\exists k \in \mathbb{Z} - \{0\}) \ b = ka$ and $(\exists l \in \mathbb{Z} - \{0\}) \ c = lb$. Therefore, $c = (lk)a$ and $lk \in \mathbb{Z} - \{0\}$, and so $(\exists m \in \mathbb{Z} - \{0\}) \ c = ma$. Hence $a \preceq c$ (transitivity).
        \item Therefore, the set $S$ with the relation $a \preceq b$ if $a$ divides $b$ is a poset.
    \end{itemize}
    \item Hasse Diagram - way of representing a finite partially ordered set, in the form of a drawing of its transitive reduction
    \item Hasse diagrams only need $1$ to draw lines to prime numbers when the rule is divisibility
    \item Total Ordering (Linear Ordering)
    \begin{itemize}
        \item Relation $\preceq$ on set $A$ if it is partial ordering and satisfies the property of comparability
        \begin{itemize}
            \item $\forall a,b \in A, a \preceq b \vee b \preceq a$
        \end{itemize}
        \item In a totally ordered set, any two distinct elements are comparable
    \end{itemize}
    \item Maximal Element - element $m$ of poset $P$ if there is no element $x \in P$ such that $m \prec x$
    \begin{itemize}
        \item $\forall x \in P, m=x \vee m \npreceq x$
        \item Poset can have multiple maximal elements or none at all
        \item Maximal elements may not necessarily surpass every other element
    \end{itemize}
    \item Consider $P=\{\emptyset, \{1\}, \{2\}, \{3\}, \{1,2\}\}$ ordered by $\subseteq$
    \begin{itemize}
        \item Maximal elements - $\{1,2\},\{3\}$
        \item No element is strictly greater than $\{1,2\}$ or $\{3\}$
    \end{itemize}
    \item Example: No Maximal Elements
    \begin{itemize}
        \item $P=\mathbb{Z}^+ = \{1,2, \dots\}$ ordered b $\le$
        \item $m \in \mathbb{Z}^+$ is maximal
        \item $m < m+1$ and $m+1 \in \mathbb{Z}^+$, giving a contradiction
    \end{itemize}
    \item Theorem - Every non-empty finite poset has at least one maximal element
    \item Greatest Element - $g \in S : \forall s \in S, s \preceq g$
    \begin{itemize}
        \item There can be only one greatest element
    \end{itemize}
    \item Theorem - If $\preceq$ is a total ordering on $A$, then every non-empty finite subset of $A$ has a greatest element
\subsection{Lecture Notes}
    \item Negation - $\neg A$ is statement "A is not valid"
    \begin{itemize}
        \item $x \neq y \equiv \neg(x=y)$
    \end{itemize} 
    \item Conjunction - $A \wedge B$ is statement "both A and B hold"
    \item Disjunction - $A \vee B$ is statement "either A holds or B holds"
    \item $\neg(A \vee B) = \neg A \wedge \neg B$
    \item Implies - $A \Rightarrow B$ is statement "If A, then B"
    \begin{itemize}
        \item $A(t),B(t) \rightarrow \{ A \Rightarrow B \} (t)$
        \item $A(t),B(f) \rightarrow \{ A \Rightarrow B \} (f)$
        \item $A(f),B(t) \rightarrow \{ A \Rightarrow B \} (t)$
        \begin{itemize}
            \item $A$ might not be the only thing affecting $B$
        \end{itemize}
        \item $A(f),B(f) \rightarrow \{ A \Rightarrow B \} (t)$
    \end{itemize}
    \item $(A \Rightarrow B) \equiv \neg A \vee B$
    \item Iff - $A \Longleftrightarrow B$ is statement "A if and only if B"
    \item Inverse - $\neg A \Rightarrow \neg B$
    \item Converse - $B \Rightarrow A$
    \item Contrapositive - $\neg B \Rightarrow \neg A$
    \item Negation - $\neg(A \Rightarrow B) \equiv A \nRightarrow B$
    \item Contrapositive has the same truth table as the original statement $A \Rightarrow B \equiv \neg B \Rightarrow \neg A$
    \begin{itemize}
        \item $(A \Rightarrow B) \equiv \neg A \vee B$
        \item Look at truth table and compare $\neg A$ and $B$
    \end{itemize}
    \item Statement: $(\forall a,b \in \mathbb{Z})(a+b \ge 15 \Rightarrow a \ge 8 \vee b \ge 8)$
    \begin{itemize}
        \item Contrapositive: $(\forall a,b \in \mathbb{Z})(a < 8 \wedge b < 8 \Rightarrow a+b < 15)$
        \item Proof: $a < 8 \wedge b < 8 \Rightarrow a \le 7 \wedge b \le 7 \Rightarrow a+b \le 14 < 15$
    \end{itemize}
    \item $p \Rightarrow (q \wedge \neg q) \equiv \neg p$
    \item Basic Quantifiers
    \begin{itemize}
        \item Exists - $\exists$ is statement "there exists"
        \item For All - $\forall$ is statement "for all"
    \end{itemize}
    \item Mixed Quantifiers - multiple quantifiers
    \begin{itemize}
        \item $(\forall x)(\exists y) \quad x>y$ is true
        \item $(\exists y)(\forall x) \quad x>y$ is false
    \end{itemize}
    \item Negations with quantifiers
    \begin{itemize}
        \item Negation for $(\forall x \in \mathbb{R})P(x)$ is $(\exists x \in \mathbb{R})\neg P(x)$: if at least one real $x$ makes $P(x)$ false
        \item Negation for $(\exists x \in \mathbb{R})P(x)$ is $(\forall x \in \mathbb{R})\neg P(x)$: if for all real $x$, $P(x)$ is false
    \end{itemize}
    \item Consider $(\forall y \in \mathbb{R})(\exists x \in \mathbb{R}) x+2 = y$
    \begin{itemize}
        \item Negation is $(\forall x \in \mathbb{R})(\exists y \in \mathbb{R}) x+2 \neq y$
    \end{itemize}
    \item Proof by Negation/Proof by Contradiction (in classical logic; not constructive logic) - proving a statement is true by proving its negation is false
    \item Statement - $(\forall n \in \mathbb{Z}) n^2 \text{ is even} \Rightarrow n \text{ is even}$
    \begin{itemize}
        \item Negation: $(\exists n \in \mathbb{Z})n^2 \text{ is even} \nRightarrow n \text{ is even}$
    \end{itemize}
\end{itemize}

\section{Tutorial 1 (09/10/2025)}
\begin{itemize}
    \item Thivan Gunawardana
    \begin{itemize}
        \item Originally from Sri Lanka
        \item Studied maths at Imperial for undergrad, theoretical physics at Oxford for masters, and came back to Imperial for PhD in mathematical physics
    \end{itemize}
    \item Key realisation: $a \preceq a$ is not equivalent to $a=a$
    \begin{itemize}
        \item $a=a$ is always true for any $a \in A$ regardless of whether or not any order relation is defined on $A$
        \item Simply put, $=$ is itself a separate relation which every set has.
        \item Reflexivity says that $a$ is also related to itself by this other relation $\preceq$
        \item For example, $1=1$ is true in $\mathbb{Z}$ but $1<1$ is not true, so the relation is not reflexive.
    \end{itemize}
    \item Consider the well-known game on the set $A = \{$rock, paper, scissors$\}$. Rock beats scissors, scissors beats paper, paper beats rock. Is $a \preceq b$ if $b$ does not lose to $a$ a total ordering, partial ordering, or neither 
    \begin{itemize}
        \item It is not transitive, so neither.
        \item Reflexivity breaks if we take $a \preceq b$ as $b$ beats $a$
    \end{itemize}
    \item Exercise 1. Provide an example of a non–totally ordered poset that has a greatest element?
    \begin{itemize}
        \item Many possibilities. Take for example the subset relation $\subseteq$ on the power set $P(S)$ of the set $S = \{1, 2\}$ (or any other set). Then $S = \{1, 2\}$ is the greatest element of $P(S)$ and it’s not a total order because $\{1\}$ and $\{2\}$ aren’t comparable
    \end{itemize}
    \item Exercise 2. Construct a poset that has exactly one maximal element. (Bonus: make it not a greatest element.)
    \begin{itemize}
        \item Many possibilities again. The key is that for the only maximal element to not be a greatest element, the poset must be infinite! For example, take the usual $\le$ on the positive integers excluding 2 and let $1 \le 2$. Then 2 becomes the only maximal element but it is not a greatest element.
    \end{itemize}
    \item Exercise 3. Can you prove that every nonempty finite poset has at least one maximal element?
    \begin{itemize}
        \item Proof by contradiction. Suppose that $A$ is a nonempty finite poset of size $N$ with no maximal element. Take an element $a_1 \in A$. Then $\exists a_2 \in A$ such that $a_1 \prec a_2$
        \item Continue this way until we get to $a_N$: $a_1 \prec a_2 \prec \cdots \prec a_{N-1} \prec a_N$
        \item But $a_N$ cannot be maximal either, so $\exists a_{N+1}$ such that $a_N \prec a_{N+1}$
        \item All of the $a_i$s must be distinct by transitivity, and so we have now found $N+1$ different elements in $A$. But we said $A$ has only $N$ elements. This is a contradiction, so $A$ must have at least one maximal element.
    \end{itemize}
    \item Exercise 4. Let's consider the following statements: $P = $ "Paul is happy", $Q =$ "Paul paints a picture", $R =$ "Richard is happy"
    \begin{itemize}
        \item $(P \wedge Q) \Rightarrow \neg R:$ "If Paul is happy and paints a picture then Richard isn't happy"
        \item $P \Rightarrow Q$: "If Paul is happy, then he paints a picture"
        \item $P \Rightarrow Q$: "Paul is happy only if he paints a picture"
    \end{itemize}
    \item $Q$ if $P$ is written $P \Rightarrow Q$ while $Q$ only if $P$ is written $Q \Rightarrow P$
    \item Exercise 7. Use a truth table to prove that $A \wedge (A \vee B) = A$
    \[
    \begin{array}{|c|c|c|c|c|}
        \hline
        A & B & A \vee B & A \wedge (A \vee B) & A \\
        \hline
        T & T & T & T & T \\
        T & F & T & T & T \\
        F & T & T & F & F \\
        F & F & F & F & F \\
        \hline
    \end{array}
    \]
    \item Exercise 9. Jia finds two trunks A and B in a cave. She knows that each of them either contains a treasure or a fatal trap. On trunk A is written: “At least one of these two trunks contains a treasure.” On trunk B is written: “In A there’s a fatal trap.” Jia knows that either both the inscriptions are true, or they are both false. Can Jia choose a trunk being sure that she will find a treasure? If this is the case, which trunk should she open?
    \begin{itemize}
        \item For each trunk, let $T$ (true) denote it contains treasure and $F$ (false) denote it contains a fatal trap. Then, we can rewrite the statements on the two trunks as:
        \begin{itemize}
            \item Trunk A: $A \vee B$
            \item $\neg A$
        \end{itemize}
        \item We can set up a truth table to go through all combos:
        \[
        \begin{array}{|c|c|c|c|}
            \hline
            A & B & A \vee B & \neg A \\
            \hline
            T & T & T & F \\
            T & F & T & F \\
            F & T & T & T \\
            F & F & F & T \\
            \hline
        \end{array}
        \]
        \item We can see that there is only one combo for which both statements have the same truth values, so this must be the only possibility! This is the third row which corresponds to only B having a treasure so Jia should open B.
    \end{itemize}
\end{itemize}

\section{Lecture 2 (14/10/2025)}
\subsection{Slide 2 (Analysis: Numbers)}
\begin{itemize}
    \item Natural Numbers - $\mathbb{N} = \{1,2,3,\dots\}$
    \item Integers - $\mathbb{Z} = \{\dots, -3,-2,-1,0,1,2,3,\dots \}$
    \item Rational Number - $\mathbb{Q} = \{\frac{p}{q}: p \in \mathbb{Z} \wedge q \in \mathbb{Z}-\{0\}\}$
    \item Prime Numbers - LCM, HCF, coprime, Prime Factorisation Theorem
    \item $0.999999\dots = 1$
    \begin{itemize}
        \item $0.999999 = x$
        \item $9.999999\dots = 10x$
        \item $9+0.9999\dots = 10x$
        \item $9+x = 10x$
        \item $9x=9$
        \item $x=1$
        \item $\sum_{i=1}^n 0.9 \cdot 10^{-n}$
    \end{itemize}
    \item Proposition: Every rational number has a finite decimal expansion or an infinite recurring decimal expansion
    \begin{itemize}
        \item Definition: every rational number can be expressed in the form $\frac{p}{q}$ where $p$ and $q$ are integers and $q \neq 0$
        \item Assume that the fraction is in its lowest terms (i.e., $p$ and $q$ are coprime)
        \item We can confine ourselves to $p,q>0$ - the decimal expansion of the positive ration $\frac{p}{q}>0$ can be turned into the decimal expansion of the negative rational $\frac{p}{q} >0$ merely via multiplication by $-1$
        \item We can also confine ourselves to $p<q$, because for $p>q$, we can write $\frac{p}{q} = N 
        + \frac{p'}{q}$ where $N \in \mathbb{N}$ and $p' \in \mathbb{N}$ and $p'<q$
        \item Divide $p$ by $q$. We know that because $p < q$, we start with $0.\dots$ and to obtain the first entry after the decimal place, we need to find the smallest $n \in \mathbb{Z}^+$ such that $10^n p > q$. We then place $n-1$ zeros after the decimal place and consider $\frac{10^n p}{q}$
        \item If the remainder is 0, then the decimal expansion terminates and is finite
        \item Otherwise, we obtain the non-zero remainder $r_1$ via $\frac{p \times 10^n}{q} = a_1 + \frac{r_1}{q}$, where $r_1$ must belong to $\{1,2,\dots,q-1\}$ and $a_1 \in \{1, \dots 9\}$.
        \item We therefore have $\frac{p}{q} = -. \underset{{n-1 \text{ zeros}}}{\underbrace{0 \dots 0}} a_1 \dots$
        \item To see what comes after $a_1$, we need to find the first $n \in \mathbb{N}$ such that $10^n r_1 > q$ and compute $\frac{10^n r_1}{q}$:
        \begin{itemize}
            \item $\frac{10^n r_1}{q} = a_2 + \frac{r_2}{q}$
        \end{itemize}
        \item There are three probabilities:
        \begin{itemize}
            \item If $r_2=0$, the decimal expansion terminates and is clearly finite
            \item $r_2=r_1$, then we obtain a recurring decimal $\frac{p}{q} = 0.0 \dots 0\underset{zeros}{\underbrace{\overline{a_1 \cdots} a_2}}$
            \item Otherwise, $r_2 = \{1,2, \dots, q-1\} - \{r_1\}$ and we consider the first $n \in \mathbb{N}$ such that $10^n r_2 > q$
        \end{itemize}
        \item As before, we have three possibilities for the remainder $r_3$
        \begin{itemize}
            \item $r_3 =0$, the decimal expansion terminates and is clearly finite
            \item $r_3=r_1$ or $r_3 = r_2$, and so we obtain a recurring decimal
            \item $r_3 = \{1,2,\dots,q-1\}-\{r_1,r_2\}$ and we consider the first $n \in \mathbb{N}$ such that $10^n r_3 >q$
        \end{itemize}
        \item We keep on executing the above loop
        \item Given that the number of possible remainders is finite, eventually a remainder must be 0 or repeated
        \item When a zero remainder is obtained, then the the loop ends and the decimal expansion is finite. When a remainder is repeated, then we know how the loop will continue, allowing us to write down the recurring decimal expansion, which is infinite
    \end{itemize}
    \item Irrationality of $\sqrt{2}$
    \begin{itemize}
        \item $1 < x^2 < 4 \Rightarrow 1<x<2$
        \item $14^2 < 10^2x^2 < 15^2 \Rightarrow 14 < 10x < 15 \Rightarrow 1.4 < x < 1.5$
        \item $141^2 < 100^2x^2 < 142^2 \Rightarrow 141 < 100x < 142 \Rightarrow 1.41 < x < 1.42$
        \item In general: for a given $n \in \mathbb{N}$, find $m \in \mathbb{N}$ such that $m^2 < 10^{2n}x^2 < (m+1)^2$ and then go through the following steps:
        \begin{itemize}
            \item $m^2 < 10^{2n}x^2 < (m+1)^2 \Rightarrow m < 10^nx < m+1 \Rightarrow m \times 10^{-n} < x < (m+1) \times 10^{-n}$
        \end{itemize}
        \item Note that the distance between $m \times 10^{-n}$ and $(m+1) \times 10^{-n}$ is $10^{-n}$ and so the maximum distance between $m \times 10^{-n}$ and $x$ is strictly less than $10^{-n}$, i.e., $x-m \times 10^{-n} < 10^{-n}$
        \item Similarly, $(m+1) \times 10^{-n} - x < 10^{-n}$
        \item Therefore, $\vert m \times 10^{-n} -x \vert < 10^{-n}$ and $\vert (m-1) \times 10^{-n} -x \vert < 10^{-n}$
    \end{itemize}
    \item Proof by contradiction (reductio ad absurdum)
    \begin{itemize}
        \item Proposition: $\sqrt{2} \notin \mathbb{Q}$
        \item Assume $\sqrt{2}$ is rational
        \begin{itemize}
            \item $(\exists p,q \in \mathbb{N}) \left( HCF(p,q) = 1 \wedge \sqrt{2} = \frac{p}{q} \right)$
        \end{itemize}
        \item We can see that $(\exists p,q \in \mathbb{N}) \left( HCF(p,q) = 1 \wedge 2q^2 = p^2 \right)$
        \item Therefore $p^2$ is even and so $p$ is even.
        \begin{itemize}
            \item $(\exists r \in \mathbb{N})(2r=p)$
        \end{itemize}
        \item But $2q^2=p^2$, and so
        \begin{itemize}
            \item $2q^2=4r^2 \Rightarrow q^2 = 2r^2$
        \end{itemize}
        \item We see that $q^2$ is even, which implies $q$ is even
        \item Therefore, $p$ and $q$ are both even which means they have a common factor of 2, which violates our initial assumption
        \item Therefore, our initial assumption is false and $\sqrt{2}$ is not a rational number
    \end{itemize}
    \item The set $\mathbb{Q}$ of rational numbers is incomplete, with gaps like that at $\sqrt{2}$
    \begin{itemize}
        \item We can complete $\mathbb{Q}$ by filling in its gaps via what is known as a Dedekind cut (not examinable)
    \end{itemize}
    \item $Set = \{ O, I\}$
    \begin{itemize}
        \item $O+O=O$
        \item $O+I=I$
        \item $I+O=I$
        \item $I+I=O$
        \item $O \times O = O$
        \item $O \times I = O$
        \item $I \times O = O$
        \item $I \times I = I$
        \item Cayley Table
        $\begin{array}{c | c | c}
            + & O & I \\
            \hline
            O & O & I \\
            \hline
            I & I & O
        \end{array}$
        \item The set $\{O,I\}$ defined as above forms a field
        \begin{itemize}
            \item It is called the binary field and denoted by $\mathbb{F}_2$
        \end{itemize}
        \item The definition $I+I=O$ implies that the additive inverse of $I$ is $I$ and we can represent this fact as $-I=I$
    \end{itemize}
    \item Field - set $\mathbb{F}$ with two binary operations $+$ and $\times$ that satisfy the following properties or axioms:
    \begin{itemize}
        \item Closure under Addition - For every $a,b \in \mathbb{F}$, $a+b \in \mathbb{F}$
        \item Closure under Multiplication - For every $a,b \in \mathbb{F}, a \times b \in \mathbb{F}$
        \item Associative Law of Addition - For every $a,b,c \in \mathbb{F}, (a+b)+c = a+ (a+b)$
        \item Associative Law of Multiplication - For every $a,b,c \in \mathbb{F}$, $(a \times b) \times c = a \times (b \times c)$
        \item Commutative Law of Addition - For every $a,b \in \mathbb{F}, a+b = b+a$
        \item Commutative Law of Multiplication - For every $a,b \in \mathbb{F}, a\times b = b \times a$
        \item Existence of Additive Identity - There exists an element $0 \in \mathbb{F}$ such that for every $a \in \mathbb{F}, a+0 = a$
        \item Existence of Multiplicative Identity - There exists an element $1 \in \mathbb{F}$ such that $1 \neq $ and for every $a \in \mathbb{F}, a \times 1 = a$
        \item Existence of Additive Inverses - For every $a \in \mathbb{F}$, there exists an element $-a \in \mathbb{F}$ such that $a + (-a)=0$
        \item Existence of Multiplicative Inverses - For every $a \in \mathbb{F}$ with $a \neq 0$, there exists an element $a^{-1} \in \mathbb{F}$ such that $a \times a^{-1} = 1$
        \item Distributive Law - For every $a,b,c \in \mathbb{F}, a \times (b+c) = (a \times b) + (a \times c)$
    \end{itemize}
    \item Lemma 1 - In a field $\mathbb{F}$, the additive identity satisfies the following property: $(\forall x \in \mathbb{F})x \times 0 = 0$
    \begin{itemize}
        \item Start with the Distributive Law - $(\forall a,b,c \in \mathbb{F})a \times (b+c)=(a \times b) + (a \times c)$ with $a=x$, $b=x$, and $c=0$
        \item $x \times (x+0)=(x \times x) + (x \times 0)$
        \item We now use the Existence of Additive Identity - $x+0=x$
        \begin{itemize}
            \item $x \times x = (x \times x) + (x \times 0)$
            \item $x^2 = x^2 + (x \times 0)$
        \end{itemize}
        \item Use the Existence of Additive Inverses to see that $x^2+(-x^2)=0$ and Associativity of Addition together with Commutativity of Addition
        \begin{itemize}
            \item $x^2 + (-x^2) = (x^2 + (x \times 0)) + (-x^2)$
            \item $x^2 + (-x^2) = 
            ((x \times 0) + x^2) + (-x^2)$
            \item $x^2 + (-x^2) = (x \times 0) + (x^2 + (-x^2))$
            \item $0 = (x \times 0) + 0$
            \item $0 = (x \times 0)$
        \end{itemize}
    \end{itemize}
    \item There are important properties of the reals which do not follow from the field axioms alone
    \begin{itemize}
        \item For example, we cannot show that $(\forall x \in \mathbb{F} - \{0\})x \neq -x$. Let us denote this statement by $P$
        \item If $P$ did indeed follow from the field axioms, it would be true for all fields, i.e. $(\forall \text{fields}) P \text{ is true}$
        \item The negation of this statement is $(\exists \text{ a field}) P \text{ is false}$
        \item We know the negation is true because the statement $P$ is false in the field $\mathbb{F}_2$. This is instance of showing a claim is false via a counterexample
    \end{itemize}
\end{itemize}

\section{Lecture 3 (16/10/2025)}
\begin{itemize}
    \item Lemma 2 - In a field $\mathbb{F}, (\forall x \in \mathbb{F})$, the additive inverse of $x$ is $-1 \times x$, i.e. $(\forall x \in \mathbb{F}) -x = -1 \times x$
    \begin{itemize}
        \item Suppose $x \in \mathbb{F}$, where $\mathbb{F}$ is a field
        \begin{flalign*}
            x+(-1 \times x) &= \overset{\text{multiplicative identity}}{\overbrace{x \times 1}} + \overset{\text{commutativity of } \times}{\overbrace{x \times (-1)}} \\
            &= x \times (1 + (-1)) \quad \text{Distributive Law}
            &= x \times 0 \quad \text{Additive inverse}
            &= 0 \quad \text{Lemma 1}
        \end{flalign*}
        \item Using additive inverse of $x$, denoted by $-x$, we have
        \begin{flalign*}
            (x + (-1 \times x)) + (-x) = -x \quad &\text{Additive inverse} \\
            x+(-1 \times x) + (-x) = -x \quad &\text{Associativity of addition} \\
            x+(-x) + (-1 \times x) = -x \quad &\text{Commutativity of addition} \\
            (x+(-x)) + (-1 \times x) = -x \quad &\text{Associativity of addition} \\
             -1 \times x = -x \quad &\text{Additive inverse}
        \end{flalign*}
    \end{itemize}
    \item Lemma 3 - The additive inverse of the multiplicative identity satisfies the following equation $(-1)^2 = 1$
    \begin{itemize}
        \item By definition, the additive inverse of the multiplicative identity is $-1$, where $1+(-1)=0$
        \begin{flalign*}
            (\forall a,b,c \in \mathbb{F}) \ a \times (b \times c) = (a \times b) + (a \times c) \qquad &\text{Distributive Property} \\
            (-1) \times ((-1) + 1) = (-1 \times -1) + (-1 \times 1) \qquad &\text{Substitute in } a=-1,b=-1,c=1 \\
            (-1) \times (1 - 1) = (-1 \times -1) + (-1 \times 1) \qquad &\text{Commutativity of Addition} \\
            (-1) \times 0 = (-1)^2 + -1 \qquad &\text{Multiplicative identity} \\
            0 = (-1)^2 + -1 \qquad &\text{Lemma 1} \\
            0 + 1 = (-1)^2 + (-1) + 1 \qquad &\text{Add 1 to both sides} \\
            1 = (-1)^2 + (-1) + 1 \qquad &\text{Additive Identity} \\
            1 = (-1)^2 \qquad &\text{Additive Inverse}
        \end{flalign*}
    \end{itemize}
    \item Trichotomy - the distinguishing feature of a toset that every pair of elements can be ordered
    \item Ordered Field - field $\mathbb{F}$ with a relation $\preceq$ that makes $\mathbb{F}$ into a totally ordered set such that
    \begin{itemize}
        \item $x \preceq y$ implies $x+z \preceq y+z$ (addition preserves order)
        \item $x \preceq y$ and $0 \preceq z$ imply $x \times z \preceq y \times z$ (multiplication preserves order when multiplying by a positive number)
        \item In $\mathbb{F}_2$ addition does not preserve order, so $\mathbb{F}_2$ cannot be an ordered field
    \end{itemize}
    \item Lemma 4 - Suppose $\mathbb{F}$ is an orderer field. Then $(\forall x \in \mathbb{F} - \{0\}) x \neq -x $
    \begin{itemize}
        \item Let $\mathbb{F}$ be an ordered field and let $x \in \mathbb{F} - \{0\}$
        \item We consider the two elements $x$ and $0$, which can be compared because they are both elements of the field, which is a toset
        \item Either $x \prec 0 \vee x = 0 \vee 0 \prec x$. Case $x=0$ ruled out by assumption
        \item Case 1. $x \prec 0$
        \begin{flalign*}
            x \prec 0 \\
            x + x \prec 0 + x \qquad &\text{Addition preserves order} \\
            x+x \prec x \qquad &\text{Existence of Additive Identity} \\
            x+x \prec 0 \qquad &x+x \prec 0 + x \wedge x \prec 0 \text{ Transitivity} \\
            (x+x) + (-x) \prec 0 + (-x) \qquad &\text{Addition preserves order} \\
            x +(x + (-x)) \prec 0 + (-x) \qquad &\text{Associativity of addition} \\
            x + 0 \prec 0 + (-x) \qquad &\text{Additive inverse} \\
            x \prec -x \qquad &\text{Additive identity}
        \end{flalign*}
        \item Case 2. $0 \prec x$
        \begin{flalign*}
            0 \prec x \\
            0 + x \prec x + x \qquad &\text{Addition preserves order} \\
            x \prec x+x \qquad &\text{Existence of Additive Identity} \\
            0 \prec x+x \qquad &0+x \prec x+x \wedge 0 \prec x \text{ Transitivity} \\
            0 + (-x) \prec x+x + (-x) \qquad &\text{Addition preserves order} \\
            0 + (-x) \prec x+ (x + (-x)) \qquad &\text{Associativity of addition} \\
            0 + (-x) \prec x +0 \qquad &\text{Additive inverse} \\
            -x \prec x \qquad &\text{Additive identity}
        \end{flalign*}
        \item Proved that $x \prec -x \vee -x \prec x \Rightarrow x \neq -x$
    \end{itemize}
    \item Lemma 5 - Let $\mathbb{F}$ be an ordered field and $x \in \mathbb{F}$. Then $x^2 \succeq 0$
    \begin{itemize}
        \item By trichotomy, either $x \prec 0, x = 0, x \succ 0$
        \item Case 1. $x=0$
        \begin{flalign*}
           x \times x = 0 \times 0 \qquad &\text{Positive multiplication preserves order} \\
           x^2 = 0 \qquad &\text{Lemma 1}
        \end{flalign*}
        \item Case 2. $x \succ 0$
        \begin{flalign*}
            x^2 \succ 0 \times x \qquad &\text{Positive multiplication preserves order} \\
            x^2 \succ 0 \qquad &\text{Lemma 1}
        \end{flalign*}
        \item Case 3. $x \prec 0$
        \begin{flalign*}
            x + (-x) \prec 0 + (-x) \qquad &\text{Addition preserves order} \\
            0 \prec 0 + (-x) \qquad &\text{Additive inverse} \\
            0 \prec -x \qquad &\text{Additive identity} \\
            0 \times (-x) \prec (-x) \times (-x) \qquad &\text{Positive multiplication preserves order} \\
            0 \prec (-x)^2 \qquad &\text{Lemma 1} \\
            \Rightarrow (-x)^2 = (-x) \times (-x) \\
            = (-1 \times x) \times (-1 \times x) \qquad &\text{Lemma 2} \\
            = (x \times -1) \times (-1 \times x) \qquad &\text{Commutativity of multiplication} \\
            x \times (-1 \times (-1 \times x)) \qquad &\text{Associativity of multiplication} \\
            =x \times ((-1 \times -1) \times x) \qquad &\text{Associativity of multiplication} \\
            =x \times (1 \times x) \qquad &\text{Lemma 3} \\
            =x \times x \qquad &\text{Multiplicative identity} \\
            =x^2 \\
            \Rightarrow 0 \prec x^2 \qquad \text{Substitute back in}
        \end{flalign*}
        \item We have proven that $x^2 = 0 \vee x^2 \succ 0 \Rightarrow x^2 \succeq 0$
    \end{itemize}
    \item In $S=\{1,2\} \subset \mathbb{Z}^+,$ 2 is the least upper bound and the maximum of $S$
    \begin{itemize}
        \item $2,3,4,\dots \subset \mathbb{Z}^+$ are all upper bounds
    \end{itemize}
    \item Upper Bound - Let $X$ be a totally ordered set and let $S \subseteq X$. An upper bound for $S$ is an element $B \in X$ such that $(\forall s \in S)s \preceq B$. If the subset $S$ has an upper bound, then we say that $S$ is bounded above
    \begin{itemize}
        \item We need X to be totally orderd so comparison $\preceq$ is possible
        \item Upper bound doesn't have to be unique
        \item Upper bound for $S$ does not have to be in $S$
    \end{itemize}
    \item Least Upper Bound - Upper bound $\underline{B}$ for $S$ is a least upper bound/supremum if nothing smaller than $\underline{B}=\sup (S)$ is an upper bound
    \begin{itemize}
        \item $(\forall s \in S)s \preceq \underline{B}$ ($\underline{B}$ is an upper bound)
        \item $(\forall y \prec \underline{B})(\exists a \in S) a \succ y$ (every element smaller than $\underline{B}$ is not an upper bound
    \end{itemize}
    \item Maximum - if $\sup S \in S$, $\sup S$ is a maximum of $S$
    \item Lower Bound - Let $X$ be a totally ordered set and let $S \subseteq X$. A lower bound for $S$ is an element $B \in X$ such that $(\forall s \in S)s \succeq B$. If the subset $S$ has an lower bound, then we say that $S$ is bounded below
    \item Greatest Lower Bound - Lower bound $\overline{b}$ for $S$ is a greatest lower bound/infimum if nothing larger than $\underline{B}=\inf (S)$ is an lower bound
    \begin{itemize}
        \item $(\forall s \in S)s \succeq \overline{b}$ ($\overline{b}$ is a lower bound)
        \item $(\forall y \succ \overline{b})(\exists a \in S) a \prec y$ (every element larger than $\overline{b}$ is not a lower bound)
    \end{itemize}
    \item Minimum - if $\inf S \in S$, $\inf S$ is a minimum of $S$
    \item $X = \mathbb{Q}$ and $A={x \in X: 0 \le x \le 1}$ has a supremum (and maximum) of 1 and infimum (and minimum) of 0
    \item $X = \mathbb{Q}$ and $A={x \in X: 0 < x < 1}$ has $\sup(S) = 1$ but does not have a maximum or minimum
    \item $X = \mathbb{Q}$ and $A={x \in X: x^2 < 2}$ does not have a supremum
    \begin{itemize}
        \item 1.5 is an upper bound, and so is 1.415 and 1,4143. For any rational number greater than $\sqrt{2}$, I can always find a smaller rational number which is strictly greater than $\sqrt{2}$. This is why there is not least upper bound/supremum. Note $\sqrt{2}$ is not an upper bound since $\sqrt{2} \notin X$.
    \end{itemize}
    \item Least Upper Bound Property - Toset $X$ has least upper bound property if every non-empty subset of $X$ that is bounded above has a supremum
    \item Greatest Lower Bound Property - Toset $X$ has greatest lower bound property if every non-empty subset of $X$ that is bounded below has an infimum
    \item The rationals are an ordered field, but do not satisfy the least upper bound property because it has "gaps"
    \begin{itemize}
        \item It is simple to check that an ordered field with the least upper bound property has the greatest lower bound property.
    \end{itemize}
    \item Real Numbers - $\mathbb{R}$ is an ordered field with the least upper bound property
    \begin{itemize}
        \item The least upper bound property ensures there are no 'gaps'
        \item Such a field is unique (up to isomorphism)
    \end{itemize}
    \item Lemma 6 (Archimedean Property v1) - Let $\mathbb{F}$ be an ordered field with the least upper bound property (real numbers). Then the set of natural numbers $\mathbb{N} = \{1,2,3,\dots\} \subset \mathbb{F}$ is not bounded above
    \begin{itemize}
        \item Assume $\{1,2,3,\dots\}$ is bounded above
        \item We also know $\mathbb{F}$ has the least upper bound property and $\{1,2,3, \dots \}$ is a subset of $\mathbb{F}$
        \item $\{1,2,3, \dots\}$ has a least upper bound, which we denote by $x$
        \item Therefore, $x-1$ is not an upper bound for $\{1,2,3,\dots\}$ and so we can find $n \in \{1,2,3,\dots\}$ such that $n > x-1$
        \item But $n+1 \in \{1,2,3, \dots\}$ and $n+1 > x$, so we have a contradiction
        \item Therefore, the set $\{1,2,3,\dots\} \subset \mathbb{F}$ is not bounded above
    \end{itemize}
    \item Absolute Value Function - $\vert \cdot \vert : \mathbb{R} \rightarrow \mathbb{R}$
    \[ \vert x \vert = \left\{ \begin{array}{cc}
         x, &x \ge 0 \\
         -x, &x < 0
    \end{array} \right. \]
    \item If $x \le y + \epsilon$ for every$\epsilon > 0:x \le y$
    \begin{itemize}
        \item Start from $(\forall \epsilon > 0) x \le y + \epsilon$
        \item We know $\mathbb{R}$ is a totally ordered set, and so we have either $x \le y$ or $x > y$. We now rule out the case $x>y \Longleftrightarrow x-y > 0$, and so we can pick $\epsilon = \frac{1}{2}(x-y)$
        \[ (\forall \epsilon > 0) x \le y + \epsilon \Rightarrow x \le y + \frac{1}{2} (x-y) \Longleftrightarrow x-y \le 0 \Longleftrightarrow x \le y \]
    \end{itemize}
    \item $\vert x-y \vert < \epsilon$ for every $\epsilon >0 : x = y$
    \begin{itemize}
        \item Assume $x \neq y$, which implies $\frac{1}{2} \vert x-y\vert > 0$
        \[ (\forall \epsilon > 0) \vert x-y \vert < \epsilon \Rightarrow \vert x-y \vert  < \frac{1}{2} \vert x-y \vert \Longleftrightarrow \vert x-y \vert < 0\]
        \item This contradicts $x \neq y \Rightarrow x=y$
    \end{itemize}
    \item Theorem 1 (Triangle Inequality) - If $x,y \in \mathbb{R}$, then $\vert x+y \vert \le \vert x \vert + \vert y \vert$
    \begin{flalign*}
        &- \vert x \vert \le x \le \vert x \vert \wedge - \vert y \vert \le y \le \vert y \vert \\
        \Rightarrow &- \vert x \vert - \vert y \vert \le x+y \le \vert x \vert + \vert y \vert \\
        \Longleftrightarrow &-(\vert x \vert + \vert y \vert) \le x+y \le \vert x \vert + \vert y \vert \\
        \Longleftrightarrow &\vert x + y \vert \le \vert x \vert + \vert y \vert
    \end{flalign*}
    \item Induction
    \begin{itemize}
        \item The construction of natural numbers underlies the Principle of Mathematical Induction, which we can state as:
        \[ \forall P(P(1) \wedge (\forall k \in \mathbb{N})P(k) \Rightarrow P(k+1) \Rightarrow (\forall n \in \mathbb{N})P(n)) \]
        \item $P(n)$ is some
        statement, such as $\sum_{i=1}^n i = \frac{1}{2}n(n+1)$
        \begin{itemize}
            \item P(1)
            \[ \sum_{i=1}^1 i = \frac{1}{2}(1)(1+1) = 1 \]
            \item Assume $P(k)$ is true
            \[ \sum_{i=1}^k i = \frac{1}{2}(k)(k+1) \]
            \item Prove $P(k+1)$ is true
            \[ \sum_{i=1}^{k+1} i = \sum_{i=1}^k i + k+1 = \frac{1}{2}(k)(k+1)+k+1 = \frac{k(k+1)}{2}+\frac{2(k+1)}{2} = \frac{k^2+3k+2}{2}= \frac{(k+1)(k+2)}{2} \]
            \item Therefore, by the principle of mathematical induction:
            \[ (\forall n \in \mathbb{Z}^+) \sum_{i=1}^n i =\frac{n(n+1)}{2}\]
        \end{itemize}
    \end{itemize}
\end{itemize}

\section{Tutorial 2 (16/10/2025)}
\begin{itemize}
    \item Exercise 1. Find $\sqrt{7}$ to 3 decimal places by hand without aid from a machine.
    \begin{itemize}
        \item Start with
        \[2^2 < 7 < 3^2 \Rightarrow 2 < \sqrt{7} < 3
        \]
        \item Multiplying by $100=10^2$, we can refine our estimate further:
        \[
        26^2 < 7 \times 10^2 < 27^2 \Rightarrow 2.6 < \sqrt{7} < 2.7
        \]
        \item Continuing in this way, we find
        \[
        264^2 < 7 \times 10^4 < 265^2 \Rightarrow 2.64 < \sqrt{7} < 2.65
        \]
        \[
        2645^2 < 7 \times 10^6 < 2646^2 \Rightarrow 2.645 < \sqrt{7} < 2.646
        \]
    \end{itemize}
    \item Exercise 2. Write computer code in R or Python to find the square root of a natural number to any required degree of accuracy. You cannot use the square root function!
    \begin{itemize}
        \item For this, you are allowed to use ChatGPT to generate the Python code. You can then also run the code via ChatGPT (you just have to tell it to run the code for whichever inputs you specify).
    \end{itemize}
    \item Exercise 3. Prove that $\sqrt{3}$ is irrational.
    \begin{itemize}
        \item We do this in the same way as proving the irrationality of $\sqrt{2}$
        \item Suppose for a contradiction that $\sqrt{3} \in \mathbb{Q}$ so that $(\exists a,b \in \mathbb{N}) \ HCF(a,b) = 1 \wedge \sqrt{3} = \frac{a}{b}$
        \item $3 = \frac{a^2}{b^2} \Rightarrow a^2 = 3b^2$
        \item Thus $3$ divides $a^2$. Since $3$ is prime, this implies $3$ divides $a$.
        \item Writing $a = 3k$ where $k \in \mathbb{N}$, we find that $b^2 = 3k^2$, which in turn implies that $b = 3m$ for some $m \in \mathbb{N}$
        \item But this contradicts the assumption that $HCF(a,b)=1$. Therefore, $\sqrt{3} \notin \mathbb{Q}$ and is in fact irrational.
        \item This proof also generalises for $\sqrt{p}$ where $p$ is a prime since only then can we claim that if $p \mid a^2 \Rightarrow p \mid a$
    \end{itemize}
    \item Exercise 5. Two quantities are in the golden ratio if their ratio is the same as the ratio of their sum to the larger of the two quantities. Expressed algebraically, for quantities $a$ and $b$ with $a > b > 0$,
    \[
    \frac{a+b}{a} = \frac{a}{b} = \varphi
    \]
    where $\varphi$ denotes the golden ratio.
    \begin{itemize}
        \item (a) Prove that $\varphi^2 = \varphi + 1$ and hence derive an expression for $\varphi$.
        \begin{itemize}
            \item Note that from the given definition,
            \[
            \varphi = \frac{a}{b} = 1 + \frac{b}{a} = 1 + \frac{1}{\varphi}
            \]
            \item From this, it follows immediately that
            \[
            \varphi^2 = \varphi + 1
            \]
            \item We solve the quadratic equation and get
            \[
            \varphi = \frac{1 \pm \sqrt{5}}{2}
            \]
            \item But $\varphi > 0$ since $\varphi = \frac{a}{b}$ for $a,b > 0$. Since $\sqrt{5} > 1$, the only positive value for $\varphi$ is
            \[
            \varphi = \frac{1+\sqrt{5}}{2}
            \]
        \end{itemize}
        \item (b) Prove that $\varphi = \frac{1+\sqrt{5}}{2}$ is irrational.
        \begin{itemize}
            \item To do this properly, we must prove three things:
            \item (1) $\sqrt{5}$ is irrational, which we proved in exercises 3/4
            \item (2) If $a$ is rational and $b$ is irrational, then $a+b$ is irrational.
            \begin{itemize}
                \item Let $a \in \mathbb{Q}$ and $b \notin \mathbb{Q}$. Suppose for contradiction that $a+b \in \mathbb{Q}$.
                \item Then, $a+b = c \in \mathbb{Q}$. This implies that $b = c + (-a) \in \mathbb{Q}$
                \item But this contradicts the assumption that $b \notin \mathbb{Q}$. Therefore, $a+b \notin \mathbb{Q}$.
            \end{itemize}
            \item (3) If $a$ is rational and $b$ is irrational, then $ab$ is irrational.
            \begin{itemize}
                \item Let $a \in \mathbb{Q}$ and $b \notin \mathbb{Q}$. Then $a = \frac{p}{q}$ so $\frac{1}{a} = \frac{q}{p} \in \mathbb{Q}$ as well. Suppose for contradiction that $ab \in \mathbb{Q}$.
                \item Then, $ab = c \in \mathbb{Q}$. This implies that $b = c\frac{1}{a} \in \mathbb{Q}$
                \item But this contradicts the assumption that $b \notin \mathbb{Q}$. Therefore, $ab \notin \mathbb{Q}$.
            \end{itemize}
        \end{itemize}
        \item (c) Explain why
        \[
        \varphi = 1+ \frac{1}{1+\frac{1}{1+\frac{1}{1+\ddots}}}
        \]
        \begin{itemize}
            \item We recall from (a) that $\varphi = 1 + \frac{1}{\varphi}$
            \item We can iterate by substituting the above as $\varphi$ on the RHS. This gives
            \[
            \varphi = 1 + \frac{1}{1+\frac{1}{\varphi}}
            \]
            \item We can keep repeating this, which yields the continued fraction stated in the question. Note that this gives the positive root because the construction remains positive at each step.
        \end{itemize}
    \end{itemize}
    \item Exercise 7. Find out what the reverse triangle inequality is, state it, and prove it.
    \begin{itemize}
        \item The reverse triangle inequality states that for $x,y \in \mathbb{R}$,
        \[
        \vert x - y \vert \ge \vert \vert x \vert - \vert y \vert \vert
        \]
        \item To prove this, it is sufficient to show that for $x,y \in \mathbb{R}$,
        \[
        \vert x - y \vert \ge \vert x \vert - \vert y \vert \qquad \wedge \qquad \vert x-y \vert \ge \vert y \vert - \vert x \vert
        \]
        \item The above is equivalent to
        \[
        \vert x \vert \le \vert x-y \vert + \vert y \vert \qquad \wedge \qquad \vert y \vert \le \vert x-y \vert + \vert x \vert
        \]
        \item Take the first of these, which is
        \[
        \vert x \vert \le \vert x-y \vert + \vert y \vert
        \]
        \item Setting $u = x-y$ and $v = y$, we can use the usual triangle inequality, which states that for $u,v \in \mathbb{R}$,
        \[
        \vert u + v \vert \le \vert u \vert + \vert v \vert \Rightarrow \vert (x-y) + y \vert = \vert x \vert \le \vert x-y \vert + \vert y \vert
        \]
        \item Now setting $u = y-x$ and $v = x$, we get
        \[
        \vert (y-x) + x \vert = \vert y \vert \le \vert y - x \vert + \vert x \vert = \vert x-y \vert + \vert x \vert
        \]
    \end{itemize}
    \item Exercise 8. Let $a,b \in \mathbb{R}^+$, and let us donte
    \[
    QM = \sqrt{\frac{a^2+b^2}{2}}, \quad AM = \frac{a+b}{2}, \quad GM=\sqrt{ab}, \quad HM = \left( \frac{a^{-1}+b^{-1}}{2} \right)^{-1}
    \]
    where $QM$ is the quadratic mean, $AM$ is the arithmetic mean, $GM$ is the geometric mean, and $HM$ is the harmonic mean. Prove that
    \[
    QM \ge AM \ge GM \ge GM
    \]
    with equalities if and only if $a=b$.
    \begin{itemize}
        \item $AM^2 = \frac{(a+b)^2}{4} = \frac{a^2+b^2}{4} + \frac{ab}{2}$
        \item Now recall the useful inequality
        \[
        (a-b)^2 \ge 0 \Rightarrow a^2+b^2 \ge 2ab \Rightarrow ab \le \frac{a^2+b^2}{2}
        \]
        with equality if and only if $a=b$. Then it follows that
        \[
        AM^2 \le \frac{a^2+b^2}{4} + \frac{a^2+b^2}{4} = \frac{a^2+b^2}{2} = QM^2
        \]
        with equality if and only if $a=b$. Since $AM$ and $QM$ are both positive, we get that $AM \le QM$ with equality if and only if $a=b$.
        \item $AM^2 = \frac{(a+b)^2}{4} = \frac{a^2+b^2}{4} + \frac{ab}{2}$
        \item Now recall the useful inequality\[
        (a-b)^2 \ge 0 \Rightarrow a^2+b^2 \ge 2ab \Rightarrow ab \le \frac{a^2+b^2}{2}
        \]
        with equality if and only if $a=b$. 
        \[
        AM^2 \ge \frac{ab}{2} + \frac{ab}{2} = ab = GM^2
        \]
    \end{itemize}
\end{itemize}

\section{Lecture 4 (21/10/2025)}
\begin{itemize}
\subsection{Slide 3 (Analysis: Sequences, Series, and Convergence)}
    \item Economic time series are sequences
    \begin{itemize}
        \item Wealth $(x_t),$ consumption $(c_t)$, inflation $(\pi_t)$, stock prices $(p_t)$
        \item Each data, agents decide how much to save, consume, or invest
        \item These choices determine next date's variables: $x_t+1 = g(x_t, \text{decisions})$
    \end{itemize}
    \item Economy is a system of recursive equations
    \item Real analysis provides precise language to describe limits of such sequences
    \item Baseline Savings Model
    \begin{itemize}
        \item Dates $\mathbb{N}_0,$ income per period (1), saving rate $s \in (0,1)$, net interest $r>0$
        \begin{itemize}
            \item $x_{n+1} = g(x_n) := s[(1+r)x_n+1] = kx_n+s, \qquad k := s(1+r), \quad x_0 \ge 0$
            \item Fixed point: $x^*=g(x^*) \Longleftrightarrow x^* = \frac{s}{1-k} (k \neq 1)$
            \item Classification: $k < 1$ contraction $\Rightarrow x_n \rightarrow x^*; k=1$ linear growth; $k>1$ exponential growth
            \item Closed form: $x_n = x^* + (x_0 - x^*)k^n$
        \end{itemize}
    \end{itemize}
    \item Sequence - a function $a: \mathbb{N} \rightarrow \mathbb{R}$
    \begin{itemize}
        \item Write $a_n,(a_n),(a_n)_1^\infty, (a_n)_{n=1}^\infty$ instead of $a(n)$
    \end{itemize}
    \item Definition (Convergence of sequence)
    \begin{itemize}
        \item Let $(a_n)$ be a sequence and $l \in \mathbb{R}$
        \item $a_n$ converges to $l$, tends to $l$, or $a_n \rightarrow l$, if for all $\varepsilon > 0$, there is some $N \in \mathbb{N}$ such that whenever $n > N$, we have $\vert a_n - l \vert < \varepsilon$
        \item $(\forall \varepsilon > 0)(\exists N \in \mathbb{N})(\forall n \ge N) \vert a_n - l \vert < \varepsilon$
        \item We say $l$ is the limit of $(a_n)$
    \end{itemize}
    \item $\epsilon$-neighborhood - Given a real number $a \in \mathbb{R}$ and a positive number $\epsilon > 0$, the set $V_\epsilon(a) = \{ x \in \mathbb{R} : \vert x-a \vert < \epsilon \}$
    \begin{itemize}
        \item $V_\epsilon(a)$ consists of all of those points whose distance from $a$ is less than $\epsilon$
    \end{itemize}
    \item Definition (Convergence of sequence: Topological Version)
    \begin{itemize}
        \item A sequence $(a_n)$ converges to $a$ if, given any $\epsilon$-neighborhood $V_\epsilon(a)$ of $a$, there exists a point int he sequence after which all of the terms are in $V_\epsilon(a)$
        \item Every $\epsilon$-neighborhood contains all but a finite number of the terms of $(a_n)$
        \item $(\forall \epsilon > 0)(\exists N \in \mathbb{N})(\forall n \ge N)x_n \in V_\epsilon(a)$
        \item This definition is equivalent to the first convergence of sequence definition
        \item Natural number $N$ in the original version of the definition is the point where the sequence $(a_n)$ enters $V_\epsilon(a)$, never to leave
        \item Value of $N$ depends on the choice of $\epsilon$
        \begin{itemize}
            \item The smaller $V_\epsilon(a)$, the larger $N$ may have to be
        \end{itemize}
        \item $a_n \rightarrow a: (a_n)$ will eventually get trapped inside every $\epsilon$-neighbourhood of $a$ 
    \end{itemize}
    \item Proposition (Archimedian Property v2) - $\lim_{n \rightarrow \infty} \frac{1}{n} \rightarrow 0$
    \begin{itemize}
        \item Natural numbers are unbounded
        \begin{flalign*}
            &\Rightarrow (\forall \epsilon > 0)(\exists N \in \mathbb{N})N > \frac{1}{\epsilon} \\
            &\Longleftrightarrow (\forall \epsilon > 0)(\exists N \in \mathbb{N}) \frac{1}{N} < \epsilon \\
            &\Rightarrow (\forall \epsilon > 0)(\exists N \in \mathbb{N}) \left\vert \frac{1}{N} - 0 \right\vert < \epsilon \\
            &\Rightarrow (\forall \epsilon > 0)(\exists N \in \mathbb{N})(\forall n \ge N) \left\vert \frac{1}{n} - 0 \right\vert < \epsilon \\
        \end{flalign*}
    \end{itemize}
    \item Template for a Proof that $(x_n) \rightarrow x$
    \begin{itemize}
        \item Let $\epsilon > 0$ be arbitrary
        \item Demonstrate a choice for $N \in \mathbb{N}$
        \begin{itemize}
            \item This step usually requires the most work, almost all of which is done prior to actually writing the formal proof
        \end{itemize}
        \item Now, show that $N$ actually works
        \item Assume $n \ge N$
        \item With $N$ well chosen, it should be possible to derive the inequality $\vert x_n - x \vert < \epsilon$
    \end{itemize}
    \item Example
    \begin{itemize}
        \item We are now going to use the definition of convergence to prove that $\frac{1}{\sqrt{n}} \rightarrow 0$
        \item Start playing around with the definition by choosing $\epsilon = 0.1$
        \item How large does $N$ have to be so that $\left\vert \frac{1}{\sqrt{N}} - 0 \right\vert < 0.1$
        \item $N=100$ is not quite large enough, but $N=101$ is large enough
        \item Proof
        \begin{itemize}
            \item In general, for every $\epsilon > 0$, we can pick $N > \frac{1}{\epsilon^2}$ (because the natural numbers are unbounded), and so $\frac{1}{\sqrt{N}} < \epsilon$ and $(\forall n \ge N) \frac{1}{\sqrt{n}} < \epsilon$
            \item $(\forall \epsilon > 0)(\exists N \in \mathbb{N})(\forall n \ge N) \left\vert \frac{1}{\sqrt{n}} - 0 \right\vert < \epsilon$
        \end{itemize}
    \end{itemize}
    \item Show $\left( \frac{n+1}{n} \right) \rightarrow 1$
    \begin{itemize}
        \item Consider the inequality $\frac{n+1}{n} -1 < \epsilon$ and then play around with it
        \item Proof
        \begin{itemize}
            \item Since $(\forall n \in \mathbb{N}) \frac{n+1}{n} > 1$, we have
            \[
            \left\vert \frac{n+1}{n} - 1 \right\vert < \epsilon \Longleftrightarrow \frac{n+1}{n} - 1 < \epsilon \Longleftrightarrow 1 + \frac{1}{n} - 1 < \epsilon \Longleftrightarrow n > \frac{1}{\epsilon}
            \]
            \item Define $a_n = \frac{n+1}{n}$ and we know we can pick $N \in \mathbb{N}$ such $N > \frac{1}{\epsilon}$
            \[ (\forall \epsilon > 0)(\exists N \in \mathbb{N})(\forall n \ge N) \vert a_n -1 \vert < \epsilon \]
        \end{itemize}
    \end{itemize}
    \item Proving Non-Convergence (Proof by Contradiction)
    \begin{itemize}
        \item Sequence $(a_n)=(-1)^n$
        \item Assume sequence converges to some limit $L \in \mathbb{R}$
        \[ (\forall \epsilon > 0)(\exists N \in \mathbb{N})(\forall n \ge N) \vert (-1)^n - L \vert < \epsilon \]
        \item Hence, we have $\left\vert (-1)^N - L \right\vert < \epsilon$
        \item Now, $n \ge N \Rightarrow \vert (-1)^n - L \vert < \epsilon$
        \item When $n \pmod 2 = 0$, then $\vert 1 - L \vert < \epsilon$ and $1 - \epsilon < L < 1 + \epsilon$
        \item When $n \pmod 2 = 1$, then $\vert -1 - L \vert < \epsilon$ and $-1 - \epsilon < L < -1 + \epsilon$
        \item When $\epsilon \le 1$, the two intervals are disjoint, so $L$ cannot exist. Therefore, the sequence $a_n = (-1)^n$ does not have a limit and hence does not converge.
    \end{itemize}
\end{itemize}

\section{Tutorial 3 (23/10/2025)}
\begin{itemize}
    \item Exercise 3.9. Compute, without proofs, the suprema and infima of the following sets (let's take them as subsets of the real numbers):
    \begin{itemize}
        \item (a) $S = \{ n \in \mathbb{N} : n^2 < 10 \}$
        \begin{itemize}
            \item $S = \{1,2,3\}$ so $\sup(S)=3$ and $\inf(S)=1$
        \end{itemize}
        \item (b) $S = \{ \frac{n}{(m+n)} : m,n \in \mathbb{N} \}$
        \begin{itemize}
            \item Note that every element is positive. Setting $n=1$, $(\forall m \in \mathbb{N}) \ \frac{1}{(m+1)} \in S$. By the Archimedian property, we find that $\inf(S) = 0$.
            \item Now note that
            \[
            \frac{n}{m+n} = 1 - \frac{m}{m+n}
            \]
            \item So, $\sup(S) = 1 - \inf(S) = 1$
        \end{itemize}
        \item (c) $S = \{ \frac{n}{(2n+1)} : n \in \mathbb{N} \}$
        \begin{itemize}
            \item Once again, we note that
            \[
            \frac{n}{(2n+1} = \frac{1}{2} - \frac{1}{2(2n+1)} < \frac{1}{2} \forall n \in \mathbb{N}
            \]
            \item By the Archimedian property, we deduce that $\sup(S) = \frac{1}{2}$. We can also see that $\inf(S) = \frac{1}{2} - \frac{1}{2(2+1)} = \frac{1}{3}$.
        \end{itemize}
        \item (d) $S = \{ \frac{n}{m} : m,n \in \mathbb{N} \text{ with } m+n \le 10 \}$
        \begin{itemize}
            \item $\inf(S) = \frac{1}{9}$ because out of the finitely many valid pairs $(n,m)$, every other one has a bigger or equal $n$ and a strictly smaller $m$. For similar reasons, $\sup(S) = \frac{9}{1} = 9$
        \end{itemize}
    \end{itemize}
    \item Consider the sequence
    \[
    a_n = \left\{ \begin{array}{lr}
        1, & n \text{ is prime} \\
        \frac{1}{n}, & \text{otherwise}
    \end{array} \right.
    \]
    Does this converge? Can you intuitively explain why?
    \begin{itemize}
        \item No it does not converge. This is because there are infinitely many primes and so no matter how far into the sequence you go, you will always find $n \in \mathbb{N}$ such that $a_n = 1$. Even though the rest of the terms are getting arbitrarily close to zero, no complete tail of the sequence is ever arbitrarily close to zero because of the ones.
    \end{itemize}
    \item Exercise 4.2. Show that the constant sequence $(a,a,a,\dots)$ converges to $a$.
    \begin{itemize}
        \item We want to show that $(\forall \epsilon >0)(\exists N \in \mathbb{N})(\forall n \ge N) \ \vert a_n - a \vert < \epsilon$
        \item Start by fixing some arbitrary $\epsilon > 0$
        \item Note that $a_n = a$ for all $n$ and so $\vert a_n - a \vert = 0 < \epsilon$ for all $n \in \mathbb{N}$
        \item So we can simply pick $N=1$ (or any natural number) in the limit definition and see that it is automatically satisfied.
        \item Since $\epsilon > 0$ was arbitrary, the above works for all $\epsilon > 0$ and so $a_n \rightarrow a$
    \end{itemize}
    \item Exercise 4.3. Let $x_n \ge 0 \forall n \in \mathbb{N}$
    \begin{itemize}
        \item (a) If $x_n \rightarrow 0$, show that $\sqrt{x_n} \rightarrow 0$.
        \begin{itemize}
            \item Start by fixing $\epsilon > 0$. Then $\epsilon^2 > 0$ as well.
            \item Since $x_n \rightarrow 0$, $(\forall \epsilon > 0)(\exists N \in \mathbb{N})(\forall n \ge N) \ \vert x_n - 0 \vert = x_n < \epsilon^2$
            \item Then $(\forall n \ge N) \ \vert \sqrt{x_n} - 0 \vert = \sqrt{x_n} < \sqrt{\epsilon^2} = \epsilon.$
            \item We conclude $(\forall \epsilon > 0)(\exists N \in \mathbb{N})(\forall n \ge N) \ \vert \sqrt{x_n} - 0 \vert < \epsilon \Rightarrow \sqrt{a_n} \rightarrow 0$
        \end{itemize}
        \item (b) If $x_n \rightarrow x$, show that $\sqrt{x_n} \rightarrow \sqrt{x}$
        \begin{itemize}
            \item We start with some rough work:
            \begin{itemize}
                \item We want to get to $\vert \sqrt{x_n} - \sqrt{x} \vert < \epsilon$ for $n \ge N$
                \item We want to involve $\vert x_n - x \vert$ somehow. We remember that
                \[
                \vert \sqrt{x_n} - \sqrt{x} \vert = \frac{\vert x_n - x \vert}{\vert \sqrt{x_n} + \sqrt{x} \vert} \le \frac{\vert x_n - x \vert}{\sqrt{x}}
                \]
                since $x_n \ge 0$
                \item So we define $\epsilon$ as $\epsilon \sqrt{x}$ for everything to work out.
            \end{itemize}
            \item Proof:
            \begin{itemize}
                \item $(\forall \epsilon > 0)(\exists N \in \mathbb{N})(\forall n \ge N) \ \vert x_n - x \vert < \epsilon$
                \item We notice that we can replace $\epsilon$ with $\epsilon \sqrt{x}$ as $x > 0$ since $(\forall n \in \mathbb{N}) \ x_n \ge 0$
                \item $(\forall \epsilon > 0)(\exists N \in \mathbb{N})(\forall n \ge N) \ \vert x_n - x \vert < \epsilon \sqrt{x}$
                \item Thus, since $x_n \ge 0$, we see that $\forall n \ge N$,
                \[
                \vert \sqrt{x_n} - \sqrt{x} \vert = \frac{\vert x_n - x \vert}{\vert \sqrt{x_n} + \sqrt{x} \vert} \le \frac{\vert x_n - x \vert}{\sqrt{x}} \le \frac{\epsilon\sqrt{x}}{\sqrt{x}} = \epsilon
                \]
                \item We can therefore say that $(\forall \epsilon > 0)(\exists N \in \mathbb{N})(\forall n \ge N) \ \vert \sqrt{x_n} - \sqrt{x} \vert < \epsilon \Rightarrow \sqrt{x_n} \rightarrow \sqrt{x}$
            \end{itemize}
        \end{itemize}
    \end{itemize}
    \item Exercise 4.5. Show that limits, if they exist, must be unique. In other words, assume $\lim a_n = l_1$ and $\lim a_n = l_2$, and prove that $l_1 = l_2$
    \begin{itemize}
        \item Suppose for contradiction that $l_1 \neq l_2$. Let us assume without loss of generality that $l_1 < l_2$
        \item Then, consider $\epsilon = \frac{l_2 - l_1}{2} > 0$. Since $a_n \rightarrow l_1$,
        \[
        (\exists N_1 \in \mathbb{N})(\forall n \ge N_1) \ \vert a_n - l_1 \vert < \epsilon \Rightarrow l_1-\epsilon < a_n < l_1 + \epsilon
        \]
        \item In particular, $a_n < l_1 + \epsilon = \frac{l_1 + l_2}{2}$.
        \item Similarly, since $a_n \rightarrow l_2$,
        \[
        (\exists N_2 \in \mathbb{N})(\forall n \ge N_2) \ \vert a_n - l_2 \vert < \epsilon \Rightarrow l_2-\epsilon < a_n < l_2 + \epsilon
        \]
        \item In particular, $a_n > l_2 - \epsilon = \frac{l_1 + l_2}{2}$.
        \item Now set $N = \max(N_1,N_2)$. Then $\forall n \ge N$, we have that both
        \[
        a_n < \frac{l_1+l_2}{2} \text{ and } a_n > \frac{l_1+l_2}{2}
        \]
        must be true.
        \item This is clearly a contradiction. Therefore, by proof by contradiction, it must be the case that $l_1=l_2$.
    \end{itemize}
    \item Exercise 4.4. (Squeeze Theorem) Show that if $x_n \le y_n \le z_n$ for all $n \in \mathbb{N}$, and if $\lim x_n = \lim z_n = l$, then $\lim y_n = l$ as well.
    \begin{itemize}
        \item We start with the definitions:
        \[
        (\forall \epsilon > 0)(\exists N_1 \in \mathbb{N})(\forall n \ge N_1) \ \vert x_n - l \vert < \epsilon
        \]
        \item This means that $l - \epsilon < x_n < l + \epsilon$
        \[
        (\forall \epsilon > 0)(\exists N_2 \in \mathbb{N})(\forall n \ge N_2) \ \vert z_n - l \vert < \epsilon
        \]
        \item This means that $l - \epsilon < z_n < l + \epsilon$
        \item We know that $x_n \le y_n \le z_n$. Letting $N = \max(N_1,N_2)$, we have
        \[
        (\exists N \in \mathbb{N})(\forall n \ge N) \ l - \epsilon < x_n \le y_n \le z_n < l + \epsilon
        \]
        \item Hence, $(\forall n \ge N) \ l - \epsilon < y_n < l + \epsilon$.
        \item We conclude $(\forall \epsilon > 0)(\exists N \in \mathbb{N})(\forall n \ge N) \ \vert y_n - l \vert < \epsilon \Rightarrow y_n \rightarrow l$
    \end{itemize}
    \item Quick Trick Question: Recall the definition of convergence. What would happen if we changed the definition so that it ended $\vert a_n - l \vert \le \epsilon$?
    \begin{itemize}
        \item Nothing happens! The two definitions are equivalent.
        \item Clearly $\vert a_n - l \vert < \epsilon \Rightarrow \vert a_n - l \vert \le \epsilon$
        \item For the other direction, let $\epsilon' = \frac{\epsilon}{2} > 0$. So, $(\exists N \in \mathbb{N})(\forall n \ge N) \ \vert a_n - l \vert \le \epsilon' = \frac{\epsilon}{2} < \epsilon$
        \item That is, $\vert a_n - l \vert \le \epsilon \Rightarrow \vert a_n - l \vert < \epsilon$.
    \end{itemize}
    \item So why do we use the strict inequality?
    \begin{itemize}
        \item It turns out that open sets are nicer to work with in general, especially when doing topology (where we throw away the notion of distance but keep the notion of an open interval). Openness corresponds to having a strict inequality so we like to use strict inequalities in this (and other similar definitions).
    \end{itemize}
\end{itemize}

\section{Lecture 5 (28/10/2025)}
\begin{itemize}
    \item Definition - Boundedness
    \begin{itemize}
        \item A sequence $(a_n)$ is bounded if $(\exists L,U \in \mathbb{R})(\forall n \in \mathbb{N})L \le a_n \le U$
    \end{itemize}
    \item Proposition - a sequence $(a_n)$ is bounded if and only if $(\exists C \in \mathbb{R})(\forall n \in \mathbb{N}) \vert a_n \vert \le C$
    \begin{itemize}
        \item Proof: We assume that $(\exists C \in \mathbb{R}) (\forall n \in \mathbb{N}) \vert a_n \vert \le C$
        \item Now,
        \begin{flalign*}
            &(\exists C \in \mathbb{R})(\forall n \in \mathbb{N}) \ \vert a_n \vert \le C \\
            \Longleftrightarrow &(\exists C \in \mathbb{R})(\forall n \in \mathbb{N}) \ -C \le a_n \le C
        \end{flalign*}
        \item By setting $L=-C$ and $U=C$, we can see that
        \[ (\exists L,U \in \mathbb{R})(\forall n \in \mathbb{N}) \ L \le a_n \le U \]
        \item Now we assume
        \[ (\exists L,U \in \mathbb{R})(\forall n \in \mathbb{N}) \ L \le a_n \le U \]
        \item We set $C=\max\{\vert L \vert, \vert U \vert \}$
        \begin{flalign*}
            &(\exists L,U \in \mathbb{R})(\forall n \in \mathbb{N}) \ L \le a_n \le U \\
            \Rightarrow &(\exists L,U \in \mathbb{R})(\forall n \in \mathbb{N}) \ \vert a_n \vert \le C
        \end{flalign*}
    \end{itemize}
    \item Proposition - If $(a_n)$ is a convergent sequence, then $(a_n)$ is bounded
    \begin{itemize}
        \item We want to show that
        \begin{flalign*}
            &(\forall \epsilon > 0)(\exists N \in \mathbb{N})(\forall n \ge N) \vert a_n - a \vert < \epsilon \\
            \Rightarrow &(\exists L,U \in \mathbb{N})(\forall n \in \mathbb{N}) \ L \le a_n \le U
        \end{flalign*}
        \item Set up as follows:
        \begin{flalign*}
            &(\forall \epsilon > 0)(\exists N \in \mathbb{N})(\forall n \ge N) \vert a_n - a \vert < \epsilon \\
            \Rightarrow &(\forall \epsilon > 0)(\exists N \in \mathbb{N})(\forall n \ge N) a - \epsilon < a_n < a+ \epsilon
        \end{flalign*}
        \item Let us set $\epsilon = 1$:
        \begin{flalign*}
            (\forall n \ge N) a - 1 < a_n < a+ 1
        \end{flalign*}
        \item Now for $n < N_1$, we need to consider $\{a_1, \dots, a_{N_1-1}\}$, which will have a maximum element $M$ and minimum element $m$. Therefore,
        \[ (\forall n \le N_1 - 1) \ m \le a_n \le M \]
        \item Hence,
        \begin{flalign*}
            &(\forall n \ge N_1) \ a-1 \le a_n \le a+1 \wedge (\exists m,M \in \mathbb{R})(\forall n \le N_1 - 1) \ m \le a_n \le M \\
            \Rightarrow &(\exists m,M \in \mathbb{R})(\forall n \in \mathbb{N}) \ \min\{a-1,m\} \le a_n \le \max \{ a+1, M \} \\
            \Rightarrow &(\exists L,U \in \mathbb{R})(\forall n \in \mathbb{N}) \ L \le a_n \le U \\
        \end{flalign*}
    \end{itemize}
    \item Theorem (Algebraic Limit Theorem) - Let $\lim a_n = a$ and $\lim b_n = b$. Then,
    \begin{itemize}
        \item $(\forall c \in \mathbb{R}) \lim (ca_n) = ca$
        \begin{itemize}
            \item We know that
            \[ (\forall \epsilon > 0)(\exists N \in \mathbb{N})(\forall n \ge N) \vert a_n - a \vert < \epsilon \]
            \item Now we exploit the power of the fact that above statement works $\forall \epsilon > 0$
            \begin{flalign*}
                &(\forall \frac{\epsilon}{\vert c \vert} > 0)(\exists N \in \mathbb{N})(\forall n \ge N) \vert a_n - a \vert < \frac{\epsilon}{\vert c \vert} \\
                \Rightarrow &(\forall \epsilon > 0)(\exists N \in \mathbb{N})(\forall n \ge N) \vert c \vert \vert a_n - a \vert < \epsilon
            \end{flalign*}
            \item Now we note that $\vert c \vert \vert a_n - a \vert = \vert c(a_n-a) \vert$, because the product of two real numbers, ignoring signs before multiplication is the same as when we ignore the signs after multiplication.
            \begin{flalign*}
                &(\forall \epsilon > 0)(\exists N \in \mathbb{N})(\forall n \ge N) \vert c \vert \vert a_n - a \vert < \epsilon \\
                \Longleftrightarrow &(\forall \epsilon > 0)(\exists N \in \mathbb{N})(\forall n \ge N) \vert c(a_n - a) \vert < \epsilon \\
                \Longleftrightarrow &(\forall \epsilon > 0)(\exists N \in \mathbb{N})(\forall n \ge N) \vert ca_n - ca \vert < \epsilon \\
                \Longleftrightarrow &ca_n \rightarrow ca
            \end{flalign*}
        \end{itemize}
        \item $\lim (a_n+b_n)=a+b$
        \begin{itemize}
            \item We know that
            \begin{flalign*}
                &(\forall \epsilon > 0)(\exists N_a \in \mathbb{N})(\forall n \ge N_a) \vert a_n - a \vert < \frac{\epsilon}{2} \\
                \wedge &(\forall \epsilon > 0)(\exists N_b \in \mathbb{N})(\forall n \ge N_b) \vert b_n - b \vert < \frac{\epsilon}{2} \\
            \end{flalign*}
            \item Setting $N = \max(N_a,N_b)$, we see that
            \begin{flalign*}
                &(\forall \epsilon > 0)(\exists N \in \mathbb{N})(\forall n \ge N) \vert b_n - b \vert < \frac{\epsilon}{2} \wedge \vert a_n - a \vert < \frac{\epsilon}{2} \\
                \Rightarrow &(\forall \epsilon > 0)(\exists N \in \mathbb{N})(\forall n \ge N) \vert a_n - a \vert + \vert b_n - b \vert < \epsilon
            \end{flalign*}
            \item Now, we know from the Triangle Inequality, we have
            \[ \vert a_n + b_n - (a+b) \vert = \vert a_n - a + b_n - b \vert \le \vert a_n - a \vert + \vert b_n - b \vert \]
            \item Therefore, by transitivity, 
            \[ (\forall \epsilon > 0)(\exists N \in \mathbb{N})(\forall n \ge N) \vert a_n + b_n - (a+b) \vert < \epsilon \Longleftrightarrow a_n + b_n \rightarrow a + b \]
        \end{itemize}
        \item $\lim (a_nb_n) = ab$
        \begin{itemize}
            \item We have
            \begin{flalign*}
                &(\forall \epsilon_a > 0)(\exists N_a \in \mathbb{N})(\forall n \ge N_a) \vert a_n - a \vert < \epsilon_a \\
                \wedge &(\forall \epsilon_b > 0)(\exists N_b \in \mathbb{N})(\forall n \ge N_b) \vert b_n - b \vert < \epsilon_b \\
            \end{flalign*}
            \item We note that
            \begin{flalign*}
                \vert a_nb_n - ab \vert &= \vert a_nb_n - ab_n + ab_n - ab \vert \\
                &\le \vert a_nb_n - ab_n \vert + \vert ab_n - ab \vert \qquad \text{Triangle Inequality} \\
                &= \vert b_n \vert \vert a_n - a \vert + \vert a \vert \vert b_n - b \vert \\
            \end{flalign*}
            \item We also know that if $b_n$ is convergent, then it is bounded, i.e.
            \[ (\exists C \in \mathbb{R})(\forall n \in \mathbb{N}) \vert b_n \vert \le C \]
            \item Therefore,
            \[ \vert a_nb_n - ab \vert \le C\vert a_n-a \vert + \vert a \vert \vert b_n-b \vert\]
            \item Therefore, setting $M = \max\{N_a,N_b\}$, we see that 
            \[ (\forall \epsilon_a, \epsilon_b > 0)(\exists M \in \mathbb{N})(\forall n \ge M) \ \vert a_nb_n - ab \vert < C\epsilon_a + \vert a \vert \epsilon_b \]
            \item We now choose
            \[ \epsilon_a = \frac{1}{2C+1} \epsilon < \frac{1}{2C} \epsilon\]
            and
            \[ \epsilon_b = \frac{1}{2\vert a \vert + 1} \epsilon < \frac{1}{2\vert a \vert} \epsilon \]
            Therefore,
            \[ (\forall \epsilon > 0)(\exists M \in \mathbb{N})(\forall n \ge M) \ \vert a_nb_n - ab \vert < C \frac{1}{2C+1}\epsilon + \vert a \vert \frac{1}{2\vert a \vert +1} \epsilon < \frac{\epsilon}{2} + \frac{\epsilon}{2} \]
            and so
            \[ (\forall \epsilon > 0)(\exists M \in \mathbb{N})(\forall n \ge M) \vert a_nb_n - ab \vert < \epsilon \Rightarrow a_nb_n \rightarrow ab \]
        \end{itemize}
        \item $\lim(a_n/b_n) = a/b, b \neq 0$
        \begin{itemize}
            \item We just need to prove that
            \[ (\forall b \neq 0) \left( b_n \rightarrow b \Rightarrow \frac{1}{b_n} \rightarrow \frac{1}{b} \right) \]
            and then use $\lim a_nb_n = ab$
            \item We assume that
            \[( \forall \epsilon > 0)(\exists N_b \in \mathbb{N})(\forall n \ge N_b) \vert b_n - b \vert < \epsilon \]
            \item Now
            \[ \left\vert \frac{1}{b_n} - \frac{1}{b} \right\vert = \frac{\vert b-b_n \vert}{\vert b \vert \vert b_n \vert} \]
            \item From the fact that $b_n \rightarrow b$, we have $(\exists N_1 \in \mathbb{N})(\forall n \ge N_1) \vert b_n - b \vert < c$, where $0 < c < \vert b \vert$. Now the below proves that $\vert b_n - b \vert < c \Rightarrow \vert b_n \vert > \vert b \vert - c > 0$. 
            \begin{flalign*}
                &\text{Case 1.} b \ge 0 \\
                &\vert b \vert = b \\
                &b-c < b_n < b+c \\
                &b_n \le \vert b_n \vert \\
                &\vert b_n \vert > b - c = \vert b \vert - c \\
                &\text{Case 2. } b < 0 \\
                &\vert b \vert = -b \\
                &b-c < b_n < b+c \\
                &-b+c > b_n > -b-c \Rightarrow b_n > \vert b \vert - c\\
                &b_n \le \vert b_n \vert \\
                &\vert b_n \vert > \vert b \vert - c \\
            \end{flalign*}
            Therefore,
            \[ (\forall c \in (0,\vert b \vert))(\exists N_1 \in \mathbb{N})(\forall n \ge N_1) \ \frac{1}{\vert b_n \vert} < \frac{1}{\vert b \vert - c} \]
            \item Also, from the fact that $b_n \rightarrow b$, we have
            \[ (\forall \epsilon > 0)(\exists N_2 \in \mathbb{N})(\forall n \neq N_2) \ \vert b_n - b \vert < \vert b \vert (\vert b \vert - c) \epsilon \]
            \item By setting $N = \max\{N_1,N_2\}$, we see that
            \[ (\forall \epsilon > 0)(\exists N \in \mathbb{N})(\forall n \ge N) \frac{\vert b-b_n \vert}{\vert b \vert \vert b_n \vert} < (\vert b \vert - c)\epsilon \frac{1}{\vert b \vert - c} = \epsilon\]
            and so
            \[ (\forall \epsilon > 0)(\exists N \in \mathbb{N})(\forall n \ge N) \left\vert \frac{1}{b_n} - \frac{1}{b} \right\vert < \epsilon\]
        \end{itemize}
    \end{itemize}
\end{itemize}

\section{Lecture 6 (30/10/2025)}
\begin{itemize}
    \item Theorem (Sequence Squeeze Theorem)
    \[ ((\forall n \in \mathbb{N}) a_n \le x_n \le b_n) \wedge a_n \rightarrow L \wedge b_n \rightarrow L \Rightarrow x_n \rightarrow L \]
    \begin{itemize}
        \item We know that
        \begin{flalign*}
            &(\forall \epsilon > 0)(\exists N_a \in \mathbb{N})(\forall n \ge N_a) \ \vert a_n - L \vert < \epsilon \\
            \wedge &(\forall \epsilon > 0)(\exists N_b \in \mathbb{N})(\forall n \ge N_b) \ \vert b_n - L \vert < \epsilon \\
            \Longleftrightarrow &(\forall \epsilon > 0)(\exists N_a \in \mathbb{N})(\forall n \ge N_a) L-\epsilon < a_n < L+\epsilon \\
            \wedge &(\forall \epsilon > 0)(\exists N_b \in \mathbb{N})(\forall n \ge N_b) L-\epsilon < b_n < L+\epsilon \\
        \end{flalign*}
        \item Setting $N = \max \{N_a,N_b\}$, and noting that $(\forall n \in \mathbb{N}) \ a_n \le x_n \le b_n$, we see that
        \begin{flalign*}
            &(\forall \epsilon > 0)(\exists N \in \mathbb{N})(\forall n \ge N) L-\epsilon < a_n \le x_n \le b_n < L+\epsilon \\
            \Rightarrow &(\forall \epsilon > 0)(\exists N \in \mathbb{N})(\forall n \ge N) L-\epsilon < x_n < L+\epsilon \\
            \Longleftrightarrow &(\forall \epsilon > 0)(\exists N \in \mathbb{N})(\forall n \ge N) \vert x_n - L \vert < \epsilon \\
            \Longleftrightarrow &x_n \rightarrow L
        \end{flalign*}
    \end{itemize}
    \item Prove that $\frac{1}{n^3+\sqrt{n}+\pi} \rightarrow 0$
    \begin{itemize}
        \item $(\forall n \in \mathbb{N}) 0 \le \frac{1}{n^3+\sqrt{n}+\pi} \le \frac{1}{n}$
        \item Define sequences $a_n=0$ and $b_n = \frac{1}{n}$
        \item Now $(\forall n \in \mathbb{N}) \ \vert a_n - 0 \vert = 0 \Rightarrow a_n \rightarrow 0$
        \item The natural numbers are unbounded by the Archimedes Principle, so $\forall \epsilon > 0$, we can pick $N > \frac{1}{\epsilon}$
        \item Now $N > \frac{1}{\epsilon} \Longleftrightarrow \frac{1}{N} < \epsilon$
        \item Hence, $\forall \epsilon > 0$, we can pick $N > \frac{1}{\epsilon}$ and $\forall n \ge N$, we have $b_n - 0 = \vert b_n - 0 \vert < \epsilon$
        \item Therefore, $b_n \rightarrow 0$
        \item Applying the Sequence Squeeze Theorem, we see that $\frac{1}{n^3+\sqrt{n}+\pi} \rightarrow 0$
    \end{itemize}
    \item The limiting process is also well-behaved with respect to the order operation
    \item Theorem (Order Limit Theorem) - Assume $\lim a_n = a$ and $\lim b_n = b$
    \begin{itemize}
        \item Prove: $((\forall n \in \mathbb{N}) a_n \ge 0) \Rightarrow a \ge 0$
        \begin{itemize}
            \item Using proof by contradiction, assume $a < 0$
            \item By the definition of convergence
            \[ (\forall \epsilon) (\exists N \in \mathbb{N})(\forall n \ge N) \ \vert a_n - a \vert < \epsilon \]
            \item Now draw real line with $a < 0$ and $0$
            \item When $\epsilon = \vert a \vert$, the definition of convergence tells that there exists $N_{\vert a \vert} \in \mathbb{N}$ such that $(\forall n \ge N_{\vert a \vert}) \ \vert a_n - a \vert < \vert a \vert$ From the above figure, we can see that this means that $(\forall N_{\vert a \vert}) \ a_n < 0$. This contradicts the assumption that $(\forall n \in \mathbb{N}) \ a_n \ge 0$
            \item Using symbols, we have
            \begin{flalign*}
                &(\forall \epsilon) (\exists N \in \mathbb{N})(\forall n \ge N) \ \vert a_n - a \vert < \epsilon \\
                \Rightarrow &(\exists N_{\vert a \vert} \in \mathbb{N})(\forall n \ge N_{\vert a \vert}) \ \vert a_n - a \vert < \epsilon = \vert a \vert \\
                \Rightarrow &(\exists N_{\vert a \vert} \in \mathbb{N})(\forall n \ge N_{\vert a \vert})) \ a_n < 0
            \end{flalign*}
            \item Without the figure, the last step is hard
        \end{itemize}
        \item Prove: $((\forall n \in \mathbb{N}) a_n \le b_n) \Rightarrow a \le b$
        \begin{itemize}
            \item Using proof by contradiction, assume $a > b$
            \item By the definition of convergence
            \[ (\forall \epsilon) (\exists N \in \mathbb{N})(\forall n \ge N) \ \vert a_n - a \vert < \epsilon \]
            \item Now draw real line with $a > b$
            \item When $\epsilon = \vert a-b \vert$, the definition of convergence tells that there exists $N_{\vert a \vert} \in \mathbb{N}$ such that $(\forall n \ge N_{\vert a \vert}) \ \vert a_n - a \vert < \vert a-b \vert$ From the above figure, we can see that this means that $(\forall N_{\vert a \vert}) \ a_n > b_n$. This contradicts the assumption that $(\forall n \in \mathbb{N}) \ a_n \le b_n$
            \item Using symbols, we have
            \begin{flalign*}
                &(\forall \epsilon) (\exists N \in \mathbb{N})(\forall n \ge N) \ \vert a_n - a \vert < \epsilon \\
                \Rightarrow &(\exists N_{\vert a \vert} \in \mathbb{N})(\forall n \ge N_{\vert a \vert}) \ \vert a_n - a \vert < \epsilon = \vert b-a \vert \\
                \Rightarrow &(\exists N_{\vert a \vert} \in \mathbb{N})(\forall n \ge N_{\vert a \vert})) \ a_n > b_n
            \end{flalign*}
        \end{itemize}
        \item Prove: $((\exists c \in \mathbb{R})(\forall n \in \mathbb{N}) \ c \le b_n) \Rightarrow c \le b$. Similarly, $((\forall n \in \mathbb{N}) \ a_n \le c) \Rightarrow a \le c$
        \begin{itemize}
            \item Using proof by contradiction, assume $c > b$
            \item By the definition of convergence
            \[ (\forall \epsilon) (\exists N \in \mathbb{N})(\forall n \ge N) \ \vert b_n - b \vert < \epsilon \]
            \item Now draw real line with $c > b$
            \item When $\epsilon = \vert c-b \vert$, the definition of convergence tells that there exists $N_{\vert b \vert} \in \mathbb{N}$ such that $(\forall n \ge N_{\vert b \vert}) \ \vert b_n - b \vert < \vert c-b \vert$ From the above figure, we can see that this means that $(\forall N_{\vert a \vert}) \ c > b_n$. This contradicts the assumption that $(\forall n \in \mathbb{N}) \ c \le b_n$
            \item Using symbols, we have
            \begin{flalign*}
                &(\forall \epsilon) (\exists N \in \mathbb{N})(\forall n \ge N) \ \vert b_n - b \vert < \epsilon \\
                \Rightarrow &(\exists N_{\vert b \vert} \in \mathbb{N})(\forall n \ge N_{\vert b \vert}) \ \vert b_n - b \vert < \epsilon = \vert c-b \vert \\
                \Rightarrow &(\exists N_{\vert b \vert} \in \mathbb{N})(\forall n \ge N_{\vert b \vert})) \ c > b_n
            \end{flalign*}
        \end{itemize}
    \end{itemize}
    \item Definitions
    \begin{itemize}
        \item A sequence $(a_n)$ is monotone increasing if $(\forall n \in \mathbb{N}) \ a_n \le a_{n+1}$
        \item A sequence $(a_n)$ is strictly monotone increasing if $(\forall n \in \mathbb{N}) \ a_n < a_{n+1}$
        \item A sequence $(a_n)$ is monotone decreasing if $(\forall n \in \mathbb{N}) \ a_n \ge a_{n+1}$
        \item A sequence $(a_n)$ is strictly monotone decreasing if $(\forall n \in \mathbb{N}) \ a_n > a_{n+1}$
        \item A sequence is monotone if it is either increasing or decreasing
        \item A sequence is strictly monotone if it is either strictly increasing or strictly decreasing
    \end{itemize}
    \item Show that the sequence $a_n = -\frac{1}{n}$ is strictly monotone increasing and converges to 0
    \begin{itemize}
        \item $(\forall n \in \mathbb{N}) \ a_{n+1} = -\frac{1}{n+1} > -\frac{1}{n} = a_n$
        \item Therefore $(a_n)$ strictly monotone increasing
        \item $\forall \epsilon > 0$ we can pick $N \in \mathbb{N}$ such that $N > \frac{1}{\epsilon} \Longleftrightarrow \forall \epsilon > 0$ we can pick $N \in \mathbb{N}$ such that $\frac{1}{N} < \epsilon$. Therefore,
        \[ (\forall \epsilon > 0)(\exists N \in \mathbb{N}) \left\vert - \frac{1}{n} \right\vert < \epsilon \]
        \item So, we have proved that $-\frac{1}{n} \rightarrow 0$
    \end{itemize}
    \item Theorem - Monotone Convergence Theorem
    \item If a sequence is monotone and bounded, then it converges
    \begin{itemize}
        \item A monotone sequence is either monotone increasing or monotone decreasing
        \item Case 1. Assume $(a_n)_n$ is monotone increasing and bounded
        \begin{itemize}
            \item We assume $(a_n)_n$ is monotone increasing and bounded. We exploit the boundedness of the sequence, which means it has a supremum, i.e., least upper bound property, which we label as $s$:
            \[ s = \sup \{ a_n : n \in \mathbb{N} \} \]
            \item We draw a picture of monotone increasing and bounded sequence with supremum $s$
            \item Using the definition of a supremum, we know that for all $\epsilon > 0, s-\epsilon$ is not an upper bound of the sequence $(a_n)_n$. Therefore, there exists $N \in \mathbb{N}$ such that $a_N > s-\epsilon$
            \[ (\exists s \in \mathbb{R}) s = \sup \{ a_n : n \in \mathbb{N} \} \Rightarrow (\forall \epsilon > 0)(\exists N \in \mathbb{N}) \ s - \epsilon < a_N \]
            \item Now, we know that because $(a_n)_n$ is monotone increasing, $(\forall n \ge N)(a_N \le a_n)$. We also know that $a_n$ must be less than or equal to the supremum $s$, and so $a_N \le a_n \le s < s + \epsilon$.
            \item Putting this all together, we have
            \begin{flalign*}
                &(a_n)_n \text{ is bounded and monotone increasing sequence} \\
                \Rightarrow &(\exists s \in \mathbb{R}) s = \sup \{ a_n : n \in \mathbb{N} \} \Rightarrow (\forall \epsilon > 0)(\exists N \in \mathbb{N}) \ s - \epsilon < a_N
            \end{flalign*}
            and
            \begin{flalign*}
                &(a_n)_n \text{ is bounded and monotone increasing sequence} \\
                \Rightarrow &(\exists s \in \mathbb{R}) s = \sup \{ a_n : n \in \mathbb{N} \} \Rightarrow (\forall \epsilon > 0)(\exists N \in \mathbb{N}) \ s - \epsilon < a_n < s + \epsilon \\
                \Longleftrightarrow &(\exists s \in \mathbb{R}) s = \sup \{ a_n : n \in \mathbb{N} \} \Rightarrow (\forall \epsilon > 0)(\exists N \in \mathbb{N}) \ \vert a_n - s \vert < \epsilon
            \end{flalign*}
            which tell us that $\lim a_n = s$. A monotone increasing and bounded sequence converges to its supremum.
        \end{itemize}
        \item Case 2. Assume $(a_n)_n$ is monotone decreasing and bounded
        \begin{itemize}
            \item We assume $(a_n)_n$ is monotone decreasing and bounded. We exploit the boundedness of the sequence, which means it has a infimum, i.e., greatest lower bound property, which we label as $b$:
            \[ b = \inf \{ a_n : n \in \mathbb{N} \} \]
            \item We draw a picture of monotone decreasing and bounded sequence with infimum $b$
            \item Using the definition of a infimum, we know that for all $\epsilon > 0, b+\epsilon$ is not a lower bound of the sequence $(a_n)_n$. Therefore, there exists $N \in \mathbb{N}$ such that $a_N < b+\epsilon$
            \[ (\exists b \in \mathbb{R}) b = \inf \{ a_n : n \in \mathbb{N} \} \Rightarrow (\forall \epsilon > 0)(\exists N \in \mathbb{N}) \ s - \epsilon < a_N \]
            \item Now, we know that because $(a_n)_n$ is monotone increasing, $(\forall n \ge N)(a_N \le a_n)$. We also know that $a_n$ must be less than or equal to the supremum $s$, and so $a_N \le a_n \le s < s + \epsilon$.
            \item Putting this all together, we have
            \begin{flalign*}
                &(a_n)_n \text{ is bounded and monotone increasing sequence} \\
                \Rightarrow &(\exists s \in \mathbb{R}) s = \sup \{ a_n : n \in \mathbb{N} \} \Rightarrow (\forall \epsilon > 0)(\exists N \in \mathbb{N}) \ s - \epsilon < a_N
            \end{flalign*}
            and
            \begin{flalign*}
                &(a_n)_n \text{ is bounded and monotone increasing sequence} \\
                \Rightarrow &(\exists s \in \mathbb{R}) s = \sup \{ a_n : n \in \mathbb{N} \} \Rightarrow (\forall \epsilon > 0)(\exists N \in \mathbb{N}) \ s - \epsilon < a_N < s + \epsilon \\
                \Longleftrightarrow &(\exists s \in \mathbb{R}) s = \sup \{ a_n : n \in \mathbb{N} \} \Rightarrow (\forall \epsilon > 0)(\exists N \in \mathbb{N}) \ \vert a_n - s \vert < \epsilon
            \end{flalign*}
            which tell us that $\lim b_n = b$. A monotone decreasing and bounded sequence converges to its infimum.
        \end{itemize}
    \end{itemize}
    \item Definition - Partial Sums and Series
    \begin{itemize}
        \item The $m$'th partial sum of the sequence $(b_n)$ is given by
        \[ s_m = \sum_{n=1}^m b_n \]
        and the sequences of partial sums $(s_m)$ is a series.
        \item We say that the series $\sum_{n=1}^\infty b_n$ converges to $b$ if the seqeunce $(s_m)$ converges to $b$
        \item $\sum_{n=1}^\infty b_n = b$
    \end{itemize}
    \item Series - sequence generated by a sum of a sequence
    \item Prove that the series $\sum_{n=1}^\infty \frac{1}{n^2}$ converges
    \begin{itemize}
        \item Define the partial sum:
        \[ s_m = \sum_{n=1}^m \frac{1}{n^2} \]
        \item We can see that $(\forall m \in \mathbb{N}) \ s_{m+1} > s_m$, and so $(s_m)$ is strictly monotone increasing.
        \item Proof that $(s_m)$ is bounded above (obviously bounded below by 1)
        \begin{flalign*}
            s_m &= 1 + \frac{1}{2 \cdot 2} + \frac{1}{3 \cdot 3} + \dots + \frac{1}{m \cdot m} \\
            &< 1 + \frac{1}{2 \cdot 1} + \frac{1}{3 \cdot 2} + \dots + \frac{1}{m(m-1)} \\
            &< 1 + \left( 1- \frac{1}{2} \right) + \left( \frac{1}{2} - \frac{1}{3} \right) + \dots + \left( \frac{1}{m-1} - \frac{1}{m} \right) \qquad \text{Cancels out (telescopic sum)} \\
            &= 1 + 1 - \frac{1}{m} < 2
        \end{flalign*}
        \item Thus, 2 is an upper bound for the sequence of partial sums, so by the Monotone Convergence Theorem, since the series $(s_m)$ is bounded and strictly monotonically increasing, the series $\sum_{n=1}^\infty \frac{1}{n^2}$ converges to some unknown limit less than 2
    \end{itemize}
\end{itemize}

\section{Tutorial 4 (30/10/2025)}
\begin{itemize}
    \item We can also argue that if $a_n$ is monotone decreasing, by definition, $(\forall n \in \mathbb{N}) a_n \ge a_{n+1}$, we can construct a sequence $-a_n = (-1) \times a_n$ that is monotone increasing $(\forall n \in \mathbb{N}) -a_n \le -a_{n+1}$, which converges, so $a_n$ converges by the Algebraic Limit Theorem.
    \item Exercise 4.15
    \begin{itemize}
        \item (a) Explain why $\sqrt{xy} \le \frac{x+y}{2}$ for any two positive real numbers $x$ and $y$.
        \begin{itemize}
            \item We do the same as before (refer to Tutorial 2):
            \[
            \sqrt{xy} \le \frac{x+y}{2} \Longleftrightarrow xy \le \frac{(x+y)^2}{4} \Longleftrightarrow (x-y)^2 \ge 0
            \]
            Since the last inequality is always true, and since each implication goes both ways, it follows that the first inequality is true as well.
        \end{itemize}
        \item (b) Now let $0 \le x_1 \le y_1$ and define
        \[
        x_{n+1} = \sqrt{x_ny_n} \text{ and } y_{n+1} = \frac{x_n+y_n}{2}
        \]
        Show that $\lim x_n$ and $\lim y_n$ both exist and are equal.
        \begin{itemize}
            \item With recurrence relations like this, we can often apply the Monotone Convergence Theorem. We need to show that $x_n$ and $y_n$ are both monotone and bounded.
            \item From (a), we have that $(\forall n \ge 1) \ x_{n+1} \le y_{n+1}$. Including that $0 \le x_1 \le y_1$, we see that $(\forall n \ge \mathbb{N}) \ x_n \le y_n$
            \item So we immediately see that $\forall n \in \mathbb{N}$, $x_{n+1} = \sqrt{x_ny_n} \ge \sqrt{x_n^2} = x_n$ (since $x_n \ge 0$).
            \item This means that $x_n$ is monotone increasing
            \item Likewise, $y_{n+1} - y_n = \frac{x_n - y_n}{2} \le 0$ and so $(\forall n \in \mathbb{N})$
            \item This means that $y_n$ is monotone decreasing
            \item Now we need to show that both sequences are bounded
            \item Note that $\forall n \in \mathbb{N}$, we have shown that $0 \le x_1 \le x_n \le y_n \le y_1$
            \item So by induction, the sequence $x_n$ is bounded above by $y_1$ and the sequence $y_n$ is bounded below by $x_1$.
            \item By the Monotone Convergence Theorem, since both $x_n$ and $y_n$ are monotone and bounded, they both converge
            \item Suppose $x_n \rightarrow x$ and $y_n \rightarrow y$. Then, taking the limit of both sides of $y_{n+1} = \frac{x_n + y_n}{2}$, and using the algebra of limits, we get
            \[
            y = \frac{x+y}{2}
            \]
            which implies that $x=y$.
        \end{itemize}
    \end{itemize}
\end{itemize}

\section{Lecture 7 (04/11/2025)}
\begin{itemize}
    \item Subsequence - Let $(a_n)$ be a sequence of real numbers, and let $n_1 < n_2 < n_3 \cdots$ be a strictly increasing sequence of natural numbers. Then the sequence $(a_{n_1},a_{n_2},a_{n_3}, \dots)$ is called a subsequence of $(a_n)$ and is denoted by $(a_{n_k})$ where $k \in \mathbb{N}$ indexes the subsequence
    \begin{itemize}
        \item A nice way to think about subsequences is to start from a sequence $a(n): \mathbb{N} \rightarrow \mathbb{R}$ and a strictly monotone increasing sequence $m(n): \mathbb{N} \rightarrow \mathbb{N}$. Then $a(m(n))$ is a subsequence of $a(n)$.
        \item It is also important to see that because $m$ maps from $\mathbb{N}$ to $\mathbb{N}$, $m(n)$ will not converge.
    \end{itemize}
    \item Proposition - $a_n \rightarrow a \Longleftrightarrow \text{(for every subsequence } a_{n_k}) \ a_{n_k} \rightarrow a$
    \begin{itemize}
        \item Prove: $a_n \rightarrow a \Rightarrow \text{(for every subsequence } a_{n_k}) \ a_{n_k} \rightarrow a$
        \begin{itemize}
            \item We start by assuming $\lim a_n = a$
            \[ (\forall \epsilon > 0)(\exists N \in \mathbb{N})(\forall n \ge N) \ \vert a_n - a \vert < \epsilon \]
            \item Now, for every monotone increasing sequence $m : \mathbb{N} \rightarrow \mathbb{N}$ such that $m(n+1) \ge m(n)$, using the Archimedean property of natural numbers, we can find an $M$ such that $m(M) > N$
            \item Therefore,
            \[ (\forall \epsilon > 0)(\exists M \in \mathbb{N})(\forall n \ge M) \ \vert a(m(n)) - a \vert < \epsilon \]
            \item Hence, $a_{m_n} \rightarrow a$
        \end{itemize}
        \item Prove: $\text{(for every subsequence } a_{n_k}) \ a_{n_k} \rightarrow a \Rightarrow a_n \rightarrow a$
        \begin{itemize}
            \item By the definition of subsequences, suppose for every $m : \mathbb{N} \rightarrow \mathbb{N}$ such that $m(n+1) \ge m(n)$,
            \[ (\forall \epsilon > 0)(\exists M \in \mathbb{N})(\forall n \ge M) \ \vert a(m(n)) - a \vert < \epsilon \]
            \item If this is true for every $m$, we can just choose $m(n)=n$ and get
            \[ (\forall \epsilon > 0)(\exists N \in \mathbb{N})(\forall n \ge N) \ \vert a_n - a \vert < \epsilon \]
            \item Hence, $a_n \rightarrow a$
            \item In other words, since $a_n$ is a subsequence of itself, $\text{(for every subsequence } a_{n_k}) \ a_{n_k} \rightarrow a \Rightarrow a_n \rightarrow a$ is obvious
        \end{itemize}
       
    \end{itemize}
    \item Peak - $a_m$ is a peak of the sequence $(a_n)$ if $(\forall n > m) \ a_m \ge a_n$
    \item Lemma for BW Theorem - Every sequence contains a monotone subsequence
    \begin{itemize}
        \item A sequence will either have a fixed number of peaks, say $N \in \mathbb{N}$ or the number of peaks, like the natural numbers, will never end
        \item If the number of peaks never end, let $a_{n_k}$ be the $k$th peak. Then for every $k \in \mathbb{N}$, we have $a_{n_k} \ge a_{n_{k+1}}$. Hence $(a_{n_k})$ is a monotone decreasing subsequence.
        \item Now we consider the case where there are a fixed number of peaks. Without loss of generality, let $a_K$ be the final peak. Now, $(\forall n \ge K) \ a_n \le a_K$. Define $a_{n_1} = a_{K+1}$ and let $a_{n_2}$ be the next term in the sequence which is greater than or equal to $a_{n_1}$. Such a term must exist, because $a_{n_1}$ is not a peak. Similarly, let $a_{n_3}$ be the next term in the sequence which is greater than or equal to $a_{n_2}$, which must exists, because $a_{n_2}$ is not a peak. Continuing in this way, we construct a monotone increasing subsequence.
    \end{itemize}
    \item Theorem (Bolzano-Weierstrass Theorem) - Every bounded sequence contains a convergent subsequence
    \begin{itemize}
        \item Proof: from the Lemma for BW Theorem, we know that every bounded sequence contains a bounded and monotone subsequence, which by the Monotone Convergence Theorem, must converge
    \end{itemize}
    \item Cauchy Sequence - sequence $(a_n)$ that is
    \[
    (\forall \epsilon > 0)(\exists N \in \mathbb{N})(\forall m,n \ge N) \vert a_n - a_m \vert < \epsilon
    \]
    \item Theorem - Every convergent sequence is a Cauchy sequence
    \begin{itemize}
        \item Proof.
        \item Assume $(x_n) \rightarrow x$
        \[ (\forall \epsilon > 0)(\exists N \in \mathbb{N})(\forall n \ge N) \vert x_n - x \vert < \epsilon \]
        \item Now we note that
        \begin{flalign*}
            &(\forall \epsilon > 0)(\exists N \in \mathbb{N})(\forall n \ge N) \vert x_n - x \vert < \frac{\epsilon}{2} \\
            \Rightarrow &(\forall \epsilon > 0)(\exists N \in \mathbb{N})(\forall n,m \ge N) \vert x_n - x \vert < \frac{\epsilon}{2} \wedge \vert x_m - x \vert < \frac{\epsilon}{2} \\
            \Rightarrow &(\forall \epsilon > 0)(\exists N \in \mathbb{N})(\forall n,m \ge N) \vert x_n - x \vert + \vert x_m - x \vert < \epsilon \\
        \end{flalign*}
        \item From the Triangle Inequality, we know that $\vert x + y \vert \le \vert x \vert + \vert y \vert$. Hence,
        \begin{flalign*}
            &(\forall \epsilon > 0)(\exists N \in \mathbb{N})(\forall n,m \ge N) \vert x_n - x_m \vert \le \vert x_n - x \vert + \vert -(x_m - x) \vert = \vert x_n - x \vert + \vert x_m - x \vert < \epsilon \\
            \Rightarrow &(\forall \epsilon > 0)(\exists N \in \mathbb{N})(\forall n,m \ge N) \vert x_n - x_m \vert < \epsilon
        \end{flalign*}
    \end{itemize}
    \item Lemma - Cauchy sequences are bounded
    \begin{itemize}
        \item Suppose $(x_n)$ is a Cauchy sequence, y.e.
        \[
        (\forall \epsilon > 0)(\exists N \in \mathbb{N})(\forall n,m \ge N) \ \vert x_n - x_m \vert < \epsilon
        \]
        \item Setting $\epsilon = 1$, we obtain
        \[
        (\exists N \in \mathbb{N})(\forall n,m \ge N) \ \vert x_n - x_m \vert < 1
        \]
        \item We now set $m = N$ and so
        \begin{flalign*}
            &(\exists N \in \mathbb{N})(\forall n \ge N) \ \vert x_n - x_N \vert < 1 \\
            \Rightarrow &(\exists N \in \mathbb{N})(\forall n \ge N) \ x_N - 1 < x_n < x_N + 1 \\ 
        \end{flalign*}
        \item Therefore,
        \[
        (\forall n \in \mathbb{N}) L \le x_n \le U
        \]
        where
        \begin{flalign*}
            &L = \min(\{ x_1, \dots, x_{N-1},x_N - 1\}) \\
            &U = \max(\{x_1, \dots, x_{N-1},x_N + 1 \})
        \end{flalign*}
    \end{itemize}
    \item Theorem (Cauchy Criterion) - A sequence converges if any only if it is a Cauchy sequence
    \[
    (\forall \epsilon > 0)(\exists N \in \mathbb{N})(\forall n \ge N) \ \vert x_n - x \vert < \epsilon \Longleftrightarrow (\forall \epsilon > 0)(\exists N \in \mathbb{N})(\forall n,m \ge N) \ \vert x_n - x_m \vert < \epsilon 
    \]
    \begin{itemize}
        \item Proof.
        \item We already know that if a sequence is convergent, then it is a Cauchy sequence
        \item It remains to prove that if a sequence is a Cauchy sequence, then it converges
        \item Suppose $(x_n)$ is a Cauchy seqeunce. We know that every Cauchy sequence is bounded, so by the Bolzano-Weierstrass Theorem, we see that every Cauchy sequence has a convergent subseqeunce
        \[
        (\exists x \in \mathbb{R}) \ \lim x_{n_k} = x
        \]
        \item We know that
        \[
        (\forall \epsilon > 0)(\exists N \in \mathbb{N})(\forall n,m \ge N) \ \vert x_n - x_m \vert < \frac{\epsilon}{2}
        \]
        and
        \[
        (\forall \epsilon > 0)(\exists M \in \mathbb{N}) (\forall n_k \ge M) \ \vert x_{n_k} - x \vert < \frac{\epsilon}{2}
        \]
        \item Therefore, by setting $K = \max(N,M)$, we know
        \begin{flalign*}
            &(\forall \epsilon > 0)(\exists K \in \mathbb{N})(\forall n,n_k \ge K) \ \vert x_n - x_{n_k} \vert < \frac{\epsilon}{2} \wedge  \ \vert x_{n_k} - x \vert < \frac{\epsilon}{2} \\
            &(\forall \epsilon > 0)(\exists K \in \mathbb{N})(\forall n,n_k \ge K) \ \vert x_n - x_{n_k} \vert + \vert x_{n_k} - x \vert < \epsilon
        \end{flalign*}
        \item From the Triangle Inequality, $\vert x_n - x \vert \le \vert x_n - x_{n_k} \vert + \vert x_{n_k} - x \vert$, and so
        \[
        (\forall \epsilon > 0)(\exists K \in \mathbb{N})(\forall n \ge K) \ \vert x_n - x \vert < \epsilon
        \]
    \end{itemize}
\subsection{Slide 4 (Linear Algebra: Vectors and Matrices)}
    \item Linear algebra is the study of vectors and the rules to manipulate them
    \item Traditional "geometric vectors" are denoted $\underline{x}$
    \item General vectors and be added and scaled by scalars
    \[ \lambda \underline{x}, \qquad \underline{x} + \underline{y} \]
    \item Examples of vectors
    \begin{itemize}
        \item Geometric vectors
        \begin{itemize}
            \item Represented by arrows with magnitude and direction
            \item Operations:
            \[ \underline{x} + \underline{y}, \quad \lambda \underline{x}, \quad \lambda \in \mathbb{R} \]
        \end{itemize}
        \item Polynomials
        \begin{itemize}
            \item Can be added and scaled:
            \[ \sum_{n=0}^N a_nx^n + \sum_{n=0}^N b_nx^n = \sum_{n=0}^N (a_n + b_n)x^n, \quad \lambda \sum_{n=0}^N a_nx^n = \sum_{n=0}^N \lambda a_nx^n \]
        \end{itemize}
        \item Audio signals
        \begin{itemize}
            \item Represented as sequences of numbers
            \item Operations:
            \[ \text{Addition: Combine signals, \quad    Scaling: Amplify/attenuate} \]
        \end{itemize}
        \item Elements of $\mathbb{R}^n$
        \begin{itemize}
            \item Tuples of real numbers, e.g.:
            \[ \underline{a} = \left[ \begin{array}{c}
                1 \\
                2 \\ 
                3 \\
            \end{array} \right], \quad \underline{b} = \left[ \begin{array}{c}
                4 \\
                5 \\ 
                6 \\
            \end{array} \right] \]
        \end{itemize}
    \end{itemize}
    \item Simple Linear Regression
    \begin{itemize}
        \item Given points $(x_1,y_1),(x_2,y_2),(x_3,y_3)$, fit a straight line:
        \begin{flalign*}
            y_i &= \beta_1 + \beta_2 x_i + \epsilon_i \\
            \epsilon_i &= y_i - (\beta_1 + \beta_2 x_i)
        \end{flalign*}
        \item Minimize Mean Squared Error (MSE):
        \begin{flalign*}
        &\text{Minimize: } \frac{1}{3} \sum_{i=1}^3 \epsilon_i^2 \\
            \frac{\partial}{\partial \beta_1} \sum_{i=1}^3 \epsilon_i^2 &= -\sum_{i=1}^3 2y_i + \sum_{i=1}^3 2 (\beta_1 + \beta_2 x_i) = 0 \\
            \frac{\partial}{\partial \beta_2} \sum_{i=1}^3 \epsilon_i^2 &= -\sum_{i=1}^3 2x_iy_i + \sum_{i=1}^3 2 (\beta_1 + \beta_2 x_i)x_i = 0 \\
            -\frac{1}{3} \sum_{i=1}^3 y_i + \beta_1 + \beta_2 \frac{1}{3} \sum_{i=1}^3 x_i &= 0 \\
            -\frac{1}{3} \sum_{i=1}^3 x_iy_i + \beta_1 + \beta_2 \frac{1}{3} \sum_{i=1}^3 x_i^2 &= 0 \\ \\
            \text{Define:}& \\
            \overline{x} &= \frac{1}{3} \sum_{i=1}^3 x_i \\
            \overline{y} &= \frac{1}{3} \sum_{i=1}^3 y_i \\
            \overline{xy} &= \frac{1}{3} \sum_{i=1}^3 x_iy_i \\
            \overline{x^2} &= \frac{1}{3} \sum_{i=1}^3 x_i^2 \\
            \frac{1}{3} \sum_{i=1}^3 y_i &= \beta_1 + \beta_2 \frac{1}{3} \sum_{i=1}^3 x_i \\
            \frac{1}{3} \sum_{i=1}^3 x_iy_i &= \beta_1 + \beta_2 \frac{1}{3} \sum_{i=1}^3 x_i^2 
        \end{flalign*}
        \item Use calculus to derive normal equations:
        \begin{flalign*}
            \overline{y} &= \beta_1 + \beta_2 \overline{x} \\
            \overline{xy} &= \overline{x} \beta_1 + \overline{x^2} \beta_2
        \end{flalign*}
        \item Application: Multivariate regression, geometry of regression via matrices
    \end{itemize}
\end{itemize}

\section{Tutorial 5 (06/11/2025)}
\begin{itemize}
    \item Focus on Exercises 1-9
    \item Exercise 4.1. Let $x_1 = 2$, and define
    \[
    x_{n+1} = \frac{1}{2} \left( x_n + \frac{2}{x_n} \right)
    \]
    \begin{itemize}
        \item (a) Show that $x_n^2$ is always greater than 2, and then use this to prove that $x_n - x_{n+1} \ge 0$. Conclude that $\lim x_n = \sqrt{2}$.
        \item (ai) Show that $x_n^2$ is always greater than 2
        \begin{itemize}
            \item If you are proving $x_n^2 > 2$, you consider $x_n^2 - 2 > 0$
            \begin{flalign*}
                x_{n+1}^2 - 2 &= \frac{1}{4} \left( x_n + \frac{2}{x_n} \right)^2 - 2 = \frac{1}{4} \left( x_n^2 + 4 + \frac{4}{x_n^2} \right) - 2 = \\
                &= \frac{1}{4x_n^2} \left( x_n^4 + 4x_n^2 +4 - 8x_n^2\right) = \frac{1}{4x_n^2} \left( x_n^4 - 4x_n^2 +4\right) = \frac{1}{4x_n^2} (x_n^2 - 2)^2 \ge 0
            \end{flalign*}
            \item Note that $x_n \neq 0$ for any $n$ from the above recursive formula.
            \item Therefore, $(\forall n \in \mathbb{N}) \ x_{n+1}^2 \ge 2$ and since $x_1^2 = 4 \ge 2, (\forall n \in \mathbb{N}) \ x_n^2 > 2$ by induction.
            \item If we want to show $x_n^2 > 2$ (strict inequality), then we can employ induction with the relationship (from the previous slide).
            \[
            x_{n+1}^2 - 2 = \frac{1}{4x_n^2} (x_n^2 - 2)^2
            \]
            \item Note that $x_1^2 = 4 > 2,$ and since $x_n^2 > 2 \Rightarrow (\forall n \in \mathbb{N}) \ x_{n+1}^2 - 2 > 0$ by induction.
        \end{itemize}
        \item (aii) Prove that $x_n - x_{n+1} \ge 0$ (sequence is monotone decreasing)
        \begin{itemize}
            \item Now consider
            \[
            x_n - x_{n+1} = x_n - \frac{1}{2} \left( x_n + \frac{2}{x_n} \right) = \frac{1}{2} \left( x_n - \frac{2}{x_n} \right) = \frac{x_n^2 - 2}{2x_n}
            \]
            \item Since $x_n^2 > 2$, $2x_n > 0$ and $x_n^2 - 2 > 0$. Therefore, both the numerator and denominator is positive.
            \item Therefore, $(\forall n \in \mathbb{N}) \ x_n - x_{n+1} \ge 0$
        \end{itemize}
        \item (aiii) Conclude that $\lim x_n = \sqrt{2}$.
        \begin{itemize}
            \item Since $x_n - x_{n+1} \ge 0$, it follows that $(\forall n \in \mathbb{N}) \ x_{n+1} \le x_n$. so $x_n$ is a monotone decreasing sequence.
            \item Moreover, $x_n^2 > 2$, so $x_n$ is bounded from below by $\sqrt{2}$. It is also bounded from above by 2 since $x_1 = 2$ and the sequence is monotone decreasing.
            \item Since the sequence is bounded by $2 > x_n > \sqrt{2}$ and it is monotone decreasing, by the Monotone Convergence Theorem, $x_n$ converges
            \item Suppose that $x_n \rightarrow x$. Then, taking the limit of both sides of the recurrence relation above, and using the algebra of limits, we obtain that $x = \frac{1}{2} \left( x + \frac{2}{x} \right) \Rightarrow x = \sqrt{2}$
            \item Note that the limit can't be $-\sqrt{2}$ because we argued that $x_n > 0, \ \forall n \in \mathbb{N}$
        \end{itemize}
        \item (b) Modify the sequence $(x_n)$ so that it converges to $\sqrt{c}$.
        \begin{itemize}
            \item We can guess what the recurrence relation must be by looking at the equation we had for the limit in the previous slide:
            \[
            x = \frac{1}{2} \left( x + \frac{2}{x} \right)
            \]
            which can be written as
            \[
            x = \frac{2}{x}
            \]
            \item Thus, to get $\sqrt{c}$ instead of $\sqrt{2}$, we replace $2$ with a $c$, giving the recurrence relation
            \[
            x_1 = c, \quad x_{n+1} = \frac{1}{2} \left( x_n + \frac{c}{x_n} \right)
            \]
        \end{itemize}
    \end{itemize}
    \item Exercise 4.9. Write computer code to investigate the convergence of
    \[
    \sqrt{2}, \sqrt{2+\sqrt{2}}, \sqrt{2+\sqrt{2+\sqrt{2}}},\dots
    \]
    \begin{itemize}
        \item (i) Prove either convergence or non-convergence
        \begin{itemize}
            \item We can prove this is monotone increasing by induction. First, note that $2+\sqrt{2} > 2$, and so $x_2 > x_1$
            \item Now suppose that $x_n > x_{n-1}$ for some $n \in \mathbb{N}$. Then
            \[
            x_{n+1} = \sqrt{2 + x_n} > \sqrt{2+ x_{n-1}} = x_n
            \]
            \item Hence, $x_{n+1} > x_n, \ \forall n \in \mathbb{N}$ by manthematical induction.
            \item We now show that $x_n$ is bounded. Note that
            \[
            x_{n+1}^2 = 2 + x_n
            \]
            \item We guess that a large enough bound probably works. In fact, we can show that $x_n \le 2$ through induction.
            \item Firstly, $x_1 = \sqrt{2} < 2$
            \item Now suppose that $x_n \le 2$ for some $n \in \mathbb{N}$. Then,
            \[
            x_{n+1} = \sqrt{2 + x_n} \le \sqrt{2+2} =2 \Rightarrow x_{n+1} \le 2
            \]
            \item So $(\forall n \in \mathbb{N}) \ x_n \le 2$ by induction. We also know trivially that $\sqrt{2}$ is a lower bound since $x_1 = \sqrt{2}$ and the sequence is monotone increasing.
            \item Since the sequence is bounded and monotone increasing, it converges by the Monotone Convergence Theorem.
        \end{itemize}
        \item (ii) If the series converges, find its limit.
        \begin{itemize}
            \item Now that we proved convergence, we can assume $x_n \rightarrow x$ and take the limit of boht sides of our recurrence relation
            \[
            x_{n+1}^2 = 2 + x_n
            \]
            to find, using algebra of limits,
            \[
            x^2 = 2 + x
            \]
            \item This yields $x=2$ as the only admissible solution. Note that while the quadratic gives $x=-1$ as a solution also, but it can't be the limit since $(\forall n \in \mathbb{N}) \ x_n > \sqrt{2}$
        \end{itemize}
    \end{itemize}
    \item Exercise 4.7. Consider the so-called harmonic series
    \[
    \sum_{n=1}^\infty \frac{1}{n}
    \]
    and show that it does not converge.
    \begin{itemize}
        \item The simplest argument is that
        \[
        \sum_{n=1}^N \frac{1}{n} = 1+\frac{1}{2}+\left(\frac{1}{3}+\frac{1}{4}\right)+\left(\frac{1}{5}+\frac{1}{6}+\frac{1}{7}+\frac{1}{8}\right) + \dots + \frac{1}{N} \ge 1 + \frac{1}{2} + \frac{1}{2} + \dots + \frac{1}{2} = 1 + \frac{k}{2}
        \]
        where $k = \lfloor \log_2 (N) \rfloor$
        \item Since $\log_2 (N)$ diverges as $N \rightarrow \infty$, the harmonic series also diverges.
    \end{itemize}
    \item Exercise 4.16. Let $a_n$ be a bounded sequence.
    \begin{itemize}
        \item (a) Prove that the sequence defined by $y_n = \sup \{ a_k : k \ge n \}$ converges.
        \begin{itemize}
            \item Conceptually, $y_n$ is a sequence of the values of the supremum of the subsequence of $a_k$ from $n$ onwards, then $n+1$ onwards, and so on.
            \item Note that $(\forall n \in \mathbb{N}) \ \{ a_k : k \ge n+1 \} \subseteq \{ a_k : k \ge n \}$ and so $\sup \{ a_k : k \ge n+1 \} \le \sup \{ a_k : k \ge n\}$.
            \item Hence $(\forall n \in \mathbb{N}) \ y_{n+1} \le y_n$, implying that $y_n$ is a monotone decreasing sequence.
            \item Furthermore, note that $a_n$ is bounded so $\exists C \in \mathbb{R}$ such that $(\forall n \in \mathbb{N}) \ a_n \ge C$.
            \item Since $y_n$ only consists of elements within $a_n$, $(\forall n \in \mathbb{N}) \ y_n \ge C$ and so $y_n$ is bounded from below.
            \item Since $y_n$ is monotone decreasing and bounded, it converges by the Monotone Convergence Theorem.
        \end{itemize}
        \item (b) The limit superior $(a_n)$, or $\limsup a_n$ is defined by
        \[
        \limsup a_n = \lim y_n
        \]
        where $\lim y_n$ is the sequence from part (a) ofr this exercise. Provide a reasonable definition of lim inf and briefly explain why it always exists for any bounded sequence.
        \begin{itemize}
            \item We can define the sequence $z_n = \inf \{ a_k : k \ge n \}$. 
            \item Then, we can show similarly to the previous part that $z_n$ is monotone increasing and bounded above, which gives us convergence from the Monotone Convergence Theorem.
            \item We can then define
            \[
            \liminf a_n = \lim z_n
            \]
        \end{itemize}
        \item (c) Prove that $\liminf a_n \le \limsup a_n$ for every bounded sequence, and give an example of a sequence for which the inequality is strict.
        \begin{itemize}
            \item Note that $\forall n \in \mathbb{N}$, $z_n = \inf \{ a_k : k \ge n \} \le \sup \{ a_k : k \ge n \} = y_n$
            \item Therefore, by the Order Limit Theorem,
            \[
            \lim z_n \le \lim y_n
            \]
            \item In other words,
            \[
            \liminf a_n \le \limsup a_n
            \]
            for any sequence.
            \item For example, consider $a_n = (-1)^n$. Then $\liminf a_n = -1 < 1 = \limsup a_n$
        \end{itemize}
        \item (d) Show that $\liminf a_n = \limsup a_n$ if and only if $\lim a_n$ exists. In this case, all three share the same value.
        \begin{itemize}
            \item First, let us prove the forwards direction. Suppose that $\liminf a_n = \limsup a_n = a$.
            \item Then, note that
            \[
            z_n = \inf \{ a_k : k \ge n \} \le a_n \le \sup \{ a_k : k \ge n \} = y_n
            \]
            \item Since $z_n \rightarrow a$ and $y_n \rightarrow a$, we conclude using the squeeze theorem that $a_n \rightarrow a$ as well.
            \item Now let us prove the backwards direction. Assume that $a_n \rightarrow a$.
            \[
            (\forall \epsilon > 0)(\exists N \in \mathbb{N})(\forall n \ge N) \ \vert a_n - a \vert < \frac{\epsilon}{2}
            \]
            \item Letting $y_n = \sup \{ a_k : k \ge n \}$ as usual, I claim that $(\forall n \ge N) \ \vert y_n - a \vert \le \frac{\epsilon}{2}$
            \begin{itemize}
                \item Suppose for contradiction that my claim is false. That is, suppose $(\exists n \ge N) \ \vert y_n - a \vert > \frac{\epsilon}{2}$.
                \item This implies that $y_n > a + \frac{\epsilon}{2}$. Note that the other possibility $y_n < a - \frac{\epsilon}{2}$ is not possible since $(\forall n \in \mathbb{N}) y_n \ge a_n$.
                \item But since $y_n$ is the supremum, $\exists k \ge n$ such that $a_k \ge a + \frac{\epsilon}{2}$ (otherwise $a + \frac{\epsilon}{2}$ would be a smaller upper bound).
                \item But this implies that $\vert a_k - a \vert \ge \frac{\epsilon}{2}$ for some $k \ge n \ge N$, which contradicts how $N$ was previously defined.
                \item Thus, the original claim that $(\forall n \ge N) \ \vert y_n - a \vert \le \frac{\epsilon}{2}$ must be true.
            \end{itemize}
            \item Since
            \[
            (\forall \epsilon > 0)(\exists N \in \mathbb{N})(\forall n \ge N) \ \vert y_n - a \vert \le \frac{\epsilon}{2} < \epsilon
            \]
            we can now conclude that $y_n \rightarrow a$.
            \item By using very similar arguments, we can also show that $\liminf a_n = a$.
        \end{itemize}
    \end{itemize}
\end{itemize}

\section{Lecture 8 (11/11/2025)}
\begin{itemize}
    \item Linear System of Equations
    \begin{itemize}
        \item $\left\{ \begin{array}{c}
            x_1 - x_2 = 1 \\
            x_1 + x_2 = 1
        \end{array} \right.$
        \item Represented in matrix form:
        \[
        \left[ \begin{array}{cc}
            1 & 1 \\
            1 & -1
        \end{array} \right] \cdot \left[ \begin{array}{c}
            x_1 \\
            x_2 
        \end{array} \right] = \left[ \begin{array}{c}
            1 \\
            1 
        \end{array} \right]
        \]
        \item Calculation of vector $x$ using inverse matrix:
        \[
        A^{-1} = \left[ \begin{array}{cc}
            0.5 & 0.5 \\
            0.5 & -0.5
        \end{array} \right], \quad x= A^{-1} b = \left[ \begin{array}{c}
            1 \\
            0 
        \end{array} \right], \quad x_1 =1, \quad x_2 = 0
        \]
    \end{itemize}
    \item Matrix Addition ($m \times n$ matrices):
    \[
    \bm{A} + \bm{B} = \left[ \begin{array}{ccc}
        a_{11}+b_{11} & \cdots & a_{1n} + b_{1n} \\
        \vdots & \ddots & \vdots \\
        a_{m1} + b_{m_1} & \cdots & a_{mn} + b_{mn}
    \end{array}\right]
    \]
    \item Matrix Multiplication ($(m \times n) \times (n \times k)$ matrices):
    \[
    \bm{A}\bm{B} = \left[ \begin{array}{ccc}
        \sum_{i=1}^n a_{1i} b_{i1}  & \cdots & \sum_{i=1}^n a_{1i} b_{ik} \\
        \vdots & \ddots & \vdots \\
        \sum_{i=1}^n a_{mi} b_{i1}  & \cdots & \sum_{i=1}^n a_{mi} b_{ik} \\
    \end{array}\right]
    \]
    \item Matrix Transpose ($m \times n$ matrix):
    \[
    \bm{A^T} = \left[ \begin{array}{cccc}
        a_{11} & a_{21} & \cdots & a_{m1} \\
        a_{12} & a_{22} & \cdots & a_{m2} \\
        \vdots & \vdots & \ddots & \vdots \\
        a_{1n} & a_{2n} & \cdots & a_{mn} \\
    \end{array}\right]
    \]
    \item Eigenvalue/Eigenvector - $\bm{A} \underline{e} = \lambda \underline{e}$ where $\lambda$ is the eigenvalue and $\underline{e}$ is the eigenvector
    \begin{itemize}
        \item Applications: stability analysis, principle component analysis (PCA), quantum mechanics
        \item Example:
        \[ \bm{A} = \left[ \begin{array}{cc}
            2 & 1 \\
            1 & 2
        \end{array} \right], \quad \det(\bm{A} - \lambda \bm{I}) = 0
        \]
    \end{itemize}
    \item Exercises
    \begin{itemize}
        \item 1. Compute the matrix product:
        \begin{flalign*}
            \bm{A} = \left[ \begin{array}{cc}
                1 & 2 \\
                3 & 4
            \end{array} \right], \quad \bm{B} = \left[ \begin{array}{cc}
                0 & 1 \\
                1 & 0
            \end{array} \right], \quad \bm{A}\bm{B} = \left[ \begin{array}{cc}
                1 \cdot 0 + 2 \cdot 1 = 2 & 1 \cdot 1 + 2 \cdot 0 = 1 \\
                3 \cdot 0 + 4 \cdot 1 = 4 & 3 \cdot 1 + 4 \cdot 0 = 3
            \end{array} \right]
        \end{flalign*}
        \item 2. Solve the system:
        \begin{flalign*}
            &\bm{A} \underline{x} = \underline{b}, \quad \bm{A} = \left[ \begin{array}{cc}
                1 & 1 \\
                1 & -1
            \end{array} \right], \quad \underline{b} = \left[ \begin{array}{cc}
                2 \\
                0
            \end{array} \right] \\
            &\bm{A}\bm{A^{-1}} = \bm{I} \\
            &\bm{A}\bm{A^{-1}} = \left[ \begin{array}{cc}
                1 & 1 \\
                1 & -1
            \end{array} \right] \cdot \left[ \begin{array}{cc}
                a & b \\
                c & d
            \end{array} \right] = \left[ \begin{array}{cc}
                a+c & b+d \\
                a-c & b-d
            \end{array} \right] = \left[ \begin{array}{cc}
                1 & 0 \\
                0 & 1
            \end{array} \right] \\
            &a+c = 1, \quad b+d = 0, \quad a-c = 0, \quad b-d = 1 \\
            &2a = 1 \Rightarrow a = 0.5 \\
            &2b = 1 \Rightarrow b = 0.5 \\
            &0.5+c = 1 \Rightarrow c = 0.5 \\
            &0.5+d = 0 \Rightarrow d = -0.5 \\
            &\bm{A^{-1}} = \left[ \begin{array}{cc}
                0.5 & 0.5 \\
                0.5 & -0.5
            \end{array} \right] \\
            &\underline{x} = \bm{A}^{-1} \underline{b} = \left[ \begin{array}{cc}
                0.5 & 0.5 \\
                0.5 & -0.5
            \end{array} \right] \left[ \begin{array}{cc}
                2 \\
                0
            \end{array} \right] = \left[ \begin{array}{cc}
                1 \\
                1
            \end{array} \right]
        \end{flalign*}
        \item 3. Find eigenvalues of $\bm{A} = \left[ \begin{array}{cc}
            6 & 4 \\
            4 & 6
        \end{array} \right]$
        \begin{flalign*}
            &\bm{A} \underline{e} = \lambda \underline{e} \text{ for some } \bm{A} \neq \bm{0} \\
            &\bm{A}\underline{e} - \lambda \underline{e} = \bm{0} \Rightarrow  (\bm{A} -\lambda \bm{I})\underline{e} = \bm{0} \\
            &\text{Suppose } (\lambda \bm{I} - \bm{A}) \text{ is invertible} \\
            &\underline{e} = \bm{I} \underline{e} \\
            &=((\lambda \bm{I} - \bm{A})^{-1} (\lambda \bm{I} - \bm{A})) \underline{e} \\
            &= (\lambda \bm{I} - \bm{A})^{-1} ((\lambda \bm{I} - \bm{A}) \underline{e}) \\
            &= (\lambda \bm{I} - \bm{A})^{-1} \cdot \bm{0} \\
            &= \bm{0} \\
            &\text{This claims that } \bm{A} = \bm{0} \text{, which is a contradiction} \\
            &\text{If a matrix is not invertible, then its determinant is equal to } 0 \\
            &\det(\lambda \bm{I} - \bm{A}) = \det(\bm{A} - \lambda \bm{I}) = 0 \\
            &\det(\left[ \begin{array}{cccc}
                a_{11} & a_{12} & \cdots & a_{1m} \\
                a_{21} & a_{22} & \cdots & a_{2m} \\
                \vdots & \vdots & \ddots & \vdots \\
                a_{n1} & a_{n2} & \cdots & a_{nm} 
            \end{array} \right] - \lambda \left[ \begin{array}{cccc}
                1 & 0 & \cdots & 0 \\
                0 & 1 & \cdots & 0 \\
                \vdots & \vdots & \ddots & \vdots \\
                0 & 0 & \cdots & 1 
            \end{array} \right]) = 0 \\
            &\det(\left[ \begin{array}{cc}
                6 & 4 \\
                4 & 6
            \end{array} \right] - \left[ \begin{array}{cc}
                 \lambda & 0 \\
                 0 & \lambda 
            \end{array} \right]) = \det(\begin{array}{cc}
                6-\lambda & 4 \\
                4 & 6-\lambda
            \end{array}) = (6-\lambda)(6-\lambda) - 4^2 = 0 \Rightarrow \lambda = 2
        \end{flalign*}
        \item Characteristic Polynomial of $\bm{A}$ - $p(\lambda) = \det(\lambda \bm{I} - \bm{A})$
        \item Characteristic Equation of $\bm{A}$ - $\det(\lambda \bm{I} - \bm{A}) = 0$
    \end{itemize}
    \item Matrix - With $M,N \in \mathbb{N}$ a real-valued $(M,N)$ (or $M \times N$) matrix $\bm{A}$ is an $M \cdot N$-tuple of elements $a_{ij}, i \in \{1, \dots, M\}, j \in \{1,\dots, N\}$, which is ordered according to a rectangular scheme consisting of $M$ rows and $N$ columns:
    \[
    \bm{A} = \left[ \begin{array}{cccc}
                a_{11} & a_{12} & \cdots & a_{1m} \\
                a_{21} & a_{22} & \cdots & a_{2m} \\
                \vdots & \vdots & \ddots & \vdots \\
                a_{n1} & a_{n2} & \cdots & a_{nm} 
    \end{array} \right], \quad a_{ij} \in \mathbb{R}
    \]
    \begin{itemize}
        \item By convention $(1,N)$-matrices are called rows and $(M,1)$-matrices are called columns (row/column vectors)
        \item $\mathbb{R}^{M \times N}$ is the set of all real-valued $(M,N)$-matrices
    \end{itemize}
    \item Example matrix multiplication
    \begin{itemize}
        \item For $A = \left[ \begin{array}{ccc}
            1 & 2 & 3 \\
            3 & 2 & 1
        \end{array} \right] \in \mathbb{R}^{2 \times 3}, \quad B = \left[ \begin{array}{cc}
            0 & 2 \\
            1 & -1 \\
            0 & 1
        \end{array} \right] \in \mathbb{R}^{3 \times 2}$, we obtain
        \[
        AB = \left[ \begin{array}{ccc}
            1 & 2 & 3 \\
            3 & 2 & 1
        \end{array} \right] \left[ \begin{array}{cc}
            0 & 2 \\
            1 & -1 \\
            0 & 1
        \end{array} \right] = \left[ \begin{array}{cc}
            2 & 3 \\
            2 & 5
        \end{array} \right] \in \mathbb{R}^{2 \times 2}
        \]
        and
        \[
        BA = \left[ \begin{array}{cc}
            0 & 2 \\
            1 & -1 \\
            0 & 1
        \end{array} \right] \left[ \begin{array}{ccc}
            1 & 2 & 3 \\
            3 & 2 & 1
        \end{array} \right] = \left[ \begin{array}{ccc}
            6 & 4 & 2 \\
            -2 & 0 & 2 \\
            3 & 2 & 1
        \end{array} \right] \in \mathbb{R}^{3 \times 3}
        \]
        From this example, we can see that matrix multiplication is not commutative, i.e., $AB \neq BA$.
    \end{itemize}
    \item Identity Matrix - In $\mathbb{R}^{n \times n}$, we define the identity matrix
    \[
    \bm{I}_n := \left[ \begin{array}{cccc}
                1 & 0 & \cdots & 0 \\
                0 & 1 & \cdots & 0 \\
                \vdots & \vdots & \ddots & \vdots \\
                0 & 0 & \cdots & 1 
            \end{array} \right] \in \mathbb{R}^{n \times n}
    \]
    \item System of linear equations with $M$ equations and $N$ unknowns
    \begin{flalign*}
        A_{11}x_1 + \dots + A_{n1}x_n &= b_1 \\
        &\hspace{6mm} \vdots \\
        A_{n1}x_1 + \dots + A_{nm}x_n &= b_m
    \end{flalign*}
    which we can write in vector matrix form as
    \begin{flalign*}
        &\bm{A} \vec{x} = \vec{b} \\
        &\bm{A} = \left[ \begin{array}{cccc}
                a_{11} & a_{12} & \cdots & a_{1n} \\
                a_{21} & a_{22} & \cdots & a_{2n} \\
                \vdots & \vdots & \ddots & \vdots \\
                a_{m1} & a_{m2} & \cdots & a_{mn} 
        \end{array} \right] \\
        &\vec{x} = \left[ \begin{array}{ccc}
            x_1 & \cdots & x_n \\
        \end{array} \right]^\top \\
        &\vec{b} = \left[ \begin{array}{ccc}
            b_1 & \cdots & b_m \\
        \end{array} \right]^\top
    \end{flalign*}
    \item This is rewritten as
    \[
    \sum_{n=1}^N x_n A_{\cdot, n} = \underline{b}
    \]
    \begin{itemize}
        \item We can see that this is a linear combination of the $N$ row vectors $A_{\cdot, 1}, \dots, A_{\cdot, N}$
    \end{itemize}
    \item There are no solutions, one solution, or a set of solutions
    \item No solution system of linear equation
    \begin{flalign*}
        &\left[ \begin{array}{cc}
            1 & 0 \\
            0 & 1 \\
            0 & 0
        \end{array} \right] \left[ \begin{array}{cc}
            x_1 \\
            x_2
        \end{array} \right] = \left[ \begin{array}{cc}
            0 \\
            0 \\
            1
        \end{array} \right] \\
        \Rightarrow &x_1 \left[ \begin{array}{c}
            1 \\
            0 \\
            0
        \end{array} \right] + x_2 \left[ \begin{array}{c}
            0 \\
            1 \\
            0
        \end{array} \right] = \left[ \begin{array}{cc}
            0 \\
            0 \\
            1
        \end{array} \right] \\ \\
        &x_1 = 0, \quad x_2 = 0, \quad 0 = 1
    \end{flalign*}
    \begin{itemize}
        \item It is obvious there is no solution since $0 \neq 1$.
    \end{itemize}
    \item Scalar Product - scalar product of the column vectors $\vec{a} = \left[ \begin{array}{c}
        a_1 \\
        \vdots \\
        a_n
    \end{array} \right], \quad \vec{b} = \left[ \begin{array}{c}
        b_1 \\
        \vdots \\
        b_n
    \end{array}\right]$ is given by
    \[ \vec{a} \cdot \vec{b} = \sum_{i=1}^n a_ib_i \]
    We say $\vec{a}$ and $\vec{b}$ are orthogonal if $\vec{a} \cdot \vec{b} = \vec{0}$
    \item Linear Independence - set of vectors $\{\vec{a_1}, \dots, \vec{a_k}\} \subset \mathbb{R}^n$ is linearly independent if
    \[
    \sum_{i=1}^k x_i\vec{a}_i = 0 \Longleftrightarrow (\forall n \in \{ 1, \dots, k\} ) \ x_n = 0
    \]
    \item Consider the following system of linear equations:
    \begin{flalign*}
        &\left[ \begin{array}{ccc}
            1 & 1 & 1 \\
            0 & 1 & 1  
        \end{array} \right] \left[ \begin{array}{c}
            x_1 \\
            x_2 \\
            x_3
        \end{array} \right] = \left[ \begin{array}{c}
            3 \\
            0
        \end{array} \right] \\
        \Longleftrightarrow &x_1 + x_2 + x_3 = 3 \\
        &\hspace{8.5mm} x_2 + x_3 = 0 \\
        &\vec{x} = \left[ \begin{array}{c}
            x_1 \\
            x_2 \\
            x_3
        \end{array} \right] = \left[ \begin{array}{c}
            3 \\
            -\lambda \\
            \lambda
        \end{array} \right] = \left[ \begin{array}{c}
            3 \\
            0 \\
            0
        \end{array} \right] + \lambda \left[ \begin{array}{c}
            0 \\
            -1 \\
            1
        \end{array} \right] \\
        &\vec{x} \in \left\{ \left[ \begin{array}{c}
            3 \\
            0 \\
            0
        \end{array} \right] + \lambda \left[ \begin{array}{c}
            0 \\
            -1 \\
            1
        \end{array} \right] : \lambda \in \mathbb{R} \right\}
    \end{flalign*}
    \begin{itemize}
        \item The solution set describes the general solution and any one solution is called a particular solution.
        \item Note that $\left[ \begin{array}{c}
            3 \\
            0 \\
            0
        \end{array} \right] \cdot \left[ \begin{array}{c}
            0 \\
            -1 \\
            0
        \end{array} \right] = 0$ and
        \[
        \left[ \begin{array}{ccc}
            1 & 1 & 1 \\
            0 & 1 & 1  
        \end{array} \right] \left[ \begin{array}{c}
            0 \\
            -1 \\
            0
        \end{array} \right] = \left[ \begin{array}{c}
            0 \\
            0
        \end{array} \right]
        \]
        \item The three vectors formed from the columns of $\left[ \begin{array}{ccc}
            1 & 1 & 1 \\
            0 & 1 & 1  
        \end{array} \right]$ are a trivial example of a non-linearly independent set.
        \item The solution set describes the general solution and any one solution is called a particular solution.
    \end{itemize}
    \item When $M < N$, we will always be able to write the general solution of $\underset{M \times N}{\underbrace{\bm{A}}} \underset{N \times 1}{\underbrace{\vec{x}}} = \underset{M \times 1}{\underbrace{\vec{b}}}$ in the form $\vec{a} + \sum_{i=1}^{N-M} \lambda_i \vec{c}_i$, where
    \begin{itemize}
        \item $\vec{a}$ is a particular non-unique solution such that $\bm{A} \vec{a} = \vec{b}$
        \item $\{\vec{c_1}, \dots, \vec{c}_{N-M} \}$ is a set of vectors such that $(\forall j \in \{ 1, \dots, N-M \}) \ \vec{a} \cdot \vec{c}_j = 0 \wedge \bm{A} \vec{c}_j = \vec{0}$
    \end{itemize}
    \item We shall see that
    \begin{itemize}
        \item $\vec{a}$ lies in the row space of $\bm{A}$
        \item $\{ \vec{c}_1, \dots, \vec{c}_{N-M} \}$ is a basis of the kernel of $\bm{A}$, which is the orthogonal complement of the row space of $\bm{A}$
        \item $\vec{b}$ must lie in the column space of $\bm{A}$ for a solution to exist
    \end{itemize}
    \item We shall learn what is meant by
    \begin{itemize}
        \item a vector space
        \item a subspace of a vector space
        \item a basis of a vector space
        \item the row space of a matrix
        \item the column space of a matrix
        \item the kernel or null space of a matrix
        \item the orthogonal complement of subspace within a vector space
    \end{itemize}
    \item Matrix Properties:
    \begin{itemize}
        \item Associativity:
        \[
        \forall \bm{A} \in \mathbb{R}^{m \times n}, \bm{B} \in \mathbb{R}^{n \times p}, C \in \mathbb{R}^{p \times q}: (\bm{A} \bm{B}) \bm{C} = \bm{A} (\bm{B} \bm{C})
        \]
        \item Distributivity:
        \[
        \forall \bm{A}, \bm{B} \in \mathbb{R}^{m \times n}, \bm{C}, \bm{D} \in \mathbb{R}^{n \times p}: (\bm{A} + \bm{B}) \bm{C} = \bm{A} \bm{C} + \bm{B} \bm{C}, \quad \bm{A} (\bm{C} + \bm{D}) = \bm{A} \bm{C} + \bm{A} \bm{D}
        \]
        \item Multiplication with the indetity matrix:
        \[
        \forall \bm{A} \in \mathbb{R}^{m \times n}: \bm{I}_m \bm{A} = \bm
        {A} \times \bm{I} = \bm{A}
        \]
        Note that $\bm{I}_m \neq \bm{I}_n$ for $m \neq n$
    \end{itemize}
    \item Inverse Matrices - Consider a square matrix $\bm{A} \in \mathbb{R}^{n \times n}$. Let matrix $\bm{B} \in \mathbb{R}^{n \times n}$ have the property that $\bm{A} \bm{B} = \bm{I}_n = \bm{B} \bm{A}$. $\bm{B}$ is called the inverse of $\bm{A}$ and denoted by $\bm{A}^{-1}$
    \begin{itemize}
        \item Not every matrix $\bm{A}$ possesses an inverse $\bm{A}^{-1}$
        \item If the inverse does exist, $\bm{A}$ is called regular/invertible/nonsingular
        \item If the inverse does not exist, $\bm{A}$ is called singular/noninvertible
        \item When the matrix inverse exists, it is unique:
        \[
        (\exists \bm{A}^{-1}) \bm{A}^{-1} \bm{A} = \bm{I} = \bm{A} \bm{A}^{-1} \Rightarrow \bm{A}^{-1} \text{ is unique}
        \]
    \end{itemize}
    \item Cofactor Matrix - a matrix where each element is the cofactor of the corresponding element in the original matrix
    \begin{itemize}
        \item Cofactor is calculated by finding the minor (determinant of the submatrix formed by removing the row and column of the element) and multiplying it by $(-1)^{i+j}$, where $i$ and $j$ are the row and column numbers
    \end{itemize}
    \item Adjugate Matrix - transpose of a square matrix's cofactor matrix
    \begin{itemize}
        \item The adjugate of the $2 \times 2$ matrix
        \begin{flalign*}
            A &= \left[ \begin{array}{cc}
                a & b \\
                c & d
            \end{array} \right] \\
            \text{adj} (A) &= \left[ \begin{array}{cc}
                d & -b \\
                -c & a
            \end{array} \right]
        \end{flalign*}
    \end{itemize}
    \item The determinant of a $2 \times 2$ matrix is
    \[
    \det(A) = ad - bc
    \]
    \item The $2 \times 2$ inverse matrix is given by
    \[
    A^{-1} = \frac{1}{\det(A)} \text{adj} (A)
    \]
    \item For a $3 \times 3$ matrix of $A = \left[ \begin{array}{ccc}
        a_{11} & a_{12} & a_{13} \\
        a_{21} & a_{22} & a_{23} \\
        a_{31} & a_{32} & a_{33}
    \end{array}\right]$, its cofactor matrix is
    \begin{flalign*}
        &C = \left[ \begin{array}{ccc}
            (-1)^{1+1} \left\vert \begin{array}{cc}
                a_{22} & a_{23} \\
                a_{32} & a_{33}
            \end{array}\right\vert & (-1)^{1+2} \left\vert \begin{array}{cc}
                a_{21} & a_{23} \\
                a_{31} & a_{33}
            \end{array}\right\vert & (-1)^{1+3} \left\vert \begin{array}{cc}
                a_{21} & a_{22} \\
                a_{31} & a_{32}
            \end{array}\right\vert \\
            (-1)^{2+1} \left\vert \begin{array}{cc}
                a_{12} & a_{13} \\
                a_{32} & a_{33}
            \end{array}\right\vert & (-1)^{2+2} \left\vert \begin{array}{cc}
                a_{11} & a_{13} \\
                a_{31} & a_{33}
            \end{array}\right\vert & (-1)^{2+3} \left\vert \begin{array}{cc}
                a_{11} & a_{12} \\
                a_{31} & a_{32}
            \end{array}\right\vert \\
            (-1)^{3+1} \left\vert \begin{array}{cc}
                a_{12} & a_{13} \\
                a_{22} & a_{33}
            \end{array}\right\vert & (-1)^{3+2} \left\vert \begin{array}{cc}
                a_{11} & a_{13} \\
                a_{21} & a_{23}
            \end{array}\right\vert & (-1)^{3+3} \left\vert \begin{array}{cc}
                a_{11} & a_{12} \\
                a_{21} & a_{22}
            \end{array}\right\vert \\
        \end{array} \right] \\
        &\left\vert \begin{array}{cc}
            a_{im} & a_{in} \\
            a_{jm} & a_{jn}
        \end{array} \right\vert = \det \left[ \begin{array}{cc}
            a_{im} & a_{in} \\
            a_{jm} & a_{jn}
        \end{array} \right] \\
        &\text{adj} (A) = C^T = \left[ \begin{array}{ccc}
            (-1)^{1+1} \left\vert \begin{array}{cc}
                a_{22} & a_{23} \\
                a_{32} & a_{33}
            \end{array}\right\vert & (-1)^{2+1} \left\vert \begin{array}{cc}
                a_{12} & a_{13} \\
                a_{32} & a_{33}
            \end{array}\right\vert & (-1)^{3+1} \left\vert \begin{array}{cc}
                a_{12} & a_{13} \\
                a_{22} & a_{33}
            \end{array}\right\vert \\
            (-1)^{1+2} \left\vert \begin{array}{cc}
                a_{21} & a_{23} \\
                a_{31} & a_{33}
            \end{array}\right\vert & (-1)^{2+2} \left\vert \begin{array}{cc}
                a_{11} & a_{13} \\
                a_{31} & a_{33}
            \end{array}\right\vert & (-1)^{3+2} \left\vert \begin{array}{cc}
                a_{11} & a_{13} \\
                a_{21} & a_{23}
            \end{array}\right\vert \\
            (-1)^{1+3} \left\vert \begin{array}{cc}
                a_{21} & a_{22} \\
                a_{31} & a_{32}
            \end{array}\right\vert  & (-1)^{2+3} \left\vert \begin{array}{cc}
                a_{11} & a_{12} \\
                a_{31} & a_{32}
            \end{array}\right\vert & (-1)^{3+3} \left\vert \begin{array}{cc}
                a_{11} & a_{12} \\
                a_{21} & a_{22}
            \end{array}\right\vert \\
        \end{array} \right] \\
    \end{flalign*}
    \item For a $3 \times 3$ matrix $A$, its determinant is
    \[
    \vert A \vert = \left\vert \begin{array}{ccc}
        a_{11} & a_{12} & a_{13} \\
        a_{21} & a_{22} & a_{23} \\
        a_{31} & a_{32} & a_{33}
    \end{array}\right\vert = a_{11} \left\vert \begin{array}{cc}
        a_{22} & a_{23} \\
        a_{32} & a_{33}
    \end{array} \right\vert - a_{12} \left\vert \begin{array}{cc}
        a_{21} & a_{23} \\
        a_{31} & a_{33}
    \end{array} \right\vert + a_{13} \left\vert \begin{array}{cc}
        a_{21} & a_{22} \\
        a_{31} & a_{32}
    \end{array} \right\vert
    \]
    \item The inverse matrix is given by
    \[
    A^{-1} = \frac{1}{\det(A)} \text{adj} (A)
    \]
    \item Transpose - For $\bm{A} \in \mathbb{R}^{m \times n}$, the matrix $\bm{B} \in \mathbb{R}^{n \times m}$ with $b_{ij} = a_{ji}$ is called the transpose of $\bm{A}$. We write $\bm{B} = \bm{A}^T$).
    \item In general, $\bm{A}^T$ can be obtained by writing the columns of $\bm{A}$ as rows of $\bm{A}^T$
    \item Properties of inverses and transposes:
    \begin{itemize}
        \item $\bm{A}\bm{A}^{-1} = \bm{I} = \bm{A}^{-1} \bm{A}$
        \item $(\bm{A} \bm{B})^{-1} = \bm{B}^{-1} \bm{A}^{-1}$
        \item $(\bm{A} + \bm{B})^{-1} \neq \bm{A}^{-1} + \bm{B}^{-1}$
        \item $\left( \bm{A}^T \right)^T = \bm{A}$
        \item $( \bm{A} + \bm{B})^T = \bm{A}^T + \bm{B}^T$
        \item $(\bm{A}\bm{B})^T = \bm{B}^T \bm{A}^T$
    \end{itemize}
    \item Proof that $(\bm{A} \bm{B})^T = \bm{B}^T \bm{A}^T$
    \begin{itemize}
        \item Let $\bm{A}$ be an $m \times n$ matrix and $\bm{B}$ a $n \times p$ matrix:
        \begin{flalign*}
            [(\bm{A}\bm{B})^T]_{ij} &= [\bm{A}\bm{B}]_{ji} \\
            &= \sum_{a=1}^k \bm{A}_{ja} \bm{B}_{ai} \\
            &= \sum_{a=1}^k [\bm{A}^T]_{aj} [\bm{B}^T]_{ia} \\
            &= \sum_{a=1}^k [\bm{B}^T]_{ia} [\bm{A}^T]_{aj} \\
            &= [\bm{B}^T \bm{A}^T]_{ij}
        \end{flalign*}
    \end{itemize}
    \item Symmetric Matrix - Matrix $\bm{A} \in \mathbb{R}^{n \times n}$ is symmetric if $\bm{A} = \bm{A}^T$
    \begin{itemize}
        \item Only $(N,N)$-matrices (square matrices) can be symmetric
        \item $(\bm{A}^{-1})^T = (\bm{A}^T)^{-1} = A^{-T}$
        \item Sum of symmetric matrices $\bm{A}, \bm{B} \in \mathbb{R}^{n \times n}$ is always symmetric
        \item The product is always defined but generally not symmetric
    \end{itemize}
    \item For $\lambda, \psi \in \mathbb{R}$, the following holds:
    \begin{itemize}
        \item Associativity:
        \[
        (\lambda \psi) \bm{C} = \lambda (\psi \bm{C}), \quad \bm{C} \in \mathbb{R}^{m \times n}
        \]
        \item $\lambda (\bm{B} \bm{C}) = (\lambda \bm{B}) \bm{C} = \bm{B} (\lambda \bm{C}) = (\bm{B} \bm{C}) \lambda, \quad \bm{B} \in \mathbb{R}^{m \times n}, \bm{C} = \mathbb{R}^{n \times p}$
        \item $(\lambda \bm{C})^T=\bm{C}^T \lambda^T = \bm{C}^T \lambda = \lambda \bm{C}^T, \quad (\forall \lambda \in \mathbb{R}) \lambda = \lambda^T$
        \item Distributivity:
        \begin{flalign*}
            &(\lambda + \psi) \bm{C} = \lambda \bm{C} + \psi \bm{C}, \quad \bm{C} \in \mathbb{R}^{m \times n} \\
            &\lambda (\bm{B} + \bm{C}) = \lambda \bm{B} + \lambda \bm{C}, \quad \bm{B}, \bm{C} \in \mathbb{R}^{m \times n}
        \end{flalign*}
    \end{itemize}
    \item Example: Distributivity
    \begin{flalign*}
        C &:= \left[ \begin{array}{cc}
            1 & 2 \\
            3 & 4
        \end{array} \right] \\
        (\lambda + \psi) C &= \left[ \begin{array}{cc}
            (\lambda + \psi) 1 & (\lambda + \psi)2 \\
            (\lambda + \psi)3 & (\lambda + \psi)4
        \end{array} \right] = \left[ \begin{array}{cc}
            \lambda + \psi & 2\lambda + 2\psi \\
            3\lambda + 3\psi & 4\lambda + 4\psi
        \end{array} \right] \\
        &= \left[ \begin{array}{cc}
            \lambda & 2\lambda \\
            3\lambda & 4\lambda
        \end{array} \right] + \left[ \begin{array}{cc}
            \psi & 2\psi \\
            3\psi & 4\psi
        \end{array} \right] = \lambda C +\psi C
    \end{flalign*}
\end{itemize}

\section{Tutorial 6 (13/11/2025)}
\begin{itemize}
    \item Exercise 4.19.
    \begin{itemize}
        \item (a) Show that if $(x_n)$ is a convergent sequence, then the sequence given by the averages
        \[
        y_n = \frac{x_1+x_2+\cdots+x_n}{n}
        \]
        also converges to the same limit.
        \begin{itemize}
            \item Start with rough work. Suppose $x_n \rightarrow x$. Then, consider
            \begin{flalign*}
                \vert y_n - x \vert &= \left\vert \frac{x_1+x_2+\cdots+x_n}{n} - x \right\vert \\
                &= \left\vert \frac{(x_1-x)+(x_2-x)+\cdots+(x_n-x)}{n} \right\vert \\
                &\le \frac{1}{n} \sum_{i=1}^n \vert x_i - x \vert
            \end{flalign*}
            where we used the triangle inequality at the end.
            \item We know that eventually, $\vert x_i - x \vert$ can be made arbitrarily small.
            \item Suppose that $(\forall i \ge N) \ \vert x_i - x \vert < \delta$. How small a $\delta > 0$ would be good enough?
            \item To try and answer, this, $(\forall n \ge N)$, we split the sum as
            \[
            \vert y_n - x \vert \le \frac{1}{n} \sum_{i=1}^n \vert x_i - x \vert = \frac{1}{n} \sum_{i=1}^{N-1} \vert x_i - x \vert + \frac{1}{n} \sum_{i=N}^n \vert x_i - x \vert
            \]
            \item With a fixed $N$, the first sum is just a fixed finite value, say $S_N$. Then we can write
            \[
            \vert y_n - x \vert \le \frac{S_N}{n} + \frac{1}{n} \sum_{i=N}^n \vert x_i - x \vert
            \]
            \item Since we want the whole thing to be less than $\epsilon > 0$, it is sufficient to have each of the two terms less than $\frac{\epsilon}{2}$.
            \item Since $S_N$ is a constant, we know $\frac{1}{n} \rightarrow 0$ so for a large enough $n$, $\frac{S_N}{n} < \frac{\epsilon}{2}$.
            \item For the second term to be less than $\frac{\epsilon}{2}$, each $\vert x_i - x \vert < \delta$. Then note that since $n \ge N \ge 1$,
            \[
            \frac{1}{n} \sum_{i=N}^n \vert x_i - x \vert < \frac{(n-N+1)\delta}{n} \le \delta
            \]
            \item So for this term to be less than $\frac{\epsilon}{2}$, $\delta = \frac{\epsilon}{2}$ is good enough.
            \item Proof:
            \begin{itemize}
                \item Suppose that $x_n \rightarrow x$.
                \[
                (\forall \epsilon > 0)(\exists N_1 \in \mathbb{N})(\forall i \ge N_1) \ \vert x_i - x \vert < \frac{\epsilon}{2}
                \]
                \item Now define $S_{N_1} = \sum_{i=1}^{N-1} \vert x_i - x \vert$. Then, by the Archimedean property, $\exists N_2 \in \mathbb{N}$ such that $(\forall n \ge N_2) \ \frac{S_{N_1}}{n} < \frac{\epsilon}{2}$ by simply taking $N_2 > \frac{2 S{N_1}}{\epsilon}$.
                \item Now define $N = \max(N_1,N_2)$. Then
                \begin{flalign*}
                    \vert y_n - x \vert &\le \frac{S_{N_1}}{n} + \frac{1}{n} \sum_{i=N_1}^n \vert x_i - x \vert \qquad \because \text{Triangle Inequality} \\
                    &< \frac{\epsilon}{2} + \left( \frac{n-N_1+1}{n} \right) \frac{\epsilon}{2} \le \frac{\epsilon}{2} + \frac{\epsilon}{2} = \epsilon
                \end{flalign*}
                \item Summarizing, we have found that
                \[
                (\forall \epsilon > 0)(\exists N \in \mathbb{N})(\forall i \ge N) \ \vert y_i - x \vert < \epsilon
                \]
                and so $y_n \rightarrow x$
            \end{itemize}
        \end{itemize}
        \item (b) Given an example to show that it is possible for the sequence $(y_n)$ of averages to converge even if $(x_n)$ does not.
        \begin{itemize}
            \item Consider the sequence $x_n = (-1)^n$, which does not converge.
            \item However, we observe that the sequence $y_n$ is
            \[
            y_n = \left\{ \begin{array}{rl}
                -\frac{1}{n} & \text{if } n \text{ is odd} \\
                0 & \text{if } n \text{ is even}
            \end{array} \right.
            \]
            \item Using the Archimedean property of natural numbers, we can set $N > \frac{1}{\epsilon}$ for all $\epsilon > 0$. Then we can say that
            \begin{flalign*}
                (\forall \epsilon > 0)(\exists N \in \mathbb{N}) N &> \frac{1}{\epsilon} \\
                \Rightarrow \frac{1}{N} &< \epsilon \\
                \Rightarrow \vert -\frac{1}{N}  \vert &< \epsilon
            \end{flalign*}
            \item Now since $(\forall n \in \mathbb{N}) \ \vert y_n \vert \le \vert -\frac{1}{n} \vert$, we conclude that
            \begin{flalign*}
                (\forall \epsilon > 0)(\exists N \in \mathbb{N})(\forall n \ge N) \vert y_n - 0\vert \le \ \vert -\frac{1}{n} \vert \le \vert -\frac{1}{N} \vert &< \epsilon \\
                \Rightarrow \vert y_n - 0 \vert &< \epsilon
            \end{flalign*}
            \item Therefore, $y_n \rightarrow 0$.
        \end{itemize}
    \end{itemize}
\end{itemize}

\section{Lecture 9 (18/11/2025)}
\begin{itemize}
    \item Goals:
    \begin{itemize}
        \item Rotate vectors around the origin
        \item Eigenvectors and eigenvalues
        \item Linear combination, linear independence, spanning set, span, basis
    \end{itemize}
    \item Start with $\vec{e}_1 = \left[ \begin{array}{c}
        1 \\
        0
    \end{array}\right]$ and $\vec{e}_2 = \left[ \begin{array}{c}
        0 \\
        1
    \end{array}\right]$
    \begin{itemize}
        \item $\left[ \begin{array}{c}
            x_1 \\
            x_2
        \end{array}\right] = x_1\vec{e}_1 + x_2 \vec{e}_2$
        \item Note that the columns of $A$ satisfy $A_{\cdot,1} = A\vec{e}_1, \quad A_{\cdot,2} \vec{e}_2$
        \[
        \left[ \begin{array}{c}
            a_{11} \\
            \vdots \\
            a_{m1}
        \end{array}\right] = \left[ \begin{array}{cc}
            a_{11} & a_{12} \\
            \vdots & \vdots \\
            a_{m1} & a_{m2}
        \end{array}\right] \left[ \begin{array}{c}
            1 \\
            0
        \end{array}\right], \qquad \left[ \begin{array}{cc}
            a_{12} \\ \vdots \\ a_{m2}
        \end{array}\right] = \left[ \begin{array}{cc}
            a_{11} & a_{12} \\
            \vdots & \vdots \\
            a_{m1} & a_{m2}
        \end{array}\right] \left[ \begin{array}{c}
            0 \\
            1
        \end{array}\right]
        \]
        \item Matrix represents a linear map $\alpha$ with respect to the basis $\{\vec{e}_1,\vec{e}_2\}$
        \begin{itemize}
            \item $\alpha(\vec{e}_1) = A_{\cdot,1}, \quad \alpha(\vec{e}_2) = A_{\cdot,2}, \quad A = \left[ \begin{array}{cc}
                \alpha(\vec{e}_1) & \alpha(\vec{e}_2) \\
            \end{array}\right]$
        \end{itemize}
        \item After rotation by $\theta$ radians, $\left[ \begin{array}{c}
            1 \\
            0
        \end{array}\right] \rightarrow \left[ \begin{array}{c}
            \cos(\theta) \\
            \sin(\theta)
        \end{array}\right]$ and $\left[ \begin{array}{c}
            0 \\
            1
        \end{array}\right] \rightarrow \left[ \begin{array}{c}
            -\sin(\theta) \\
            \cos(\theta)
        \end{array}\right]$
    \end{itemize}
    \item General 2D vector rotation: $\left[ \begin{array}{c}
        x \\
        y
    \end{array} \right] \overset{R(\theta)}{\rightarrow} \left[ \begin{array}{c}
        x\cos(\theta) - y \sin(\theta) \\
        x\sin(\theta) + y\cos(\theta)
    \end{array}\right]$
    \[
    R(\theta) = \left[ \begin{array}{cc}
        \cos(\theta) & -\sin(\theta) \\
        \sin(\theta) & \cos(\theta)
    \end{array}\right]
    \]
    \item Consider rotation through $\theta$ radians as a mapping:
    \[
        R(\theta) : \vec{v} \rightarrow R(\theta)\vec{v}, \quad \vec{v} \in \mathbb{R}^2
    \]
    \item For any vector $\vec{v} = v_1 \vec{e}_1 + v_2 \vec{e}_2$, the map $R(\theta)$ satisfies the linearity properties
    \begin{itemize}
        \item Additivity:
        \[
        \forall \vec{v},\vec{w} \in \mathbb{R}^2: \quad R(\theta)(\vec{v} + \vec{w}) = R(\theta) \vec{v} + R(\theta)\vec{w}
        \]
        \item Homogeneity:
        \[
        \forall \vec{v} \in \mathbb{R}^2, \forall \lambda \in \mathbb{R} : \quad R(\theta)(\lambda \vec{v}) = \lambda R(\theta) \vec{v}
        \]
        \item Linear map - any map $X: V \rightarrow U$ that satisfies additivity and homogeneity
    \end{itemize}
    \item Represent this linear map w.r.t. basis $\{\underline{e}_1,\underline{e}_2\}$
    \begin{itemize}
        \item $R_{\cdot,1} = \rho(\theta)\underline{e}_1 = \left[ \begin{array}{c}
            \cos(\theta) \\
            \sin(\theta)
        \end{array} \right]$
        \item $R_{\cdot,2} = \rho(\theta)\underline{e}_2 = \left[ \begin{array}{c}
            -\sin(\theta) \\
            \cos(\theta)
        \end{array} \right]$
        \item $\left[ \begin{array}{c}
            x \\
            y 
        \end{array} \right] = x_1 \underline{e}_1 + x_2 \underline{e}_2, \quad \rho(\theta) \left[ \begin{array}{c}
            x \\
            y 
        \end{array} \right] = x_1 \rho(\theta)\underline{e}_1 + x_2\rho(\theta)\underline{e}_2 = x_1 \left[ \begin{array}{c}
            \cos(\theta) \\
            \sin(\theta)
        \end{array} \right] + x_2 \left[ \begin{array}{c}
            -\sin(\theta) \\
            \cos(\theta)
        \end{array} \right]$
        \item This is equal to $\left[ \begin{array}{c}
            x_1 \cos(\theta) - x_2 \sin(\theta) \\
            x_1 \sin(\theta) + x_2 \cos(\theta)
        \end{array}\right] = R(\theta)\underline{x}$
    \end{itemize}
    \item Another possible basis for $\mathbb{R}^2$ is $\{\frac{1}{\sqrt{2}}(\underline{e}_1 - \underline{e}_2), \frac{1}{\sqrt{2}} (\underline{e}_1 + \underline{e}_2)\}$
    \item Standard basis for $\mathbb{R}^3: \{\underline{e}_1,\underline{e}_2, \underline{e}_3\}$
    \begin{itemize}
        \item Mapping $R_z(\theta) : \underline{v} \rightarrow R_z(\theta) \underline{v}$ for $\underline{v} \in \mathbb{R}^3$ represents rotation around the z-axis
        \[
        R_{z,\theta} = \left[ \begin{array}{ccc}
            \cos(\theta) & -\sin(\theta) & 0 \\
            \sin(\theta) & \cos(\theta) & 0 \\
            0 & 0 & 1
        \end{array} \right]
        \]
        \item Note that because we are rotating around the z-axis, the z-component remains unchanged
    \end{itemize}
    \item $\underline{e}_i \cdot \underline{e}_j = \left\{ \begin{array}{cc}
        0 & \text{if } i \neq j \\
        1 & \text{if } i = j
    \end{array} \right.$
    \begin{itemize}
        \item Orthogonal basis if $\underline{e}_i \cdot \underline{e}_j = 0$
        \item Orthonormal basis if $\underline{e}_i \cdot \underline{e}_j = 0$ and $\vert \vert e_{i} \vert \vert = \vert \vert e_{j} \vert \vert = 1$
    \end{itemize}
    \item Sequential rotations around the same axis combine additively
    \[
    R_{z,\theta_1}R_{z,\theta_2} = R_{z,(\theta_1+\theta_2)} = \left[ \begin{array}{ccc}
        \cos(\theta_1 + \theta_2) & -\sin(\theta_1,\theta_2) & 0 \\
        \sin(\theta_1+\theta_2) & \cos(\theta_1 + \theta_2) & 0 \\
        0 & 0 & 1
    \end{array} \right]
    \]
    \begin{itemize}
        \item Rotating by $\theta_1$ then $\theta_2$ is equivalent to $\theta_1 + \theta_2$
        \item Matrices which represent rotations around same axis are commutative
    \end{itemize}
    \item Inverse mapping: $R_z^{-1}(\theta) : \underline{v} \rightarrow R_z^{-1}(\theta) \underline{v}$
    \begin{itemize}
        \item Geometrically: $R_z^{-1}(\theta)=R_z(-\theta)$
        \item Inverse matrix:
        \[
        R_{z,\theta}^{-1} = \left[ \begin{array}{ccc}
            \cos(-\theta) & -\sin(-\theta) & 0 \\
            \sin(-\theta) & \cos(-\theta) & 0 \\
            0 & 0 & 1
        \end{array} \right] = \left[ \begin{array}{ccc}
            \cos(\theta) & \sin(\theta) & 0 \\
            -\sin(\theta) & \cos(\theta) & 0 \\
            0 & 0 & 1
        \end{array} \right] = R_{z,\theta}^\top
        \]
        \begin{itemize}
            \item Note that $R_{z,\theta}^{-1} = R_{z,\theta}^T$ (rotation matrices are orthogonal)
        \end{itemize}
    \end{itemize}
    \item For the rotation $R_z(\theta)$ around the z-axis, any vector of the form $\lambda(0,0,1)^T$ is invariant:
    \[
    R_{z,\theta}\lambda\left[ \begin{array}{c}
        0 \\
        0 \\
        1
    \end{array} \right] = \lambda \left[ \begin{array}{c}
        0 \\
        0 \\
        1
    \end{array} \right]
    \]
    \item Clearly $\underline{e}_3$ is an eigenvector to the 3D rotational matrix around the z-axis with eigenvalue $\lambda = 1$
    \begin{itemize}
        \item Define $E_3 = \{\mu\underline{e}_3\ : \mu \in \mathbb{R}\}$
        \item $(\forall \underline{v} \in E_3) \ R_{z,\theta} \underline{v} = \underline{v}$
        \[
        R_{z,\theta} \underline{v} = \left[ \begin{array}{ccc}
                \cos(\theta) & -\sin(\theta) & 0 \\
                \sin(\theta) & \cos(\theta) & 0 \\
                0 & 0 & 1
            \end{array} \right] \left[ \begin{array}{c}
                0 \\ 0 \\ \mu
            \end{array} \right] = \left[ \begin{array}{c}
                0 \\ 0 \\ \mu
            \end{array} \right]
        \]
        \item Every vector in $E_3$ is an eigenvector of $R_z(\theta)$ with eigenvector 1
    \end{itemize}
    \item Eigenvector - nonzero vector that changes at most by a scalar factor under a linear transformation
    \begin{itemize}
        \item Eigenvector of the $n \times n$ matrix $\bm{A} \in \mathbb{R}^{n \times n}$ is a vector $\vec{e} \in \mathbb{C}^n - \{ \bm{0} \}$ with corresponding eigenvalue $\lambda \in \mathbb{C}$ such that
        \[
        \bm{A} \underline{e} = \lambda \underline{e}
        \]
        \item Eigenvectors must be nonzero
        \item Eigenvalues and eigenvectors can be complex
        \item Eigenvectors are unique up to scalar multiplication
        \item Corresponding eigenvalue is the scaling factor
        \item For differential operators like $\frac{d}{dx}$, eigenvectors are functions (eigenfunctions):
        \[
        \frac{d}{dx} e^{\lambda x} = \lambda e^{\lambda x}
        \]
        \[
        Df = \lambda f
        \]
    \end{itemize}
    \item Example: find all eigenvalues and eigenvectors of
    \[
    \bm{A} = \left[ \begin{array}{ccc}
        0 & -1 & 0  \\
        1 & 0 & 0 \\
        0 & 0 & 1
    \end{array} \right]
    \]
    \item Starting from $\bm{A} \underline{e} = \lambda \underline{e}$, we get
    \[
    (\bm{A} - \lambda \bm{I}) \underline{e} = 0
    \]
    \begin{flalign*}
        \det(\bm{A} - \lambda \bm{I}) &= \det \left( \left[ \begin{array}{ccc}
            -\lambda & -1 & 0  \\
            1 & -\lambda & 0 \\
            0 & 0 & 1-\lambda
        \end{array} \right]\right) \\
        &= -\lambda(-\lambda(1-\lambda)-0) - (-1)(1(1-\lambda) - 0) + 0 \\
        &= \lambda^2(1-\lambda) + (1-\lambda) \\
        &= (1-\lambda)(\lambda^2+1) \\
        &= 0
    \end{flalign*}
    \begin{itemize}
        \item Therefore, $\lambda \in \{1,i,-i\}$.
        \item Now to find the eigenvectors, we solve the system of equations to find the $\ker(\bm{A} - \lambda \bm{I})$:
    \end{itemize}
    \begin{flalign*}
        (\bm{A} - \lambda \bm{I}) \underline{e} &= \bm{0} \\
        \left[ \begin{array}{ccc}
            -\lambda & -1 & 0 \\
            1 & -\lambda & 0 \\
            0 & 0 & 1-\lambda
        \end{array} \right] \left[ \begin{array}{c}
            x_1 \\
            x_2 \\
            x_3
        \end{array} \right] &= \left[ \begin{array}{c}
            0 \\
            0 \\
            0
        \end{array} \right] \\
        0 &= \left\{ \begin{array}{c}
            -\lambda x_1 - x_2 \\
            x_1 - \lambda x_2 \\
            (1-\lambda) x_3
        \end{array} \right.
    \end{flalign*}
    \item For $\lambda = 1$,
    \[
    \left\{ \begin{array}{c}
            -x_1 - x_2 = 0 \\
            x_1 - x_2 = 0 \\
            0 = 0
        \end{array} \right.
    \]
    \begin{itemize}
        \item $x_1=x_2=0$ and $x_3 \in \mathbb{R}$ is a free variable:
        \[
        \lambda = \left\{ \alpha \left[ \begin{array}{c}
            0 \\
            0 \\
            1
        \end{array} \right] : \alpha \in \mathbb{R} \backslash \{0\} \right\}
        \]
    \end{itemize}
    \item For $\lambda = i$,
    \[
    \left\{ \begin{array}{c}
            -ix_1 - x_2 = 0 \Rightarrow x_2 = -ix_1 \\
            x_1 - ix_2 = 0 \\
            (1-i)x_3 = 0 \Rightarrow x_3=0
        \end{array} \right.
    \]
    \begin{itemize}
        \item $x_1 \in \mathbb{R}$ is a free variable:
        \[
        \lambda = \left\{ \alpha \left[ \begin{array}{c}
            1 \\
            -i \\
            0
        \end{array} \right] : \alpha \in \mathbb{R} \backslash \{0\} \right\}
        \]
    \end{itemize}
    \item For $\lambda = -i$,
    \[
    \left\{ \begin{array}{c}
            ix_1 - x_2 = 0 \Rightarrow x_2 = ix_1 \\
            x_1 + ix_2 = 0 \\
            (1+i)x_3 = 0 \Rightarrow x_3=0
        \end{array} \right.
    \]
    \begin{itemize}
        \item $x_1 \in \mathbb{R}$ is a free variable:
        \[
        \lambda = \left\{ \alpha \left[ \begin{array}{c}
            1 \\
            i \\
            0
        \end{array} \right] : \alpha \in \mathbb{R} \backslash \{0\} \right\}
        \]
    \end{itemize}
\end{itemize}
\subsection{Slide 5 (Linear Algebra: Vector Spaces)}
\begin{itemize}
    \item Vector is an element of a vector space
    \item Elements of $\mathbb{R}, \mathbb{C}$ are called scalars
    \item Vector Space - vector space over scalar field $\mathbb{F}$ is a set $V$ with two binary operations:
    \begin{itemize}
        \item Vector Addition: $(v,w) \mapsto v + w$
        \item Scalar Multiplication: $(a,v) \mapsto a \cdot v$
        \item Closure Properties:
        \begin{flalign*}
            &\forall v,w \in V : v+w \in V \\
            &\forall a \in \mathbb{F}, \forall v \in V : a \cdot v \in V
        \end{flalign*}
    \end{itemize}
    \item Vector Space Axioms
    \begin{itemize}
        \item Associativity of Addition:
        \[
        (\forall u,v,w \in V) \quad u+(v+w)=(u+v)+w
        \]
        \item Commutativity of Addition:
        \[
        (\forall v,w \in V) \quad v+w = w+v
        \]
        \item Identity Element of Addition
        \[
        (\exists 0 \in V)(\forall v \in V) \quad v+0 = v
        \]
        This is called the zero vector
        \item Inverse Elments of Addition:
        \[
        (\forall v \in V)(\exists -v \in V) \quad v + (-v) = 0
        \]
        $-v$ is the additive inverse of $v$
        \item Distributivity (scalar over vector addition):
        \[
        (\forall a \mathbb{F})(\forall v,w \in V) \quad a \cdot (v+w) = a \cdot v + a \cdot w
        \]
        \item Distributivity (scalar addition):
        \[
        (\forall a,b \in \mathbb{F})(\forall v \in V) \quad (a+b) \cdot v = a \cdot v + b \cdot v
        \]
        \item Scalar Multiplication Associativity:
        \[
        (\forall a,b \in \mathbb{F})(\forall v \in V) \quad a \cdot (bv) = (ab)v
        \]
        \item Multiplicative Identity:
        \[
        (\exists 1 \in \mathbb{F})(\forall v \in V) \quad 1 \cdot v = v
        \]
    \end{itemize}
    \item Linear Combination - given vectors $\{v_1,v_2,\dots, v_n\}$ in a vector space $V$ over $\mathbb{F}$, a linear combination is
    \[
    a_1v_1+a_2v_2+\cdots+a_nv_n
    \]
    where $a_1,a_2,\dots,a_n \in \mathbb{F}$
    \begin{itemize}
        \item Example: $V = \mathbb{R}^3, \quad S=\{\underline{e}_1,\underline{e}_2\}$
        \item $5\underline{e}_1 + 3\underline{e}_2, 0\underline{e}_1 + 0\underline{e}_2, -\underline{e}_1 + 10^{10}\underline{e}_2$, etc.
    \end{itemize}
    \item Linear Independence - vectors $\{x_1,\dots,x_n\}$ are linear independent if
    \[
    \sum_{i=1}^n \lambda_i x_i = 0 \Longleftrightarrow (\forall i \in \{1,\dots,n\}) \lambda_i = 0
    \]
    Interpretation: No vector can be expressed as a linear combination of the others
    \begin{itemize}
        \item $V=\mathbb{R}^3, \quad \mathbb{F}=\mathbb{R}$
        \item $S_1 = \left\{ \left[ \begin{array}{c}
            1 \\
            0 \\
            0
        \end{array} \right], \left[ \begin{array}{c}
            0 \\
            3 \\
            0
        \end{array} \right], \left[ \begin{array}{c}
            1 \\
            -3 \\
            0
        \end{array} \right]  \right\}$
        \begin{flalign*}
            &\lambda_1 \left[ \begin{array}{c}
                1 \\
                0 \\
                0
            \end{array} \right] +  \lambda_2 \left[ \begin{array}{c}
                0 \\
                3 \\
                0
            \end{array} \right] + \lambda_3 \left[ \begin{array}{c}
                1 \\
                -3 \\
                0
            \end{array} \right] = \left[ \begin{array}{c}
                0 \\
                0 \\
                0
            \end{array} \right] \\
            &\lambda_1 + \lambda_3 = 0 \\
            &3\lambda_2 - 3\lambda_3 = 0 \\
            &\lambda_2 = \lambda_3 \\
            &\lambda_1 = -\lambda_2 \\
            &\left[ \begin{array}{c}
                \lambda_1 - \lambda_3 \\
                3\lambda_2 - 3\lambda_3 \\
                0
            \end{array} \right] = \left[ \begin{array}{c}
                0 \\
                0 \\
                0
            \end{array} \right] \\
        \end{flalign*}
        \item Since There is a non-zero solution of $\lambda_i$s, it is not linearly independent. We can see this with an example: $\lambda_1 = 1, \lambda_2 = -1, \lambda_3 = -1$
        \begin{flalign*}
            &\left[ \begin{array}{c}
                1 \\
                -3 \\
                0
            \end{array} \right] = \left[ \begin{array}{c}
                1 \\
                0 \\
                0
            \end{array} \right] - \left[ \begin{array}{c}
                0 \\
                3 \\
                0
            \end{array} \right]
        \end{flalign*}
        \item Not linearly independent
        \item $S_2 = \{\underline{e}_1,\underline{e}_2\}$ are linearly independent.
        \begin{flalign*}
            &\lambda_1 \left[ \begin{array}{c}
                1 \\
                0 \\
                0
            \end{array} \right] +  \lambda_2 \left[ \begin{array}{c}
                0 \\
                1 \\
                0
            \end{array} \right] = \left[ \begin{array}{c}
                \lambda_1 \\
                \lambda_2 \\
                0
            \end{array} \right] = \left[ \begin{array}{c}
                0 \\
                0 \\
                0
            \end{array} \right] \Rightarrow \lambda_1=\lambda_2=0
        \end{flalign*}
    \end{itemize}
\end{itemize}

\section{Lecture 10 (20/11/2025)}
\begin{itemize}
    \item Definition - Let $V$ be a vector space over $\mathbb{F}$. A subset $W \subseteq V$ is a subspace if $W$ is itself a vector space over $\mathbb{F}$ with the same operations as $V$
    \item (Proposition) $W$ is a subspace of $V$ if and only if:
    \begin{itemize}
        \item The zero vector 0 is in $W$
        \item $\forall u,v \in W : u+v \in W$ (closed under addition)
        \item $\forall a \in \mathbb{F}, \forall v \in W : a \cdot v \in W$ (closed under scalar multiplication)
        \item Note: no need to verify all 8 axioms
    \end{itemize}
    \item Definition (Span) - the span of a set $\mathcal{S} = \{v_1,\dots, v_n\}$ is all the vectors that are possible to make from a linear combination of the vectors in $\mathcal{S}$:
    \[
    span(\mathcal{S}) = \left\{ \sum_{i=1}^n \lambda_iv_i : (\forall i) \lambda_i \in \mathbb{F} \right\}
    \]
    \item Definition (Spanning Set) - $\mathcal{S}$ is a spanning set of $V$ if every vector in $V$ can be expressed as a linear combination of vectors in $\mathcal{S}$
    \item Definition (Basis) - a basis for a subspace $U$ is a linearly independent spanning set of $U$
    \item (Lemma) If $\mathcal{S} = \{ v_1,\dots, v_n\}$ is a basis for $V$ and $\mathcal{T} = \{w_1,\dots,w_n\}$ is linearly independent in $V$, then $m \le n$
    \item Example:
    \begin{itemize}
        \item $V=\mathbb{R}^3$ with standard basis $S=\{ e_1,e_2,e_3 \}$
        \item $T=\{ w_1,w_2,w_3,w_4 \}$ is linearly independent set in $V$
        \item $S$ is basis $\Rightarrow (\forall k \in \{1,2,3\})(\exists x_k) w_k = \sum_{i=1}^3 x_ke_k$ (spanning set property)
        \begin{itemize}
            \item $w_1 = x_{11}e_1 + x_{12}e_2 + x_{13}e_3$
            \item $w_2 = x_{21}e_1 + x_{22}e_2 + x_{23}e_3$
            \item $w_3 = x_{31}e_1 + x_{32}e_2 + x_{33}e_3$
            \item $w_4 = x_{41}e_1 + x_{42}e_2 + x_{43}e_3$
        \end{itemize}
        \item Write $e_1,e_2,e_3$ as linear combinations of $w_1,w_2,w_3$ (linear independence property)
        \item $w_4$ can be written as a linear combination of $w_1,w_2,w_3$, which means that $T$ is not a linearly independent set
    \end{itemize}
    \item (Proposition) If a vector space $V$ has a basis with $n$ vectors, then all other bases of $V$ contain exactly $n$ vectors
    \begin{itemize}
        \item It can't be $n+k$ since it will fail the linear independence property
        \item We will prove it won't be $n-l$ since it will fail the spanning set property
    \end{itemize}
    \item Definition (Dimension) - if a vector space $V$ has a basis containing $n \in \mathbb{N}$ vectors, the dimension of $V$ is $n$.
    \begin{itemize}
        \item Example: $\mathbb{R}^3$ has dimension 3 with standard basis $\{e_1,e_2,e_3\}$
    \end{itemize}
    \item Definition (Complement) - a complement to a subspace $X$ of $V$ is another subspace $Y$ such that $X \cap Y = \{0\}$
    \begin{itemize}
        \item $V=\mathbb{R}^2, \quad (X \in \mathbb{R}^2 : y=0 \wedge x \in \mathbb{R}), \quad (Y \in \mathbb{R}^2 : x = 0 \wedge y \in \mathbb{R})$
        \item Basically x-axis and y-axis, which are subspaces because they contain 0 and they remain within the subspace under scalar multiplication and addition
        \item Easy to prove that $X \cap Y = \{0\}$
        \item $X$ is a complement of $Y$ in $V$ and $Y$ is a complement of $X$ in $V$
        \item Also $(Z \in \mathbb{R}^2 : y=x)$ is complement of $X$ and $Y$
    \end{itemize}
    \item Definition (Direct Sum) - let $X$ and $Y$ be subspaces of $V$. If
    \[
    (\forall v \in V)(\exists !x \in X)(\exists !y \in Y) \ v = x+y
    \]
    then $V$ is the direct sum of $X$ and $Y$, written $V = X \oplus Y$
    \begin{itemize}
        \item $\exists!$ means there only exists exactly one unique element
        \item $(X \in \mathbb{R}^2 : y=0 \wedge x \in \mathbb{R}), \quad (Y \in \mathbb{R}^2 : x = 0 \wedge y \in \mathbb{R}), \quad X \oplus Y = \mathbb{R}^2$
        \item Interpretation: every vector in $V$ can be uniquely decomposed into components from $X$ and $Y$
    \end{itemize}
    
    \item Definition (Image) - let $V$ and $W$ be vector spaces over $\mathbb{F}$ and the map. $f : V \rightarrow W$ be defined. The image of $v$ under $f$ is the set
    \[
    \text{im}(f) = \{ f(v) : v \in V \}
    \]
    \begin{itemize}
        \item If $Y \subseteq W$, then the pre-image of $Y$ is
        \[
        f^{-1}(Y) = \{ v \in V \mid f(v) \in Y \}
        \]
    \end{itemize}
    \item Definition (Linear Map) - a function $f : V \rightarrow W$ is a linear map if
    \begin{gather*}
        (\forall u,v \in V)(\forall \lambda \in \mathbb{F}) \\
        f(u+v)=f(u)+f(v) \text{ and } f(\lambda v) = \lambda f(v)
    \end{gather*}
    \begin{itemize}
        \item Nonlinear map example:
        \begin{itemize}
            \item $f : \mathbb{R}^2 \rightarrow \mathbb{R}^2$
            \item $f(x,y)^T = (x^2+y^2, x^2-y^2)^T$
            \item $f(\lambda(x,y)^T) = f((\lambda x, \lambda y)^T) = (\lambda^2x^2+\lambda^2y^2,\lambda^2x^2-\lambda^2y^2)^T = \lambda^2 f(x,y)^T$
        \end{itemize}
    \end{itemize}
    \item Definition (Injective) - $f:V \rightarrow W$ is injective (one-to-one) if:
    \[
    (\forall v_1,v_2 \in V) \ v_1 \neq v_2 \Rightarrow f(v_1) \neq f(v_1)
    \]
    \begin{itemize}
        \item Distinct elements of $V$ have distinct images
        \item Consider $f : \mathbb{R} \rightarrow \mathbb{R}, \quad f : x \rightarrow x^2$
        \begin{itemize}
            \item $f(-1)=f(1)$ so not injective
        \end{itemize}
    \end{itemize}
    \item Definition (Surjective) - $f : V \rightarrow W$ is surjective (onto) if:
    \[
    (\forall w \in W)(\exists v \in V) \ f(v) = w
    \]
    \begin{itemize}
        \item Consider $f : \mathbb{R} \rightarrow \mathbb{R}, \quad f : x \rightarrow x^2$
        \begin{itemize}
            \item Suppose $w=-4$, then you can't get to $w$ through the mapping
            \item However, $f : \mathbb{R} \rightarrow \mathbb{R}^+$ is surjective
        \end{itemize}
    \end{itemize}
    \item Definition (Bijective) - $f : V \rightarrow W$ is bijective if it is both injective and surjective
    \item Definition (Range) - the range of $f:V\rightarrow W$ is:
    \[
    ran(f) = f(V) = \{ f(v) : v \in V \} \subset W
    \]
    \item Definition (Kernel) - the kernel (null space) of a linear map $L : V \rightarrow W$ is:
    \[
    ker(L) = \{ v \in V : L(v) = 0 \}
    \]
    \item Key result: a linear map $L$ is injective if and only if $ker(L) = \{0\}$
    \begin{itemize}
        \item If $ker(L)$ includes multiple elements, then multiple elements $L(v)=0$, which means it isn't injective
    \end{itemize}
    \item Example 1
    \begin{itemize}
        \item Consider $L : \mathbb{R}^3 \rightarrow \mathbb{R}^3$ defined by
        \[
        L(x,y,z)^T = (x+y, y+z,0)^T
        \]
        \item Finding the kernel:
        \begin{flalign*}
            &x+y = 0 \\
            &y+z = 0
        \end{flalign*}
        \item Solution: $y = -z$ and $x=z$
        \item Therefore: $ker(L) = \{ \lambda(1,-1,1)^T : \lambda \in \mathbb{R} \}$
        \item Finding the range:
        \begin{flalign*}
            L \left[ \begin{array}{cc}
                x \\ y \\ z
            \end{array} \right] &= \left[ \begin{array}{cc}
                x+y \\ y+z \\ 0
            \end{array} \right] = (x+y) \left[ \begin{array}{cc}
                1 \\ 0 \\ 0
            \end{array} \right] + (y+z) \left[ \begin{array}{cc}
                0 \\ 1 \\ 0
            \end{array} \right] \\
            ran(L) &= span\{(1,0,0)^T, (0,1,0)^T\}
        \end{flalign*}
    \end{itemize}
    \item Example 2
    \begin{flalign*}
        \frac{1}{2}\sigma^2 y'' + \mu y' - ry &= 0 \\
        \frac{1}{2} \sigma^2 \frac{d^2 y}{dx^2} + \mu \frac{dy}{dx} -ry &= 0 \\
        \underset{\text{Linear operator}}{\underbrace{\left( \frac{1}{2} \sigma^2 \frac{d^2}{dx^2} + \mu \frac{d}{dx} - r \right)}} y &= 0 \\
        \text{Define } L &= \frac{1}{2}\sigma^2 \frac{d^2}{dx^2} + \mu \frac{d}{dx} - r \\
        Ly &= 0
    \end{flalign*}
\end{itemize}


\section{Tutorial 7 (20/11/2025)}
\begin{itemize}
    \item Vectors are anything that when added together or multiplied by scalar (from some field), you get another vector
    \item Matrices are objects that map vectors to other vectors linearly
    \begin{itemize}
        \item We write them as a 2D rectangular grid of numbers with $M$ rows and $N$ columns, giving an $M \times N$ matrix
        \item $(M \times N) \times (N \times 1) \rightarrow M \times 1$: $M$-dimensional vector produced
        \item Matrices are very useful in representing and solving systems of linear equations, describing geometric transformations of vectors, among many other things
    \end{itemize}
    \item Linearity - addition and scalar multiplication can be done before or after an operation and still result in the same thing
    \item Lack of linear independence means that there exists scalars $\alpha_1,\dots,\alpha_n$ not all equal to zero such that $\sum_{i}^n \alpha_ix_i = \bm{0}$
    \begin{itemize}
        \item Suppose that $\alpha_j \neq 0$ for some $j$
        \[
        \bm{x}_j = -\frac{\alpha_i}{\alpha_j} \bm{x}_1 - \dots - \frac{\alpha_{j-1}}{\alpha_j} \bm{x}_{j-1} - \frac{\alpha_{j+1}}{\alpha_j} \bm{x}_{j+1} -\dots - \frac{\alpha_n}{\alpha_j} \bm{x}_n
        \]
        \item This shows that $\bm{x}_j$ can be written as a linear combination of the other vectors
    \end{itemize}
    \item For 2D, if two vectors are linearly dependent, they lie on the same line
    \item For 3D, if three vectors are linearly dependent, they lie on the same plane
    \item Exercise 5.5. (a) Are the following sets of vectors linearly independent?
    \[
    x_1 = \left[ \begin{array}{c}
        2 \\ -1 \\ 3
    \end{array} \right], \quad x_2 = \left[ \begin{array}{c}
        1 \\ 1 \\ -2
    \end{array} \right], \quad x_3 = \left[ \begin{array}{c}
        3 \\ -3 \\ 8
    \end{array} \right]
    \]
    \begin{itemize}
        \item Note that $2\bm{x}_1 - \bm{x}_2 - \bm{x}_3 = \bm{0}$ so these vectors are not linearly independent
        \item Another way is to solve the system of equations:
        \begin{flalign*}
            2\alpha_1 + \alpha_2 + 3\alpha_3 &= 0 \\
            -\alpha_1 + \alpha -3\alpha &= 0 \\
            3\alpha_1 - 2\alpha_2 + 8\alpha_3 &= 0
        \end{flalign*}
    \end{itemize}
    \item Exercise 5.5. (b) Are the following sets of vectors linearly independent?
    \[
    x_1 = \left[ \begin{array}{c}
        1 \\ 2 \\ 1 \\ 0 \\ 0
    \end{array} \right], \quad x_2 = \left[ \begin{array}{c}
        1 \\ 1 \\ 0 \\ 1 \\ 1
    \end{array} \right], \quad x_3 = \left[ \begin{array}{c}
        1 \\ 0 \\ 0 \\ 1 \\ 1
    \end{array} \right]
    \]
    \begin{itemize}
        \item Suppose that a linear combination equals zero. Then
        \[
        \alpha_1 \left[ \begin{array}{c}
            1 \\ 2 \\ 1 \\ 0 \\ 0
        \end{array} \right] + \alpha_2 \left[ \begin{array}{c}
            1 \\ 1 \\ 0 \\ 1 \\ 1
        \end{array} \right] + \alpha_3 \left[ \begin{array}{c}
            1 \\ 0 \\ 0 \\ 1 \\ 1
        \end{array} \right] = \left[ \begin{array}{c}
            0 \\ 0 \\ 0 \\ 0 \\ 0
        \end{array} \right]
        \]
        \item Notice from third row that $\alpha_1 = 0$
        \item Use this in the second row to immediately get $\alpha_2 = 0$
        \item From the final row, we get that $\alpha_3 = 0$
        \item Therefore, the vectors must be linearly independent.
    \end{itemize}
    \item Exercise 5.6. Write
    \[
    y = \left[ \begin{array}{c}
            1 \\ -2 \\ 5
        \end{array} \right]
    \]
    as a linear combination of
    \[
    x_1 = \left[ \begin{array}{c}
        1 \\ 1 \\ 1
    \end{array} \right], \quad x_2 = \left[ \begin{array}{c}
        1 \\ 2 \\ 3
    \end{array} \right], \quad x_3 = \left[ \begin{array}{c}
        2 \\ -1 \\ 1
    \end{array} \right]
    \]
    \begin{itemize}
        \item We can rewrite the problem as
        \[
        \left[ \begin{array}{ccc}
            1 & 1 & 2 \\
            1 & 2 & -1 \\
            1 & 3 & 1
        \end{array} \right] \left[ \begin{array}{c}
            \alpha_1 \\
            \alpha_2 \\
            \alpha_3
        \end{array} \right] = \left[ \begin{array}{c}
            1 \\
            -2 \\
            5
        \end{array} \right]
        \]
        \item Matrix maps vectors into linear combination of its columns
        \item Inverting matrix is hard and tedious
        \item Instead, we solve by eliminating variables from system of equations one by one using Gaussian elimination
        \[
        \alpha_1 \left[ \begin{array}{c}
            1 \\ 1 \\ 1
        \end{array} \right] + \alpha_2 \left[ \begin{array}{c}
            1 \\ 2 \\ 3
        \end{array} \right] + \alpha_3 \left[ \begin{array}{c}
            2 \\ -1 \\ 1
        \end{array} \right] = \left[ \begin{array}{c}
            1 \\ -2 \\ 5
        \end{array} \right]
        \]
        \item The unique solution is
        \[
        \alpha_1 = -6, \alpha_2 = 3, \alpha_3 = 2
        \]
    \end{itemize}
    \item Every matrix can be uniquely represented by its set of eigenvalues and eigenvectors
    \item Exercise 5.7. (a) Consider the general $2 \times 2$ matrix
    \[
    A = \left[ \begin{array}{cc}
        a_{11} & a_{12} \\
        a_{21} & a_{22}
    \end{array} \right]
    \]
    The trace of the above matrix is just the sum of all the elements along the leading diagonal, i.e.
    \[
    Tr(A) = \sum_{i=1}^2 a_{ii}
    \]
    Show that the eigenvalues of $A$ satisfy the polynomial equation $f(\lambda)=0$, where
    \[
    f(\lambda) = \lambda^2 - (Tr(A))\lambda + \det(A)
    \]
    \begin{itemize}
        \item Start with the characteristic equation
        \[
        \det(A - \lambda I) = 0
        \]
        \item Substituting in the general form of $2 \times 2$ $A$, we find that
        \begin{flalign*}
            0 &= \det(A - \lambda I) = \det\left( \left[ \begin{array}{cc}
                a_{11} - \lambda & a_{12} \\
                a_{21} & a_{22} - \lambda
            \end{array} \right] \right) \\
            &= (a_{11}-\lambda)(a_{22}-\lambda) - a_{12}a_{21} \\
            &= \lambda^2 - (a_{11} + a_{22}) \lambda + a_{11}a_{22} - a_{12}a_{21} \\
            &= \lambda^2 - (Tr(A))\lambda + \det(A) \\
            &= f(\lambda)
        \end{flalign*}
        as desired. Now we know that the eigenvalues of $A$ satisfy $f(\lambda)$.
    \end{itemize}
    \item Exercise 5.7. (b) Now show that $f(A) = \bm{0}$, where $\bm{0}$ is a $2 \times 2$ zero matrix.
    \begin{itemize}
        \item Note that $f(A) = A^2 - (Tr(A))A + \det(A)$
        \[
        = \left[ \begin{array}{cc}
            a_{11} & a_{12} \\
            a_{21} & a_{22}
        \end{array} \right]^2 - (a_{11} + a_{22}) \left[ \begin{array}{cc}
                a_{11} & a_{12} \\
                a_{21} & a_{22}
            \end{array} \right] + (a_{11}a_{22} - a_{12}a_{21})I_{2\times2}
        \]
        \item For example, the top left component simplifies as
        \[
        a_{11}^2 + a_{12}a_{21} - a_{11} (a_{11}+a_{22}) + a_{11}a_{22} - a_{12}a_{21} = 0
        \]
    \end{itemize}
    \item Similar results hold for larger matrices as well, and this is called the Cayley-Hamilton Theorem
    \begin{itemize}
        \item For any matrix $A$, we can define the polynomial $f(\lambda)=\det(A-\lambda I)$ and $f(A)=0$ always
        \item This theorem has infamous erroneous one-line proof:
        \[
        f(A)=\det(A-AI)=\det(A-A)=\det(0)=0
        \]
        \item This proof is wrong because $\lambda$ is a scalar and so we can't necessary substitute a matrix in without being careful. The determinant of a matrix must be a scalar, but $f(A)$ should be a matrix
    \end{itemize}
    \item Exercise 5.9. (a) In there dimensions, the vectors $x,y,a$ satisfy
    \[
    x + y(x \cdot y) = a
    \]
    Show that
    \[
    (x \cdot y)^2 = \frac{\vert a \vert^2 - \vert x \vert^2}{2 + \vert y \vert^2}
    \]
    \begin{itemize}
        \item We compute
        \begin{flalign*}
            \vert a \vert^2 &= a \cdot a = (x + y(x \cdot y)) \cdot (x + y(x \cdot y)) \\
            &= x \cdot x + 2(x \cdot y)^2 + (y \cdot y)(x \cdot y)^2 \\
            &= \vert x \vert^2 + (2 + \vert y \vert^2)(x \cdot y)^2 \\
            &\Rightarrow (x \cdot y)^2 = \frac{\vert a \vert^2 - \vert x \vert^2}{2 + \vert y \vert^2}
        \end{flalign*}
    \end{itemize}
    \item Exercise 5.9. (b) Use an inequality involving $x \cdot y$ and the lengths $x$ and $y$ to deduce that
    \[
    \vert x \vert \left( 1 + \vert y \vert^2 \right) \ge \vert a \vert \ge \vert x \vert
    \]
    \begin{itemize}
        \item Note that $0 \le \vert x \cdot y \vert \le \vert x \vert \vert y \vert$. Using this, we find
        \begin{flalign*}
            &0 \le \frac{\vert a \vert^2 - \vert x \vert^2}{2 + \vert y \vert^2} \le \vert x \vert^2 \vert y \vert^2 \\
            \Longrightarrow \ &0 \le \vert a \vert^2 - \vert x \vert^2 \le \vert x \vert^2 \vert y \vert^2 (2+\vert y \vert^2) \\
            \Longrightarrow \ &\vert x \vert^2 \le \vert a \vert^2 \le \vert x \vert^2 (1 + 2 \vert y \vert^2 + \vert y \vert^2) \\
            \Longrightarrow \ &\vert x \vert^2 \le \vert a \vert^2 \le \vert x \vert^2 (1 + \vert y \vert^2)^2 \\
        \end{flalign*}
        From this, taking the square root, we obtain that
        \[
        \vert x \vert \le \vert a \vert \le \vert x \vert (1 + \vert y \vert^2)
        \]
    \end{itemize}
    \item Exercise 5.9. (c) Explain the circumstances in which either of the inequalities above become equalities, and describe the relation between $x,y,a$ in these circumstances.
    \begin{itemize}
        \item Recall that $x \cdot y = \vert x \vert \vert y \vert \cos(\theta)$. So if $x$ and $y$ are orthogonal, then $x \cdot y = 0$. In this case, we obtain $a = x$, which means their lengths are also equal ($\vert a \vert = \vert x \vert$).
        \item On the other hand, if $x$ and $y$ are parellel, then $x \cdot y = \vert x \vert \vert y \vert$. Then, the length of $a$ is equal to the sum of the lengths of $x$ and $y$. This gives
        \[
        \vert a \vert = \vert x \vert + \vert x \vert \vert y \vert^2 = \vert x \vert \left( 1+ \vert y \vert^2 \right).
        \]
    \end{itemize}
\end{itemize}

\section{Lecture 11 (25/11/2025)}
\begin{itemize}
    \item $L = \frac{1}{2} \sigma^2 \frac{d^2}{dx^2} + \mu \frac{d}{dx} - r, L: V \rightarrow V$
    \begin{itemize}
        \item $V$ is the space of twice differentiable functions (cannot describe it properly; don't know enough)
        \item Want to find $Ker(L)$
        \item We also show that $L$ is a linear map
        \begin{flalign*}
            &(\forall f \in V)(\forall V \in \mathbb{C}) \ L(\lambda f) = \lambda L(f) \\
            &L(\lambda f) = \left( \frac{1}{2} \sigma^2 \frac{d^2}{dx^2} + \mu \frac{d}{dx} - r \right) \lambda f = \lambda \left( \frac{1}{2} \sigma^2 \frac{d^2}{dx^2} + \mu \frac{d}{dx} - r \right) f = \lambda L(f) \\
            &(\forall f,g \in V) L(f+g) = Lf + Lg
        \end{flalign*}
        \item $ker(L) = \{ f \in V : L(f) = 0 \in V \}$
        \begin{itemize}
            \item $\frac{1}{2} \sigma^2 f''(x) + \mu f'(x) - rf(x) = 0$
            \item Guess $f(x) = e^{kx}$
            \item $\left( \frac{1}{2} \sigma^2 k^2 + \mu k - r \right) e^{kx} = 0$
            \item $\Rightarrow \frac{1}{2} \sigma^2k^2 + \mu k -r = 0$ (characteristic equation)
            \item $\Rightarrow k^2 + \frac{2\mu}{\sigma^2} k - \frac{2r}{\sigma^2} = 0$
            \item $\Rightarrow \left( k + \frac{\mu}{\sigma^2} \right)^2 - \frac{\mu^2}{\sigma^4} - \frac{2r\sigma^2}{\sigma^4} = 0$
            \item $\Rightarrow k = \frac{-\mu \pm \sqrt{\mu^2 - 2r\sigma^2}}{\sigma^2}$
            \begin{itemize}
                \item Two solutions to go into kernel
                \item $e^{k(-x)}$ and $e^{k(+x)}$
            \end{itemize}
            \item $ker(L) = \{ A_{-}e^{k(-x)} + A_{+} e^{k(+x)} : A_{-}, A_{+} \in \mathbb{C} \}$
        \end{itemize}
    \end{itemize}
    \item Definition (Inner Product) - an inner product on a vector space $V$ over a field $\mathbb{F}$ is a map $\langle \cdot, \cdot \rangle : V \times V \rightarrow \mathbb{F}$
    \begin{itemize}
        \item Positive definiteness: $(\forall x \in V) \ \langle x, x \rangle \ge 0$ with equality iff $x = 0$
        \item Homogeneity: $(\forall a \in \mathbb{F})(\forall x,y \in V) \ \langle ax, y \rangle = a \langle x,y \rangle$
        \item Additivity: $(\forall x,y,z \in V) \ \langle x+y,z \rangle = \langle x,z \rangle + \langle y,z \rangle$
        \item Conjugate Symmetry: $(\forall x,y \in V) \ \langle x,y \rangle = \overline{\langle y,x \rangle}$
    \end{itemize}
    \item Vector space $V$ over $\mathbb{F}$ with an inner product $\langle \cdot,\cdot \rangle$ can be denoted as the inner product space: ($(V(\mathbb{F}), \langle \cdot,\cdot \rangle)$ or $(V,\langle \cdot,\cdot \rangle)$)
    \item Definition (Orthogonal) - in an inner product space $(V,\langle \cdot,\cdot \rangle)$, vectors $x,y \in V \backslash \{0\}$ are orthogonal if
    \[
    \langle x,y \rangle = 0
    \]
    \item Definition (Orthogonal Complement) - with an inner product space $(V, \langle \cdot,\cdot \rangle)$, for a subspace $U$ of $V$:
    \[
    U^{\perp} = \{ y \in V : \forall x \in U, \langle y,x \rangle = 0 \}
    \]
    \begin{itemize}
        \item Key Property: $U^{\perp}$ is also a subspace of $V$
    \end{itemize}
    \item Example:
    \begin{itemize}
        \item $V = \mathbb{R}^2, \mathbb{F} = \mathbb{R}, \quad X = \{ \left[ \begin{array}{cc}
            x \\ y
        \end{array} \right] \in \mathbb{R}^2 : x \in \mathbb{R} \wedge y = 0 \}, \quad Y = \{ \left[ \begin{array}{cc}
            x \\ y
        \end{array} \right] \in \mathbb{R}^2 : y \in \mathbb{R} \wedge x = 0 \}$
        \item $X^{\perp}=Y, \quad (X^{\perp})^{\perp} = Y^{\perp} = X$
        \item Example of Projection Theorem: $X \oplus Y = \mathbb{R}^2 \Longleftrightarrow X \oplus X^{\perp} = \mathbb{R}^2$
    \end{itemize}
    \item Definition (Norm) - a norm on $V$ is a function $\vert \vert \cdot \vert \vert : V \rightarrow \mathbb{R}$ satisfying:
    \begin{itemize}
        \item $(\forall x \in \mathbb{R}) \ \vert \vert x \vert \vert \ge 0$ with equality iff $x = 0$
        \item $\vert \vert x+y \vert \vert \le \vert \vert x \vert \vert + \vert \vert y \vert \vert$ (triangle inequality)
        \item $\vert \vert \lambda x \vert \vert = \vert \lambda \vert \vert \vert x \vert \vert$
        \item Connection: an inner product induces a norm: $\vert \vert x \vert \vert = \sqrt{\langle x,x \rangle}$
    \end{itemize}
    \item Definition (Metric) - a metric on a set $X$ is $d : X \times X \rightarrow \mathbb{R}$ satisfying:
    \begin{itemize}
        \item $d(x,y) \ge 0$ with equality iff $x = y$
        \item $d(x,y) = d(y,x)$ (symmetry)
        \item $d(x,z) \le d(x,y) + d(y,z)$ (triangle inequality)
        \item Think of the metric as an advanced measure of distance
    \end{itemize}
    \item Definition (Orthogonal Basis) - an orthogonal basis consists of mutually orthogonal vectors
    \item Definition (Orthonormal Basis) - an orthonormal basis consists of mutually orthogonal vectors, each with unit norm
    \begin{itemize}
        \item Example: the standard basis of $\mathbb{R}^2$
        \[
        \{(1,0)^T, (0,1)^T\}
        \]
        is orthonormal with respect to the dot product
    \end{itemize}
    \item Theorem (Projection Theorem) - for any vector space $V$ and subspace $U \subset V:$
    \[
    (\forall v \in V)(\exists !u \in U)(\exists !w \in U^{\perp}) \ v = u+w
    \]
    \begin{itemize}
        \item Alternative statement: $V = U \oplus U^{\perp}$
    \end{itemize}
    \item Definition (Orthogonal Projection) - the orthogonal projection of $v \in V$ onto subspace $U$ is the unique $w \in U$ such that $v = w+u$ where $u \in U^{\perp}$
    \item Example of finding orthonormal basis for vector space:
    \begin{itemize}
        \item Consider the set of polynomials of degree less than equal to 2 with real coefficients:
        \[
        V = \{ a_0 + a_1 x + a_2 x^2 : a_0,a_1,a_2 \in \mathbb{R} \} \text{ over field } \mathbb{R}
        \]
        \item Define the inner product to be $\langle a,b \rangle = \int_0^1 ab dx$.
        \item $\langle 1, a_0 + a_1x \rangle = \int_0^1 1(a_0 + a_1x)dx = \left[ a_0x + \frac{1}{2} a_1 x^2 \right]_0^1 = a_0 + \frac{1}{2}a_1 = 0 \Rightarrow a_0 = -\frac{1}{2} a_1$
        \item $a_0 + a_1x = a_0 - 2a_0x = a_0(1-2x)$
        \item Now we pick $a_0$ so that the norm is equal to 1
        \item $\vert \vert a_0 (1-2x) \vert \vert = a_0 \vert \vert 1-2x \vert \vert = a_0 \sqrt{\langle 1-2x, 1-2x \rangle} = a_0 \sqrt{\langle 1, 1-2x \rangle - \langle 2x,1-2x \rangle}$ \\
        $= a_0 \sqrt{\langle 1,1 \rangle - 2 \langle 1,x \rangle - 2 \langle x,1 \rangle + 4 \langle x,x \rangle} = a_0 \sqrt{\langle 1,1 \rangle - 4 \langle 1,x \rangle + 4 \langle x,x \rangle}$
        \item $\langle 1,1 \rangle = \int_0^1 1dx = 1$
        \item $\langle 1,x \rangle = \int_0^1 xdx = \frac{1}{2}$
        \item $\langle x,x \rangle = \int_0^1 x^2 dx = \frac{1}{3}$
        \item $\vert \vert a_0 (1-2x) \vert \vert = a_0 \sqrt{1 - 2 + \frac{4}{3}} = a_0 \sqrt{\frac{1}{3}} = \frac{a_0}{\sqrt{3}}$
        \item $a_0 = \sqrt{3}$
        \item $\langle 1, \sqrt{3}(1-2x) \rangle = 0$
        \item $B = \{1,x,x^2\}$ is a spanning set of $V$ and a linearly independent set
        \begin{itemize}
            \item Therefore, it is a basis
        \end{itemize}
        \item Find $a_0,a_1,a_2$ such that $\langle 1,a_0+a_1x+a_2x^2 \rangle = 0 \Longleftrightarrow \langle 1-2x, a_0 + a_1x + a_2x^2 \rangle = 0$
        \item $e_1(x) = 1$
        \item $e_2(x) = \sqrt{3} (2x-1)$
        \item $e_3(x) = 6 \sqrt{5} \left( x^2 - x + \frac{1}{6} \right)$
        \begin{itemize}
            \item How to find $e_3(x)$:
            \begin{flalign*}
                a_0+a_1x+a_2x^2 &= a_2(\frac{a_0}{a_2} + \frac{a_1}{a_2}x + x^2) \\
                \text{1. Find quadratic such that } \langle 1,x^2+bx+c \rangle &= 0 \\
                \langle 1, x^2 \rangle+b\langle 1,x \rangle + c\langle 1,1 \rangle &= 0 \\
                \int_0^1 x^2 dx = \frac{1}{3}, \quad \int_0^1 xdx =\frac{1}{2}, \quad \int_0^1 1 dx = 1 \\
                \frac{1}{3} + b\frac{1}{2} + c &= 0 \\
                \Rightarrow 2 + 3b + 6c &= 0 \\
                \text{2. Find quadratic such that } \langle 2x-1,x^2+bx+c \rangle &= 0 \\
                \langle 2x-1, x^2 \rangle+b\langle 2x-1,x \rangle + c\langle 2x-1,1 \rangle &= 0 \\
                2\langle x, x^2 \rangle - \langle 1,x^2 \rangle + b \left( 2\langle x,x \rangle - \langle 1,x \rangle \right) + c \left( 2 \langle x,1 \rangle - \langle 1,1 \rangle \right) &= 0 \\
                \int_0^1 x^3 dx = \frac{1}{4} \quad \int_0^1 x^2 dx = \frac{1}{3}, \quad \int_0^1 xdx =\frac{1}{2}, \quad \int_0^1 1 dx = 1 \\
                2\frac{1}{4} - \frac{1}{3} + b \left( 2\frac{1}{3} - \frac{1}{2} \right) + c \left( 2 \frac{1}{2} - 1 \right) &= 0 \\
                \Rightarrow \frac{1}{6} + \frac{1}{6}b &= 0 \\
                \Rightarrow 1 + b &= 0 \Rightarrow b = -1 \\
                \Rightarrow 2 + 3(-1) + 6c = 6c - 1 &= 0 \Rightarrow c = \frac{1}{6} \\
                x^2 - x + \frac{1}{6} \text{ is orthogonal to } e_1(x) \text{ and } e_2(x) \\
                \text{Find } \vert \vert x^2 - x + \frac{1}{6} \vert \vert &= \langle x^2-x+\frac{1}{6},x^2-x+\frac{1}{6} \rangle \\
                &= (6\sqrt{5})^{-2} \Rightarrow \vert \vert x^2-x+\frac{1}{6} \vert \vert = (6\sqrt{5})^{-1} \\
                \frac{x^2-x+\frac{1}{6}}{\vert \vert x^2-x+\frac{1}{6} \vert \vert} = \frac{x^2-x+\frac{1}{6}}{(6\sqrt{5})^{-1}} = 6\sqrt{5} (x^2-x+\frac{1}{6})
            \end{flalign*}
            \item $\{e_1(x),e_2(x),e_3(x)\}$ is an orthonormal basis of $V$ with dimensionality 3
        \end{itemize}
    \end{itemize}
\end{itemize}
\subsection{Lecture Notes}
\begin{itemize}
    \item Scalar (Dot) Product Review
    \begin{itemize}
        \item Recall: $x,y \in \mathbb{R}^n, \quad x \cdot y = \sum_{i=1}^n x_i y_i, \quad \vert \vert x \vert \vert = \sqrt{x \cdot x}$
        \item Consider $a,b \in \mathbb{R}^n$
        \item Define $\hat{b} = \frac{b}{\vert \vert b \vert \vert} \Rightarrow \vert \vert \hat{b}  \vert \vert = 1$
        \item Think about $a \cdot \hat{b}$
        \begin{itemize}
            \item If $\theta$ is angle between $a$ and $b$, we can define the length of the projection as follows: $\vert proj_{\hat{b}}(a) \vert = \vert \vert a \vert \vert \cos(\theta) = \vert a \cdot \hat{b} \vert$
            \item The horizontal component of $a$ is $(a \cdot \hat{b})\hat{b}$
            \begin{itemize}
                \item The vertical component is $a - (a \cdot \hat{b})\hat{b}$
            \end{itemize}
        \end{itemize}
        \item Define the orthogonal projection of $x$ onto $b$: $p_b : \mathbb{R}^n \rightarrow \mathbb{R}^n. \quad \forall x \in \mathbb{R}^n$
        \begin{itemize}
            \item $p_b(x) = (x \cdot \hat{b})\hat{b}$
        \end{itemize}
        \item $a = p_b(a) + a - p_b(a) = (p_b + I - p_b)a$
    \end{itemize}
\end{itemize}

\section{Lecture 12 (27/11/2025)}
\begin{itemize}
    \item Previously, we see that $(\forall x \in \mathbb{R}^n) \ Proj_b(x) = (x \cdot \hat{b}) \hat{b}), \hat{b} = \frac{b}{\vert \vert b \vert \vert}$
    \item Exercise: $(a - (a \cdot \hat{b}) \hat{b}) \cdot ((a \cdot \hat{b}) \hat{b}) = 0$
    \item $a \cdot \hat{b} = \vert \vert a \vert \vert cos(\theta)$
    \item $\vert \vert a - (a \cdot \hat{b}) \hat{b} \vert \vert = \vert \vert a \vert \vert sin(\theta)$
    \item Next: orthogonal projection onto a plane
    \begin{itemize}
        \item $b = \left[ \begin{array}{c}
            b_1 \\ b_2 \\ b_3
        \end{array}\right]$
        \item Figure out orthogonal projection of $b$ onto xy plane:
        \item $P_{xy} \left( \left[ \begin{array}{c}
            b_1 \\ b_2 \\ b_3
        \end{array}\right] \right) = \left[ \begin{array}{ccc}
            1 & 0 & 0 \\
            0 & 1 & 0 \\
            0 & 0 & 0
        \end{array}\right] \left[ \begin{array}{c}
            b_1 \\ b_2 \\ b_3
        \end{array}\right] = \left[ \begin{array}{c}
            b_1 \\ b_2 \\ 0
        \end{array}\right]$
        \item $P = \left[ \begin{array}{ccc}
            1 & 0 & 0 \\
            0 & 1 & 0 \\
            0 & 0 & 0
        \end{array}\right] , \quad P^2 = P \Rightarrow (\forall n \in \mathbb{N}) \ P^n = P$
    \end{itemize}
    \item Definition (Idempotent) - a matrix $A \in \mathbb{R}^{n \times n}$ is idempotent if $A^2 = A$
    \item Rewrite the above question of the orthogonal projection of $b$ onto the $xy$ plane as the orthogonal projection of $b = \left[ \begin{array}{c}
        b_1 \\ b_2 \\ b_3
    \end{array}\right] \in \mathbb{R}^3$ onto $span(a_1,a_2) = \{x_1a_1+x_2a_2 : x_1,x_2 \in \mathbb{R} \}$ where $a_1,a_2 \in \mathbb{R}^3$
    \item Define $A = \left[ \begin{array}{cc}
        a_1 & a_2
    \end{array}\right] $
    \begin{itemize}
        \item Note that $Ax = A\left[ \begin{array}{c}
            x_1 \\
            x_2
        \end{array}\right] = x_1a_1 + x_2a_2$
        \item Project $b$ onto column space (span of $a_1,a_2$) of $A$ (want $Ax$ to be orthogonal projection)
        \item Decompose $b = Ax + \epsilon$
        \item Want $\epsilon \perp Ax \Rightarrow \underset{=\epsilon}{\underbrace{(b-Ax)}} \cdot Ax = 0$
        \item Recall $a \cdot b = a^Tb$, so this means $(Ax)^\top \epsilon = x^\top A^\top \epsilon = 0$
        \item Want $\epsilon \cdot a_1 = 0$ and $\epsilon \cdot a_2 = 0$
        \begin{itemize}
            \item $\epsilon \cdot Ax = \epsilon \cdot (x_1a_1+x_2a_2) = 0$
            \item $a_1^T\epsilon = 0 \Longleftrightarrow a_1^T \epsilon = 0 \Longleftrightarrow A^T\epsilon = 0$
        \end{itemize}
        \item Recall $\epsilon = b - Ax \Longleftrightarrow A^T(b-Ax) = 0$
        \begin{itemize}
            \item $A^Tb - A^TAx = 0 \Rightarrow A^Tb = A^TAx$
            \item If $A^TA$ is invertible,
            \[
            x = (A^TA)^{-1} A^Tb
            \]
            \item $Ax = A(A^TA)^{-1}A^Tb$ is the orthogonal projection of $b$ onto the $span(A)$
            \item Projection matrix:
            \[
            P_A = A(A^TA)^{-1}A^T
            \]
            \begin{itemize}
                \item Properties: $P_A^2=P_a$ (idempotent) and $P_A^\top=P_A$ (symmetric)
            \end{itemize}
        \end{itemize}
    \end{itemize}
    \item Check: $a_1 = \left[ \begin{array}{c}
        1 \\ 0 \\ 0
    \end{array}\right]$ and $a_2 = \left[ \begin{array}{c}
        0 \\ 1 \\ 0
    \end{array}\right]$, $span(a_1,a_2) = \left\{ \left[ \begin{array}{c}
        x \\ y \\ z
    \end{array}\right]  \in \mathbb{R}^3 : x,y \in \mathbb{R} \wedge z =0 \right\}$
    \begin{itemize}
        \item $A = \left[ \begin{array}{cc}
            1 & 0 \\
            0 & 1 \\
            0 & 0
        \end{array}\right]$, $A^T = \left[ \begin{array}{ccc}
            1 & 0 & 0\\
            0 & 1 & 0
        \end{array}\right]$, $P_A = A(A^TA)^{-1}A^T$
    \end{itemize}
    \item $A^TA = \left[ \begin{array}{cc}
        1 & 0 \\
        0 & 1
    \end{array}\right] = \bm{I}_2$, $A(A^TA)^{-1}A^T = AIA^T = AA^T = \left[ \begin{array}{ccc}
        1 & 0 & 0 \\
        0 & 1 & 0 \\
        0 & 0 & 0
    \end{array}\right]$
    \item Linear Regression:
    \begin{itemize}
        \item Consider $n$ measurements of a variable $y = (y_1, \dots, y_n)^T \in \mathbb{R}^n$, which we wish to explain in terms of the $n$ observations of another variable $x = (x_1, \dots, x_n)^T \in \mathbb{R}^n$ via the following set of linear equations:
        \begin{flalign*}
            \left\{ \begin{array}{ccc}
                 y_1=\beta_1 &+\beta_2x_1 &+\epsilon_1  \\
                 \vdots &\vdots& \vdots \\
                 y_n = \beta_1 &+\beta_2x_n &+\epsilon_n
            \end{array} \right.           
        \end{flalign*}
        where $\epsilon = (\epsilon_1,\dots,\epsilon_n)^T \in \mathbb{R}^n$ is the vector of errors
        \item We are trying to find $\beta_1,\beta_2$ in $y=\beta_1 + \beta_2 x$ to minimize $\sqrt{\frac{1}{n}\sum_{i=1}^n \epsilon_i^2}$ (root mean squared error/MSE) or equivalently minimize $\vert \vert \epsilon \vert \vert$ for the scalar inner product
        \begin{itemize}
            \item The error is as follows: $\epsilon_n = y_n - (\beta_1+\beta_2x)$
            \item Equivalent to minimizing $\sum_{i=1}^n \epsilon_i^2 = \epsilon \cdot \epsilon$ where $\epsilon = (\epsilon_1, \dots, \epsilon_n)^T \in \mathbb{R}^n$
        \end{itemize}
        \item $y = \beta_1 \bm{1} + \beta_2 x + \epsilon$
        \begin{itemize}
            \item Think of linear regression as the projection of $y$ onto $span(\bm{1},x)$
            \item Choose $\beta_1, \beta_2$ such that $X = \left[ \begin{array}{cc}
                \bm{1} & x            \end{array}\right]
            \left[ \begin{array}{c}
                \beta_1 \\ \beta_2
            \end{array}\right] = \beta_1 \bm{1} + \beta_2 x$ is orthogonal to $\epsilon$
            \item By replacing the dot product with another inner product, you can produce other regression
            \begin{itemize}
                \item OLS: Euclidean norm (standard dot product)
                \item WLS: Weighted norm with diagonal matrix $W$
                \item GLS: General symmetric matrix $W$
            \end{itemize}
        \end{itemize}
    \end{itemize}
\end{itemize}

\section{Tutorial 8 (27/11/2025)}
\begin{itemize}
    \item To check if $W$ is a subspace, you only have to check three conditions, namely the inclusion of zero and closure under addition and scalar multiplication
    \begin{itemize}
        \item The other axioms are automatically inherited from $V$ being vector space
    \end{itemize}
    \item Quick Question: For a linear map $L: V \rightarrow W$, is the kernel
    \[
    \ker(L) = \{ v \in V : L(v) = 0 \} \subseteq V
    \]
    a subspace of $V$?
    \begin{itemize}
        \item Note that $L(0) = L(0 \cdot u) = 0 L(u) = 0$ so $0 \in \ker(L)$
        \item By linearity of $L$, $u,v \in \ker(L) \Longrightarrow L(u+v) = L(u) + L(v) = 0 \Longrightarrow u+v \in \ker(L)$
        \item Similarly, by linearity of $L$ again, $u \in \ker(L) \wedge \lambda \in \mathbb{F} \Longrightarrow L(\lambda u) = \lambda L(u) = 0 \Longrightarrow \lambda u \in \ker(L)$
    \end{itemize}
    \item Exercise 6.2. Find the orthogonal projection of the column vector $(3,1)^\top$ onto the vector $(1,1)^\top$ using the dot product.
    \begin{itemize}
        \item Remember that we are projecting onto the direction described by $b = (1,1)^\top$, so we first construct a unit vector in that direction:
        \[
        \hat{b} = \frac{b}{\vert \vert b \vert \vert} = \frac{1}{\sqrt{2}} (1,1)^\top
        \]
        \item Then, the orthogonal projection of $a = (3,1)T\top$ onto $b$ is
        \[
        proj_b(a) = (a\cdot\hat{b})\hat{b} = 2\sqrt{2} \cdot \frac{1}{\sqrt{2}}(1,1)^\top = (2,2)^\top
        \]
    \end{itemize}
    \item Exercise 6.7. (a) The general unit vector in $\mathbb{R}^2$ can be written as $\hat{n} = (\cos(\theta),\sin(\theta))^\top$ with respect to the standard basis, where $\theta \in [0,2\pi)$. Work out the matrix $P$ which represents an orthogonal projection onto $\hat{n}$ (with respect to the standard basis).
    \begin{itemize}
        \item We are going to look at this from two angles. For this part, let us consider any $v \in \mathbb{R}^2$. This can be written in the standard basis as $v = (v_1,v_2)^\top = v_1e_1+v_2e_2$. Then, by linearity,
        \[
        Pv = v_1(Pe_1) + v_2(Pe_2)
        \]
        \item Recall that a matrix maps a vector to the corresponding linear combination of its columns. Thus, the above tells us that the matrix of $P$ must look like:
        \[
        P = \left[ \begin{array}{cc}
            \uparrow & \uparrow \\
            Pe_1 & Pe_2 \\
            \downarrow & \downarrow
        \end{array}\right]
        \]
        i.e. with columns $Pe_1$ and $Pe_2$ written in standard basis themselves.
        \item We see that $e_1 \cdot \hat{n} = \cos(\theta)$ and $e_2 \cdot \hat{n} = \sin(\theta)$. So
        \[
        Pe_1 = (\cos^2(\theta), \cos(\theta)\sin(\theta))^\top
        \]
        and
        \[
        Pe_2 = (\cos(\theta)\sin(\theta),\sin^2(\theta))^\top
        \]
        \item Hence
        \[
        P = \left[ \begin{array}{cc}
            \cos^2(\theta) & \cos(\theta)\sin(\theta) \\
            \cos(\theta)\sin(\theta) & \sin^2(\theta)
        \end{array} \right]
        \]
        \item We can also say $A = (\cos(\theta), \sin(\theta))^\top$ and $P=A(A^\top A)^{-1}A^\top$ and calculate:
        \[
        \left[ \begin{array}{c}
            \cos(\theta) \\ \sin(\theta)
        \end{array}\right] \left( \left[ \begin{array}{cc}
            \cos(\theta) & \sin(\theta)
        \end{array}\right]\left[ \begin{array}{c}
            \cos(\theta) \\ \sin(\theta)
        \end{array}\right] \right)^{-1} \left[ \begin{array}{cc}
            \cos(\theta) & \sin(\theta)
        \end{array}\right]
        \]
        \[
        = (\cos^2(\theta) + \sin^2(\theta)) \left[ \begin{array}{cc}
            \cos^2(\theta) & \cos(\theta)\sin(\theta) \\
            \cos(\theta)\sin(\theta) & \sin^2(\theta)
        \end{array} \right] = \left[ \begin{array}{cc}
            \cos^2(\theta) & \cos(\theta)\sin(\theta) \\
            \cos(\theta)\sin(\theta) & \sin^2(\theta)
        \end{array} \right]
        \]
    \end{itemize}
    \item Exercise 6.7. (b) Show that $P^2 = P$.
    \begin{itemize}
        \item Here, we can use the explicit $P$ that we wrote down.
    \end{itemize}
    \begin{flalign*}
        &\left[ \begin{array}{cc}
            \cos^2(\theta) & \cos(\theta)\sin(\theta) \\
            \cos(\theta)\sin(\theta) & \sin^2(\theta)
        \end{array} \right] \left[ \begin{array}{cc}
            \cos^2(\theta) & \cos(\theta)\sin(\theta) \\
            \cos(\theta)\sin(\theta) & \sin^2(\theta)
        \end{array} \right] \\ &= \left[ \begin{array}{cc}
            \cos^4(\theta) + \cos^2(\theta)\sin^2(\theta) & \cos^3(\theta)\sin(\theta) + \cos(\theta)\sin^3(\theta) \\
            \cos^3(\theta)\sin(\theta) + \cos(\theta)\sin^3(\theta) & \sin^4(\theta) + \cos^2(\theta)\sin^2(\theta)
        \end{array} \right] \\
        &= \left[ \begin{array}{cc}
            \cos^2(\theta)(\cos^2(\theta) + \sin^2(\theta)) & \cos(\theta)\sin(\theta)(\cos^2(\theta) + \sin^2(\theta)) \\
            \cos(\theta)\sin(\theta)(\cos^2(\theta) + \sin^2(\theta)) & \sin^2(\theta)(\sin^2(\theta) + \cos^2(\theta))
        \end{array} \right] \\
        &= \left[ \begin{array}{cc}
            \cos^2(\theta) & \cos(\theta)\sin(\theta) \\
            \cos(\theta)\sin(\theta) & \sin^2(\theta)
        \end{array} \right]
    \end{flalign*}
    \begin{itemize}
        \item We can also look at the projection matrix from a second angle. Note that $\forall v \in \mathbb{R}^2$
        \begin{flalign*}
            Pv &= (v \cdot \hat{n}) \hat{n} = \hat{n} (\hat{n} \cdot v) \\
            &= \hat{n} (\hat{n}^\top v) \\
            &= (\hat{n} \hat{n}^\top) v
        \end{flalign*}
        \item Since this is true $\forall v \in \mathbb{R}^2$, we find that
        \[
        P = \hat{n}\hat{n}^\top
        \]
        \item This object is known as the outer product of $\hat{n}$ with itself. From this, we easily find
        \[
        P^2 = (\hat{n}\hat{n}^\top)(\hat{n}\hat{n}^\top) = \hat{n}(\hat{n}^\top \hat{n})\hat{n}^\top = \vert \vert \hat{n} \vert \vert^2 \hat{n}\hat{n}^\top = \hat{n}\hat{n}^\top = P 
        \]
    \end{itemize}
    \item Exercise 6.7. (c) Find the eigenvectors and eigenvalues of $P$.
    \begin{itemize}
        \item We can do this via the characteristic equation $\det(P-\lambda I) = 0$. However, we can instead observe that
        \[
        Pv=\lambda v \Longrightarrow P^2v = \lambda Pv = \lambda^2 v
        \]
        \item But since $P^2=P$, this tells us that
        \[
        \lambda v = \lambda^2 v
        \]
        \item Since eigenvectors must be non-zero, we conclude that the eigenvalues must be $\lambda = \{0,1\}$
        \item The corresponding eigenvectors can be identified by thinking about the nature of the projection operator. We are looking for a non-zero vector unchanged by projection $(\lambda=1)$ and a non-zero vector sent to zero by projection ($\lambda=0$).
        \begin{itemize}
            \item Vectors which have zero projection along $\hat{n}$ must be orthogonal to it. A unit vector orthogonal to $\hat{n}$ is $\hat{n}^\perp = (-\sin(\theta),\cos(\theta))^\top$, and so the eigenvectors with eignevalue $\lambda=0$ are proportional to $\hat{n}^\perp$.
            \item On the other hand, vectors which are unchanged by projection must already lie fully along $\hat{n}$. Thus, the eigenvectors with eigenvalue $\lambda=1$ are simply scalar multiples of $\hat{n} = (\cos(\theta),\sin(\theta))^\top$ itself.
        \end{itemize}
    \end{itemize}
    \item Exercise 6.7. (d) Find a $2 \times 2$ matrix $Q$ (not the zero matrix) such that $PQ=0$.
    \begin{itemize}
        \item There are many solutions to this. A particularly instructive solution is to first recall that
        \[
        P = \hat{n}\hat{n}^\top
        \]
        \item Let $Q$ be the projector onto the orthogonal complement, which in this case is the direction $\hat{n}^\top$. Then,
        \[
        Q = \hat{n}^\perp (\hat{n}^\perp)\top
        \]
        \item Then,
        \[
        PQ = (\hat{n}\hat{n}^\top)(\hat{n}^\perp (\hat{n}^\perp)^\top) = \hat{n}(\hat{n}^\top\hat{n}^\perp)(\hat{n}^\perp)^\top = 0
        \]
        since $\hat{n}$ and $\hat{n}^\top$ are orthogonal.
    \end{itemize}
    \item Exercise 6.4. The following matrices are representations of linear maps with respect to the standard basis of $\mathbb{R}^3$. Describe the action of each linear map in words.
    \begin{itemize}
        \item (a)
        \[
        T_1 = \left[ \begin{array}{ccc}
            \frac{1}{\sqrt{2}} & -\frac{1}{\sqrt{2}} & 0 \\
            \frac{1}{\sqrt{2}} & \sqrt{1}{\sqrt{2}} & 0 \\
            0 & 0 & 1
        \end{array} \right]
        \]
        \begin{itemize}
            \item This keeps the $z$ direction invariant (around the $z$-axis) while rotating things in the $xy$ plane $\frac{\pi}{4}$ radians counterclockwise
        \end{itemize}
        \item (b)
        \[
        T_2 = \left[ \begin{array}{ccc}
            \frac{1}{\sqrt{2}} & 0 & \frac{1}{\sqrt{2}} \\
            0 & 1 & 0 \\
            -\frac{1}{\sqrt{2}} & 0 & \frac{1}{\sqrt{2}}
        \end{array} \right]
        \]
        \begin{itemize}
            \item Rotation of $\frac{\pi}{4}$ radians clockwise about the $y$-axis
        \end{itemize}
        \item (c)
        \[
        T_3 = \left[ \begin{array}{ccc}
            -1 & 0 & 0 \\
            0 & 1 & 0 \\
            0 & 0 & 1
        \end{array} \right]
        \]
        \begin{itemize}
            \item Reflection through the $yz$-plane
        \end{itemize}
        \item (d)
        \[
        T_4 = \left[ \begin{array}{ccc}
            1 & 0 & 0 \\
            0 & -1 & 0 \\
            0 & 0 & 1
        \end{array} \right]
        \]
        \begin{itemize}
            \item Reflection through the $xz$-plane
        \end{itemize}
        \item (e)
        \[
        T_5 = \left[ \begin{array}{ccc}
            0 & 0 & -1 \\
            0 & 1 & 0 \\
            -1 & 0 & 0
        \end{array} \right]
        \]
        \begin{itemize}
            \item $T_5$ sends $x$ to $-z$ and $z$ to $-x$ without changing $y$. This is a reflection through the plan perpendicular to the vector $\hat{n} = \frac{1}{\sqrt{2}}(1,0,1)^\top$.
        \end{itemize}
    \end{itemize}
\end{itemize}

\section{Lecture 13 (02/12/2025)}
\begin{itemize}
    \item $x \in \mathbb{R}^n$, $n$ students, $x$ hours worked, $y$ exam scores
    \begin{itemize}
        \item $y = \bm{1}\beta_1 + x \beta_2 + \epsilon$
        \item O.L.S $\forall a,b \in \mathbb{R}^n$ define $\langle a,b \rangle = a \cdot b \Rightarrow \vert \vert a \vert \vert = \sqrt{a \cdot a} = \sum_{i=1}^n a_i^2$
        \item Choose $\beta = \left[ \begin{array}{c}
            \beta_1 \\
            \beta_2
        \end{array} \right]$ to minimize $\vert \vert \epsilon \vert \vert = \sum_{i=1}^n \epsilon_i^k$
        \item Geometrically, obtain the orthogonal projection of $y$ onto $\{ 1\lambda_1 + x \lambda_2 : \lambda_1, \lambda_2 \in \mathbb{R} \}$
        \item Suppose the $k$th observation is an outlier
        \begin{itemize}
            \item Choose $\beta_1$ and $\beta_2$ to minimize $\sum_{i=1}^n \epsilon_i^2 - \epsilon_k^2 = \sum_{i \in \{1, \dots, n\} \backslash \{k\}} e_i^k$
            \item This is just a special case of minimizing $\sum_{i=1}^n w_i \epsilon_i^2$
            \item All $w_i=1$ except $w_k=0$
            \item Note that $\sum_{i=1}^n w_i\epsilon_i^2 = \left[ \begin{array}{ccc}
                \epsilon_1 & \dots & \epsilon_n
            \end{array}\right] \left[ \begin{array}{cccc}
                w_1 & 0 & \dots & 0 \\
                0 & w_2 & \dots & 0 \\
                \vdots & \vdots & \ddots & \vdots \\
                0 & 0 & \cdots & w_n
            \end{array} \right] \left[ \begin{array}{c}
                \epsilon_1 \\
                \vdots \\
                \epsilon_n
            \end{array}\right]= \epsilon^T W \epsilon$ where $W = diag(w_1, \dots, w_n)$
            \item $W_{ij} = \left\{ \begin{array}{cc}
                w_i & \text{if } i =j \\
                0 & \text{otherwise}
            \end{array} \right.$
        \end{itemize}
    \end{itemize}
    \item Define $\langle a,b \rangle = a^TWb$
    \item $\vert \vert \epsilon \vert \vert = \sqrt{\langle \epsilon, \epsilon \rangle} = \sqrt{\epsilon^T W \epsilon}$
    \begin{itemize}
        \item This is called weighted least squares (WLS)
    \end{itemize}
    \item Make $W$ a real symmetric matrix ($W_{ij}=W_{ji}$)
    \[
    \left[ \begin{array}{cccc}
        w_{11} & w_{12} & \cdots & w_{1n} \\
        w_{21} = w_{12} & w_{22} & \cdots & w_{2n} \\
        \vdots & \vdots & \ddots & \vdots \\
        w_{n1}=w_{1n} & w_{n2} = w_{2n} & \cdots & w_{nn}
    \end{array} \right]
    \]
    \begin{itemize}
        \item $\vert \vert \epsilon \vert \vert = \sum_{i=1}^n \sum_{j=1}^n \epsilon_i W_{ij} \epsilon_j$
        \item Set $n=3$
        \[
        \left[ \begin{array}{ccc}
                \epsilon_1 & \epsilon_2 & \epsilon_3
            \end{array}\right] \left[ \begin{array}{ccc}
                w_{11} & w_{12} & w_{13} \\
                w_{21} & w_{22} & w_{23} \\
                w_{31} & w_{32} & w_{33}
            \end{array} \right] \left[ \begin{array}{c}
                \epsilon_1 \\
                \epsilon_2 \\
                \epsilon_3
            \end{array}\right] = \left[ \begin{array}{ccc}
                \epsilon_1 & \epsilon_2 & \epsilon_3
            \end{array} \right] \left[ \begin{array}{ccc}
                \sum_{j=1}^3 w_{ij} e_{j} \\
                w_{21}\epsilon_1 + w_{22}\epsilon_2 + w_{23}\epsilon_3 \\
                w_{31}\epsilon_1 + w_{32}\epsilon_2 + w_{33}\epsilon_3
            \end{array} \right]
        \]
        \item this is called generalized least squares (GLS)
    \end{itemize}
\end{itemize}
\subsection{Slide 6 (Multivariable Calculus: Functions of Several Variables)}
\begin{itemize}
    \item Definition (Limit Point) - a point $x$ is a limit point of a set $A$ if every $\epsilon$-neighborhood $V_\epsilon(x)$ of $x$ intersects the set at some point other than $x$, i.e.
    \[
    (\forall \epsilon > 0)(V_\epsilon(x) - \{x\}) \cap A \neq \emptyset
    \]
    \item Example: Limit Points
    \begin{itemize}
        \item Define the set $A = \{ x\in \mathbb{R} : 0 < x \le 1 \}$
        \item Case 1: $1 \in A$
        \begin{flalign*}
            V_\epsilon(1) &= \{ x : \vert x-1 \vert < \epsilon \} \\
            (V_\epsilon(1) - \{1\}) \cap A &= \{ x : 1-\epsilon < x < 1\} \neq \emptyset
        \end{flalign*}
        Hence, 1 is a limit point of $A$
        \item Case 2: $0 \notin A$
        \begin{flalign*}
            V_\epsilon(1) &= \{ x : \vert x \vert < \epsilon \} \\
            (V_\epsilon(0) - \{0\}) \cap A &= \{ x : 0 < x < \epsilon \} \neq \emptyset
        \end{flalign*}
        Hence, 0 is a limit point of $A$
    \end{itemize}
    \item Definition (Functional Limit) - Let $f: A \rightarrow \mathbb{R}$ and let $c$ be a limit point of $A$. We say $\lim_{x \rightarrow c} f(x) = L$ if
    \[
    (\forall \epsilon > 0)(\exists \delta > 0)(0 < \vert x-c \vert < \delta \rightarrow \vert f(x) - L \vert < \epsilon)
    \]
    \item Definition (Differentiability) - Let $g: A \rightarrow \mathbb{R}$. The derivative at $c \in A$ is
    \[
    g'(c) = \lim_{x \rightarrow c} \frac{g(x)-g(c)}{x-c}
    \]
    if this limit exists. If $g'(c)$ exists for all $c \in A$, we say $g$ is differentiable on $A$
    \item Theorem (Taylor's Theorem 1D) - Let $f: \mathbb{R} \rightarrow \mathbb{R}$ be $k$ times differentiable at $a$. Then
    \[
        f(x) = \sum_{n=0}^k \frac{(x-a)^n}{n!} f^{(n)}(a) + (x-a)^k h_k(x) = f(a) + (x-a)f'(a) + \frac{(x-a)^2}{2!} f''(a) + \dots + \frac{(x-a)^k}{k!}f^{(k)}(a) + (x-a)^k h_k(x)
    \]
    where $\lim_{x \rightarrow a} h_k(x) = 0$ (Peano form of the remainder)
    \item Example: Taylor Expansion of $f(x)=x^2$ around $x=1$
    \begin{itemize}
        \item $x^2 = 1 + 2(x-1) + (x-1)^2 + (x-1)^3 h_2(x), \qquad \lim_{x \rightarrow 1} h_2(x) = 0$
        \item Direct computation shows $h_2(x)=0$, so
        \[
        x^2 = 1 + 2(x-1) + (x-1)^2
        \]
        \item from a linear algebra perspective, $\{ 1,x-1,(x-1)^2\}$ is a basis for quadratic functions
    \end{itemize}
    \item Example: Taylor expansion of $\ln(1+x)$
    \begin{itemize}
        \item $f(x) = \ln(1+x)$
        \begin{flalign*}
            f'(x) &= \frac{1}{1+x} \\
            f''(x) &= -(1+x)^{-2} \\
            f^{(3)}(x) &= 2(1+x)^{-3} \\
            f^{(4)}(x) &= -6(1+x)^{-4} \\
            f^{(n)}(x) &= (-1)^{n+1} (n-1)!(1+x)^{-n} \\
        \end{flalign*}
        \item At $x=0$,
        \[
        f^{(n)}(0) = (-1)^{n+1}(n-1)!
        \]
        \item Therefore:
        \[
        \ln(1+x) = x - \frac{x^2}{2} + \frac{x^3}{3} - \cdots
        \]
    \end{itemize}
    \item For 1D, $y = f(x)$
    \begin{itemize}
        \item We know $f(a) = f'(a) a + c$
        \item Then $y= f'(a) x + (f(a) - af'(a)) = f(a) + (x-a)f'(a)$
        \item Expanding out, we get the Taylor Series expansion for 1D
    \end{itemize}
    \item Example: Taylor expansion of $f(x) = x(x+1)(x-1)=x(x^2-1)=x^3-x$ around $x=\frac{1}{\sqrt{3}}$
    \begin{itemize}
        \item $f'(x) = 3x^2 - 1 \Rightarrow f'(\pm \frac{1}{\sqrt{3}}) = 0$
        \item $f(a) + (x-a)\underset{0}{\underbrace{f'(a)}} + \frac{(x-a)^2}{2!} f''(a) = f(a) + \frac{(x-a)^2}{2!}f''(a) = \frac{1}{\sqrt{3}}(\frac{1}{3}-1) + \frac{(x-\frac{1}{\sqrt{3}})^2}{2} \frac{6}{\sqrt{3}}$
        \item Since $f''(a) = \frac{6}{\sqrt{3}} > 0$ for all $x$, $x=a$ is a local minimum
    \end{itemize}
    \item Definition (Optimum Values) - Let $I = \{ x \in \mathbb{R} : a \le x \le b \}$
    \begin{itemize}
        \item A function $f : I \rightarrow \mathbb{R}$ has:
        \begin{itemize}
            \item Global maximum at $c$ if $(\forall x \in I) f(x) \le f(c)$
            \item Global minimum at $c$ if $(\forall x \in I) f(x) \ge f(c)$
            \item Global optimum if it is either a global max or min
        \end{itemize}
        \item Local version: If $\exists \delta > 0$ s.t. in $(\forall x \in V_\delta(c))$
        \begin{itemize}
            \item Local maximum at $c$ if $f(x) \le f(c)$
            \item Local minimum at $c$ if $f(x) \ge f(c)$
        \end{itemize}
    \end{itemize}
    \item Using Taylor's Theorem to Classify Extrema
    \begin{itemize}
        \item Suppose $f$ is twice differentiable at $a$
        \[
        f(x) = f(a) + (x-a)f'(a) + \frac{(x-a)^2}{2} f''(a) + (x-a)^3 h_2(x), \qquad \lim_{x \rightarrow a} h_2(x) = 0
        \]
        \item If $f'(a) = 0$:
        \begin{itemize}
            \item If $f''(a) > 0$, it is a local minimum at $a$
            \item If $f''(a) < 0$, it is a local maximum at $a$
        \end{itemize}
    \end{itemize}
    \item We shall consider functions of the form $f: V \rightarrow W$ where $V$ and $W$ are vector spaces
    \item $f: \mathbb{R}^2 \rightarrow \mathbb{R}, \qquad f(x,y)=x^2+y^2 = \left[ \begin{array}{c}
        x \\
        y
    \end{array}\right] \cdot \left[ \begin{array}{c}
        x \\
        y
    \end{array}\right] = r^2$
    \begin{itemize}
        \item Gradient/slope depends on the direction of travel (directional derivative)
        \item Suppose we are at the point $(a,b)$. What is the slope in the $x$-direction?
        \begin{itemize}
            \item To answer this, consider moving away from $(a,b)$ in the $x$-direction from $(a,b)$ to $(a+h,b)$
            \item Height will change by $\frac{f(a+h,b)-f(a,b)}{h}$ per unit I move
            \item Define $\frac{\partial f}{\partial x}(a,b) = \lim_{h \rightarrow 0} \frac{f(a+h,b)-f(a,b)}{h}$ and $\frac{\partial f}{\partial y}(a,b) = \lim_{h \rightarrow 0} \frac{f(a,b+h)-f(a,b)}{h}$
        \end{itemize}
        \item $\frac{f(a+h,b)-f(a,b)}{h} = \frac{(a+h)^2 + b^2 - (a^2 + b^2)}{h} = \frac{(a+h)^2 - a^2}{h}  = \frac{2ah + h^2}{h} \underset{h \rightarrow 0}{\rightarrow} 2a$
        \begin{itemize}
            \item To work $\frac{\partial f(x,y)}{\partial x}$, hold $y$ constant and proceed on
        \end{itemize}
        \item $\nabla f(x,y) = \left[ \begin{array}{cc}
            \frac{\partial f(x,y)}{\partial x} & \frac{\partial f(x,y)}{\partial y}
        \end{array} \right] = \left[ \begin{array}{cc}
            2x & 2y
        \end{array} \right]$
        \begin{itemize}
            \item The gradient vector of $f$ at $(a,b)$: $\nabla f(a,b) = \left[ \begin{array}{cc}
                2a & 2b
            \end{array} \right]$
            \item Gradient is normal to the tangent of the contour at $(a,b)$
        \end{itemize}
    \end{itemize}
    
    \item Definition (Partial Derivative) - For a function $f: \mathbb{R}^n \rightarrow \mathbb{R}, x \mapsto f(x), x \in \mathbb{R}^n$ of $n$ variables $x_1,\dots,x_n$, we define the partial derivatives as:
    \begin{flalign*}
        \frac{\partial f}{\partial x_1} &= \lim_{h\rightarrow 0} \frac{f(x_1+h,x_2,\dots,x_n)-f(x)}{h}, \\
        \vdots& \\
        \frac{\partial f}{\partial x_n} &= \lim_{h\rightarrow 0} \frac{f(x_1,x_2,\dots,x_n+h)-f(x)}{h}, \\
    \end{flalign*}
    and collect them in the gradient of $f$:
    \[
    \nabla f = grad(f) = \frac{df}{dx} = \left[ \begin{array}{ccc}
        \frac{\partial f(x)}{\partial x_1} & \cdots & \frac{\partial f(x)}{\partial x_n}
    \end{array} \right] \in \mathbb{R}^{1 \times n}
    \]
    \item Definition (Jacobian) - the collection of all first-order partial derivatives of a vector-valued function $f: \mathbb{R}^n \rightarrow \mathbb{R}^m$ is called the Jacobian
    \begin{itemize}
        \item The Jacobian $J$ is an $m \times n$ matrix:
        \begin{flalign*}
            J &= \nabla_x f = \frac{df(x)}{dx} = \left[ \begin{array}{ccc}
                \frac{\partial f(x)}{\partial x_1} & \cdots & \frac{\partial f(x)}{\partial x_n}
            \end{array} \right] \\
            &= \left[ \begin{array}{ccc}
                \frac{\partial f_1(x)}{\partial x_1} & \cdots & \frac{\partial f_1(x)}{\partial x_n} \\
                \vdots && \vdots \\
                \frac{\partial f_m(x)}{\partial x_1} & \cdots & \frac{\partial f_m(x)}{\partial x_n}
            \end{array} \right], \qquad x = \left[ \begin{array}{c}
                x_1 \\ \vdots \\ x_n
            \end{array} \right], \qquad J_{ij} = \frac{\partial f_i}{\partial x_j}
        \end{flalign*}
    \end{itemize}
    \item Example Jacobian
    \begin{itemize}
        \item $f: \mathbb{R}^n \rightarrow \mathbb{R}^m$
        \item $f(x) = Ax$, where $A \in \mathbb{R}^{m \times n}, \quad x \in \mathbb{R}^n$
        \item What is the $i$th element of $f(x)$ (row multiplication of $A_{i, \cdot}$ with $x$)?
        \begin{itemize}
            \item $f_i(x) = \sum_{j=1}^n A_{ij}x_j = A_{i1}x_1 + A_{i2}x_2 + \cdots + A_{in}x_n$
        \end{itemize}
        \item $J_{ij}(x) = \frac{\partial f_i(x)}{\partial x_j} = A_ij \Rightarrow J = A$
    \end{itemize}
    \item Definition (Hessian Matrix) - consider a function $f: \mathbb{R}^n \rightarrow \mathbb{R}$. Provided all the second order partial derivatives exist, we define the matrix
    \[
    \left[ \begin{array}{cccc}
        \frac{\partial^2 f}{\partial x_1^2} & \frac{\partial^2 f}{\partial x_1\partial x_2} & \cdots & \frac{\partial^2 f}{\partial x_1\partial x_n} \\
        \frac{\partial^2 f}{\partial x_2\partial x_1} & \frac{\partial^2 f}{\partial x_2^2} & \cdots & \frac{\partial^2 f}{\partial x_2\partial x_n} \\
        \vdots & \vdots & \ddots & \vdots \\
        \frac{\partial^2 f}{\partial x_n\partial x_1} & \frac{\partial^2 f}{\partial x_n \partial x_2} & \cdots & \frac{\partial^2 f}{\partial x_n^2}
    \end{array} \right]
    \]
    The entry in the $i$th row and $j$th column is
    \[
    (H_f)_{ij} = \frac{\partial^2 f}{\partial x_i \partial x+j}
    \]
    \item Theorem (Multivariate Taylor Expansion) - Let $f: \mathbb{R}^n \rightarrow \mathbb{R}$ be twice continuously differentiable at the point $a \in \mathbb{R}^n$. Then
    \[
    f(x) = f(a) + \nabla f(a) (x-a) + \frac{1}{2!}(x-a)^T H_f(a) (x-a) + \frac{1}{3!} \sum_{i,j,k \in \mathbb{N}: i+j+k = 3} (x_i - a_i)(x_j - a_j)(x_k - a_k) h_{i,j,k}(x)
    \]
    where $\lim_{x \rightarrow a} h_{i,j,k}(x) = 0$
    \item Definition (Epsilon Neighborhood) - an epsilon neighborhood around $a \in V$, where $V$ is an inner product space (over a field $\mathbb{F}$), is defined by
    \[
    V_\epsilon(a) = \{ x \in V : \vert \vert x - a \vert \vert < \epsilon \}
    \]
    where $\vert \vert x \vert \vert \sqrt{\langle x, x \rangle}$
    \item Definition (Optimum Values in Several Variables) - Let
    \[
    I = \{ x \in \mathbb{R}^n : (\forall i \in \{1,\dots, n\}) \ a_i \le x_i \le b_i\}
    \]
    be an interval and $f: I \rightarrow \mathbb{R}$ be a function on $I$. Then $f$ has a
    \begin{itemize}
        \item Global maximum at $c$ if $(\forall x \in I) \ f(x) \le f(x)$
        \item Global minimum at $c$ if $(\forall x \in I) \ f(x) \ge f(c)$
        \item Global optimum at $c$ if either global max or min
    \end{itemize}
    \item An optimum can also be local. If $(\exists \epsilon > 0)(\forall x \in V_\epsilon(c))$:
    \begin{itemize}
        \item $f(x) \le f(c)$, then $f$ has a local maximum at $c$
        \item $f(x) \ge f(c)$, the $f$ has a local minimum at $c$
        \item We say that $f$ has a local optimum at $c$ if either local max or min
    \end{itemize}
\end{itemize}

\section{Lecture 14 (04/12/2025)}
\begin{itemize}
    \item $f: \mathbb{R}^2 \rightarrow \mathbb{R}$ be twice differentiable at $a \in \mathbb{R}^n$
    \begin{itemize}
        \item $f(x) = f(a) + \nabla f(a) (x-a) +\frac{1}{2} 
        (x-a)^T H_f(a) (x-a) + \dots$
    \end{itemize}
    \item $f(x,y) = x^2 + y^2$
    \begin{itemize}
        \item $a = \left[ \begin{array}{c}
            0 \\ 0
        \end{array} \right]$
        \item $\nabla f(x,y) = \left[ \begin{array}{cc}
            2x & 2y
        \end{array} \right]$
        \item $H_f(x) = \left[ \begin{array}{cc}
            2 & 0 \\
            0 & 2
        \end{array} \right] = 2I$
        \item $f(x) = 0^2 + 0^2 + \left[ \begin{array}{cc}
            0 & 0
        \end{array}\right] \left[ \begin{array}{c}
            x \\ y
        \end{array} \right] + \left[ \begin{array}{cc}
            x & y
        \end{array}\right] I \left[ \begin{array}{c}
            x \\ y
        \end{array} \right] = x^2 + y^2$
    \end{itemize}
    \item Because the function doesn't have higher-order terms, the Taylor series expansion will naturally truncate
    \item $f(x,y) = x^2 - y^2$
    \begin{itemize}
        \item $a = \left[ \begin{array}{c}
            0 \\ 0
        \end{array} \right]$
        \item $f(x,y) = \left[ \begin{array}{cc}
            x & y
        \end{array}\right] \left[ \begin{array}{cc}
            1 & 0 \\
            0 & -1
        \end{array} \right] \left[ \begin{array}{c}
            x \\ y
        \end{array} \right]$
        \item $\nabla f(x,y) = \left[ \begin{array}{cc}
            2x & -2y
        \end{array} \right]$
        \item $H_f(x) = \left[ \begin{array}{cc}
            2 & 0 \\
            0 & -2
        \end{array} \right]$
        \item $f(x) = 0^2 + 0^2 + \left[ \begin{array}{cc}
            0 & 0
        \end{array}\right] \left[ \begin{array}{c}
            x \\ y
        \end{array} \right] + \left[ \begin{array}{cc}
            x & y
        \end{array}\right] \left[ \begin{array}{cc}
            1 & 0 \\
            0 & -1
        \end{array} \right] \left[ \begin{array}{c}
            x \\ y
        \end{array} \right] = x^2 - y^2$
        \item Change in height if I move to $\Delta x$: $f(\Delta x, 0) - f(0,0) = \left[ \begin{array}{cc}
            \Delta x & 0
        \end{array}\right] \left[ \begin{array}{cc}
            1 & 0 \\
            0 & -1
        \end{array} \right] \left[ \begin{array}{c}
            \Delta x \\ 0
        \end{array} \right] = (\Delta x)^2 > 0$
        \item Change in horizontal distance if I move to $\Delta y$: $f(0, \Delta y) - f(0,0) = \left[ \begin{array}{cc}
            0 & \Delta y
        \end{array}\right] \left[ \begin{array}{cc}
            1 & 0 \\
            0 & -1
        \end{array} \right] \left[ \begin{array}{c}
            0 \\ \Delta y
        \end{array} \right] = -(\Delta y)^2 < 0$
    \end{itemize}
    \item $f(x,y) = x^3 - y^3 + x^2y + y^2$
    \begin{itemize}
        \item $\nabla f(x,y) = \left[ \begin{array}{cc}
            3x^2 +2xy & -3y^2 + x^2 + 2y 
        \end{array} \right]$
        \item $H_f(x) = \left[ \begin{array}{cc}
            6x + 2y & 2x \\
            2x & -6y^2 +2
        \end{array} \right]$
        \item Symmetric $\because \frac{\partial^2 f}{\partial x \partial y} = \frac{\partial^2 f}{\partial y \partial x}$
    \end{itemize}
    \item Theorem (Real Symmetric Matrices) - a matrix in $\mathbb{R}^{n \times n}$ has all real eigenvalues and $N$ orthonormal real eigenvectors if and only if it is symmetric
    \item We can locate the local optima of $f: \mathbb{R}^n \rightarrow \mathbb{R}$ by looking for points where $\nabla f = 0$
    \begin{itemize}
        \item Suppose at $a \in \mathbb{R}^n$, we have $\nabla f(a) = 0$
        \item To classify a local optimum, e look at the second order term in the Taylor expansion:
        \[
        f(x) = f(a) + \frac{1}{2!} (x-a)^T H_f(a) (x-a) + \frac{1}{3!} \sum_{i,j,k \in \mathbb{N}: i+j+k = 3} (x_i - a_i)(x_j-a_j)(x_k-a_k)h_{i,j,k}(x)
        \]
        \item The term 
        \[
        (x-a)^T H_f(a)(x-a)
        \]
        is a quadratic form and tells us about the local behavior of $f$ around $a$
    \end{itemize}
    \item We note that $H_f(a) \in \mathbb{R}^{n \times n}$ is a real symmetric matrix and so its $n$ eigenvectors $\{ e_1, \dots, e_n\}$ form an orthonormal basis for $\mathbb{R}^n$ (we can always normalize so that $(\forall i \in \{1, \dots, n\}) \lvert \lvert e_i \rvert \rvert = 1$). Let the corresponding eigenvalues be $\{ \lambda_1, \dots, \lambda_n\}$. We can represetn this with
    \[
    H_f(a) e = \lambda e
    \]
    \begin{itemize}
        \item Hence, there exists unique $b_1, \dots, b_n \in \mathbb{R}$ such that
        \begin{flalign*}
            x-a &=\sum_{i=1}^n b_ie_i \\
            (x-a)^T H_f(a) (x-a) &= (\sum_{i=1}^n b_ie_i)^T H_f(a) (\sum_{j=1}^n b_je_j) \\
            &= (\sum_{i=1}^n b_ie_i)^T \sum_{j=1}^n b_j H_f(a) e_j \\
            &= (\sum_{i=1}^n b_ie_i)^T \sum_{j=1}^n b_j \lambda_j e_j \\
            &= \sum_{i=1}^n b_i \sum_{j=1}^n \lambda_j b_j e_i^T e_j \\
            &= (b_1e_1 + \dots + b_ne_n)^\top \cdot (b_1\lambda_1 e_1 + \dots + b_n \lambda_n e_n) \\
            &\text{Since basis is orthonormal, } e_i^Te_j = 0 \text{ if } i \neq j, \qquad e_i^\top e_i = 1 \\
            &= b_1^2 \lambda_1 + b_2^2 \lambda_2  + \dots + b_n^2 \lambda_n \\
            &= \sum_{i=1}^n b_i^2 \lambda_i
        \end{flalign*}
        \item It follows that if the eigenvalues of the Hessian are all positive (with at least one strictly positive), then $(\forall (x - a) \in \mathbb{R}^n)$
        \[
        (x-a)^T H_f(a) (x-a) \ge 0
        \]
        and for at least one choice $> 0$. Therefore, $a$ is a local minimum of $f$
    \end{itemize}
    \item Classification via eigenvalues of the Hessian
    \begin{itemize}
        \item Local minimum: if all eigenvalues of the Hessian are positive (with at least one strictly positive), then $(\forall x - a \in \mathbb{R}^n)$
        \[
        (x-a)^T H_f(a) (x-a) \ge 0
        \]
        with strict inequality for some direction. Then $a$ is a local minimum
        \item Local maximum: if all eigenvalues of the Hessian are negative (with at least one strictly negative), then $(\forall x - a \in \mathbb{R}^n)$
        \[
        (x-a)^T H_f(a) (x-a) \le 0
        \]
        with strict inequality for some direction. Then $a$ is a local maximum
        \item Saddle point: if the eigenvalues of the Hessian differ in sign, then $f$ increases in the direction of an eigenvector corresponding to a positive eigenvalue and decreases in the direction of an eigenvector corresponding to a negative eigenvalue. We say $a$ is a saddle point
        \item If the eigenvalues are all zero, then we need to look at higher-order derivatives
    \end{itemize}
    \item Minimize $f(x,y) = x^2 + y^2$ subject to the constraint $x+y=1$
    \begin{itemize}
        \item By geometry, we can see $x^2+y^2$ will be minimized when the radius $r$ of the circle $x^2+y^2 = r^2$ is chosen so that the circle is tangent to the line $x+y=1$
        \item Further, we can see that the minimum value of $x^2+y^2$ is
        \[
        \left(\frac{1}{2}\right)^2 + \left( \frac{1}{2} \right)^2 =\frac{1}{2}
        \]
        which is achieved when $(x,y) = (\frac{1}{2},\frac{1}{2})$
        \item $\nabla f(x,y) = (2x,2y)$
        \item $\nabla g(x,y) = (1,1)$
        \item When $f(x,y)$ is minimized subject to $g(x,y) = 1$, $\nabla f(x,y)$ is a multiple of $\nabla g(x,y)$, i.e.
        \[
        \nabla f(x,y) = \lambda \nabla g(x,y)
        \]
        where the constant of proportionality $\lambda$ is called a Lagrange multiplier
        \item From $(2x,2y)=\lambda(1,1)$, we have $x=y=\frac{\lambda}{2}$
        \item $g(x,y) = 2\frac{\lambda}{2} = \lambda = 1$
    \end{itemize}
    \item Value Function - $V(a) = \underset{g(x,y)=a}{\min} f(x,y)$
    \begin{itemize}
        \item $\frac{\partial}{\partial a} V(a) = a = \lambda$
        \item The Lagrange multiplier $\lambda$ gives the marginal change in the value function for a relaxation in the constraint
        \item Economists call this the shadow price of the constraint
    \end{itemize}
    \item Lagrangian formulation
    \begin{itemize}
        \item We can obtain the first order condition $\nabla f(x,y) = \lambda \nabla g(x,y)$ and the constraint $g(x,y) = a$ from the following optimization problem $\min L(x,y ; \lambda)$ where the Lagrangian $L(x,y; \lambda) = f(x,y) + \lambda(a-g(x,y))$
    \end{itemize}
    \item Constrained Optimization Problem
    \begin{itemize}
        \item Minimize $f(x,y)$ such that $g(x,y)=a$
        \item When the gradient vector for $f$ is parallel to the gradient vector for $g$ ($\nabla f = \lambda \nabla g$)
        \item Equivalent unconstrained problem: $L(x,y; \lambda) = f(x,y) + \lambda(a-g(x,y))$ where $\nabla L = \nabla f - \lambda \nabla g = 0$
        \item Now note that
        \begin{flalign*}
            \frac{\partial L}{\partial x} &= \frac{\partial f(x,y)}{\partial x} - \lambda \frac{\partial g(x,y)}{\partial x} \\
            \frac{\partial L}{\partial y} &= \frac{\partial f(x,y)}{\partial y} - \lambda \frac{\partial g(x,y)}{\partial y} \\
            \frac{\partial L}{\partial \lambda} &= a - g(x,y) \\
        \end{flalign*}
        \item Therefore the first order conditions 
        \[
        \frac{\partial L}{\partial x} = 0,
        \]
        \[
        \frac{\partial L}{\partial y} = 0,
        \]
        \[
        \frac{\partial L}{\partial \lambda} = 0
        \]
        imply that
        \[
        \nabla f = \lambda \nabla g, \quad g(x,y) = a
        \]
    \end{itemize}
\end{itemize}

\section{Tutorial 9 (04/12/2025)}
\begin{itemize}
    \item Definition (Orthogonal Matrix) - a square matrix $A$ whose transpose is also its inverse:
    \[
    A^\top A = A A^\top = A^{-1}A = A^{-1}A = I
    \]
    \item Exercise 6.14. (a) Consider the vector space $\mathbb{R}^N$ with the following inner product:
    \[
    \forall x,y \in \mathbb{R}^N, \langle x,y \rangle = \sum_{n=1}^N x_ny_n
    \]
    Let $O$ be an $N \times N$ orthogonal matrix. That is, $O^\top O = I_N$. Show that
    \[
    \langle Ox,Oy \rangle = \langle x,y \rangle.
    \]
    How would you interpret this result geometrically?
    \begin{itemize}
        \item We note that
        \[
        \langle Ox,Oy \rangle = (Ox)^\top(Oy) = x^\top O^\top O y = x^\top (O^\top O) y = x^\top I_N y = x^\top y = \langle x,y \rangle
        \]
        \item Geometrically, an orthogonal matrix preserves lengths and angles.
    \end{itemize}
    \item Exercise 6.14. (b) The orthogonal matrix $O$ from above can be written as
    \[
    O = [v_1,v_2,\dots,v_N],
    \]
    where $v_1,\dots,v_N$ are column vectors in $\mathbb{R}^N$. What does the condition $O^\top O = I_N$ tell us about $\langle v_i,v_j\rangle$.
    \begin{itemize}
        \item Notice that the $ij$-entry of $AB$ is computed by taking the inner product of the $i$th row of $A$ with the $j$th column of $B$
        \item Moreover, $v_i^\top$ is the $i$th row of $O^\top$ and $v_j$ is the $j$th column of $O$.
        \item So indeed, $\langle v_i,v_j\rangle$ is nothing but the $ij$-entry of the identity matrix $I_N$.
        \item Therefore,
        \[
        \langle v_i,v_j \rangle = \left\{ \begin{array}{lr}
            1 & i=j \\
            0 & i \neq j
        \end{array} \right.
        \]
        \item In other words, the columns of an orthogonal matrix are orthonormal.
    \end{itemize}
    \item 2023 Sample Final Exam Question 2 (a) Find the orthogonal projection of the vector $b = (1,2,6)^\top$ onto the plane $x+y+z = 0$ in $\mathbb{R}^3$, where the inner product is the standard scalar product.
    \begin{itemize}
        \item The easiest way to do this is to project $b$ onto the one-dimensional orthogonal complement of the plane first, and then subtract that off.
        \item Note that the equation of the plane is $r \cdot n = 0$, where $r = (x,y,z)^\top$ and $n = (1,1,1)^\top$ is the normal vector. A unit vector in the normal direction is therefore $\hat{n} = \frac{1}{\sqrt{3}} (1,1,1)^\top$, and projecting $b$ onto this gives
        \[
        b_\perp = (b \cdot \hat{n})\hat{n} = \frac{9}{\sqrt{3}} \times \frac{1}{\sqrt{3}} (1,1,1)^\top = (3,3,3)^\top
        \]
        \item Thus, the orthogonal projection onto the plane is the rest of $b$, which is
        \[
        b_{\vert \vert} = b - (b \cdot \hat{n})\hat{n} = (1,2,6)^\top - (3,3,3)^\top = (-2,-1,3)^\top
        \]
    \end{itemize}
    \item 2023 Sample Final Exam Question 2 (b) Let $V$ be the vector space of real polynomials of degree $\le 2$ and define the inner product as follows
    \[
    \langle \cdot,\cdot \rangle : V \times V \longrightarrow \mathbb{R}
    \]
    \[
    \langle f,g \rangle \longmapsto \frac{1}{3} \sum_{k \in \{-1,0,1\}} f(k) \cdot g(k)
    \]
    Let $U$ be the subspace spanned by $1$ and $x$.
    \begin{itemize}
        \item (i) Show that $1$ and $x$ are orthogonal and use this to find an orthonormal basis for $U$.
        \begin{itemize}
            \item $\langle 1,x \rangle = \frac{1}{3}(1(-1)+1(0)+1(1))=0$ and so $1$ and $x$ are orthogonal. To get an orthonormal basis, we just need to normalise them. We have:
            \[
            \vert \vert 1 \vert \vert^2 = \langle 1,1 \rangle = \frac{1}{3}(1+1+1) = 1
            \]
            \[
            \vert \vert x \vert \vert^2 = \langle x,x \rangle = \frac{1}{3} ((-1)(-1)+(0)(0)+(1)(1)) = \frac{2}{3}
            \]
            so $\vert \vert 1 \vert \vert=1$ and $\vert \vert x \vert \vert = \sqrt{\frac{2}{3}}$. Thus an orthonormal basis for $U$ is
            \[
            \left\{ 1, \sqrt{\frac{3}{2}}x \right\}.
            \]
            \item Note that this set is clearly linearly independent since $1$ and $x$ are not proportional to each other, and it clearly spans $U$ since $U$ contains all polynomials fo the form $a+bx$. Therefore, this set is indeed a basis.
        \end{itemize}
        \item (ii) Find a set which serves as an orthonormal basis for $V$ and then prove it is indeed a basis of $V$.
        \begin{itemize}
            \item $f \in V$ takes the form $f(x) = ax^2+bx+c$. We want to find $f$ such that it is orthogonal to both $1$ and $x$. Note that
            \[
            \langle f,1 \rangle = \frac{1}{3}[(a-b+c) + c + (a+b+c)] = \frac{1}{3}(2a+3c)
            \]
            \[
            \langle f,x \rangle = \frac{1}{3} [-(a-b+c)+0\cdot c + (a+b+c)] = \frac{2b}{3}
            \]
            \item For both of the above to be zero, we need $b=0$ and $a=\frac{-3c}{2}$.
            \item Thus, vectors orthogonal to both $1$ and $x$ take the form
            \[
            f(x) = -c\left(\frac{3}{2}x^2 - 1\right).
            \]
            \item We need this to be normalised, which determines $c$. Note that
            \[
            \vert \vert f \vert \vert^2 = \langle f,f \rangle = \frac{c^2}{3} \left( \frac{1}{2}\cdot\frac{1}{2} + (-1)(-1) + \frac{1}{2}\cdot\frac{1}{2} \right) = \frac{c^2}{2}
            \]
            So
            \[
            \vert \vert f \vert \vert = 1 \Longrightarrow c^2 = 2 \Longleftrightarrow c = \pm\sqrt{2}
            \]
            \item Picking $c=-\sqrt{2}$ (either choice works), we obtain the orthonormal set
            \[
            \left\{ 1, \sqrt{\frac{3}{2}}x, \sqrt{2}\left(\frac{3}{2}x^2-1\right) \right\}.
            \]
            \item We just need to check that this is actually a basis. First suppose that
            \[
            \lambda_1(1) + \lambda_2\sqrt{\frac{3}{2}} x + \sqrt{2}\lambda_3 \left(\frac{3}{2}x^2 - 1\right) = 0.
            \]
            Then, comparing $x$ terms, $\lambda_2 = 0$, comparing $x^2$ terms, $\lambda_3 = 0$, and comparing constant terms, $\lambda_1 = 0$. Thus, our set is linearly independent
            \item For spanning, let $g(x) = kx^2 + lx + m$ be arbitrary. We want to find $\lambda_i$s such that
            \[
            \lambda_1(1) + \lambda_2\sqrt{\frac{3}{2}} x + \sqrt{2}\lambda_3 \left(\frac{3}{2}x^2 - 1\right) = kx^2+lx+m
            \]
            Clearly, this can always be solved, with
            \[
            \lambda_3 =\frac{\sqrt{2}}{3} k, \quad \lambda_2 = \sqrt{\frac{2}{3}} l, \lambda_1 = m + \frac{2}{3}k
            \]
            and so our set spans $V$. Thus, it is an orthonormal basis of $V$.
        \end{itemize}
    \end{itemize}
    \item Bonus Question: Let $(V,\langle\cdot,\cdot\rangle)$ be an inner product space. Is it true that if the vectors $\{v_1,\dots,v_n\} \subset V$ are all non-zero and pairwise orthogonal, then those vectors are also linearly independent?
    \begin{itemize}
        \item Yes. Suppose that
        \[
        \lambda_1 v_1 + \lambda_2 v_2 + \cdots + \lambda_n v_n =0
        \]
        \item Now note that $\forall i$,
        \begin{flalign*}
            &\langle \lambda_1 v_1 + \lambda_2 v_2 + \cdots + \lambda_n v_n, v_i \rangle = \langle 0,v_i \rangle = 0 \\
            \Rightarrow &\lambda_1 \langle v_1,v_i \rangle + \lambda_2 \langle v_2,v_i \rangle + \cdots + \lambda_n \langle v_n,v_i \rangle = 0 \\
            \Rightarrow &\lambda_i \vert \vert v_i \vert \vert^2 = 0 \\
            \Rightarrow &\lambda_i = 0
        \end{flalign*}
        since $\langle v_i,v_j\rangle \neq 0 \Longleftrightarrow i=j$. Hence the vectors are linearly independent.
    \end{itemize}
    \item 2023 Sample Final Exam Question 2 (c) Consider a finite dimensional vector space $V$ over a field $\mathbb{F}$ with subspaces $U$ and $W$.
    \begin{itemize}
        \item (i) What is meant by the sum $U+W$ of $U$ and $W$?
        \[
        U+W = \{ u+w : u \in U \text{ and } w \in W \}.
        \]
        \item (ii) Show that $U+W = span(U \cup W)$. When is $U+W = U \cup W$? Justify your answer.
        \begin{itemize}
            \item First suppose that $v \in U+W$. Then $\exists u \in U$ and $\exists w \in W$ such that $v=u+w$. Since $u,w \in U \cup W$, we find that $v \in span(U \cup W)$ and so conclude that
            \[
            U+W \subseteq span(U \cup W).
            \]
            \item Now suppose that $v \in span(U \cup W)$. Then $\exists \lambda_1,\dots,\lambda_n \in \mathbb{F}$, $\exists \mu_1,\dots,\mu_m \in \mathbb{F}$ and $\exists u_1,\dots,u_n \in U$, $\exists w_1,\dots,w_m \in W$ such that
            \[
            v = \sum_{i=1}^n \lambda_i u_i + \sum_{i=1}^m \mu_i w_i
            \]
            \item Since $U$ and $W$ are subspaces, they are closed under addition and scalar multiplication. So,
            \[
            u = \sum_{i=1}^n \lambda_i u_i \in U \text{ and } w = \sum_{i=1}^m \mu_i w_i \in W.
            \]
            \item Thus, $v \in span(U \cup W)$ can be written as $v = u+w$, where $u \in U$ and $w \in W$. In other words, $v \in U+W$ and so $span(U \cup W) \subseteq U+W$.
            \item Note that $U+W = U \cup W \Longleftrightarrow span(U \cup W) = U \cup W$.
            \item We will now show that this can happen if any only if $U \subseteq W$ or $W \subseteq U$.
            \item If $U \subseteq W$, then $U \cup W = W$. Furthermore, since $W$ is a vector space, $span(W)=W$. Therefore,
            \[
            U \subseteq W \Rightarrow span(U \cup W) = U \cup W \Rightarrow U+W = U \cup W.
            \]
            \item Similarly, $W \subseteq U$ also implies that $U+W = U \cup W$.
            \item For the other direction, we consider the contrapositive. Suppose that $U \nsubseteq W$ and $W \nsubseteq U$. Then $\exists u \in U$ and $\exists w \in W$ such that $u \notin W$ and $w \notin U$.
            \item Then, since $U$ and $W$ are closed under addition, it must be that $u+w \notin U$ and $u+w \notin W$. Hence $u+w \notin U \cup W$.
            \item However, $u+w \in U + W$ and so under these assuptions, $U+W \neq span(U \cup W)$.
            \item Thus, conversely, $U+W = U \cup W$ implies that $U \subseteq W$ or $W \subseteq U$.
            \item This completes the proof of both directions.
        \end{itemize}
    \end{itemize}
\end{itemize}

\section{Lecture 15 (09/12/2025)}
\begin{itemize}
    \item Consumer problem with two goods
    \begin{itemize}
        \item A consumer has the choice of purchasing two goods at the same price of 1 USD
        \[
        u(C_1,C_2) = \ln(C_1)+\ln(C_2)
        \]
        \item Budget is $A$ USD
        \item Constrained optimization problem $\underset{C_1,C_2}\max u(C_1,C_2)$ such that $g(C_1,C_2)=A$ where $g(C_1,C_2) = C_1+C_2$
        \begin{flalign*}
            \ln(C_1)+\ln(C_2) &= k \\
            \ln(C_1 C_2) &= k \\
            C_1C_2 &= e^k \\
            C_2 &= \frac{e^k}{C_1}
        \end{flalign*}
        \item Drawing contour lines of $u(C_1,C_2)$ and the constraint $g(C_1,C_2)$, we can see that the utility will be maximized when the gradient of $u$ is parallel to the gradient of $g$
        \begin{itemize}
            \item $\nabla u = \lambda \nabla g$
        \end{itemize}
        \item Lagrangian Method: $L = u(C_1,C_2) - \lambda(g(C_1,C_2) - A)$
        \begin{itemize}
            \item $\nabla L = \nabla u - \lambda \nabla g = 0 \Rightarrow \nabla u = \lambda \nabla g$
            \begin{flalign*}
                \nabla u &= \left[ \begin{array}{cc}
                C_1^{-1} & C_2^{-1}
                \end{array} \right] \\
                \lambda \nabla g &= \lambda \left[ \begin{array}{cc}
                    1 & 1
                \end{array} \right] \\
                \left[ \begin{array}{cc}
                C_1^{-1} & C_2^{-1}
                \end{array} \right] &= \left[ \begin{array}{cc}
                    \lambda & \lambda
                \end{array} \right] \\
                \Rightarrow C_1^{-1} &= \lambda = C_2^{-1} \\
                C_1 + C_2 &= 2\lambda^{-1} \\
                \text{Constraint: } C_1 + C_2 &= A \\
                2\lambda^{-1} &= A \\
                \Rightarrow \lambda &= \frac{2}{A} \\
                \Rightarrow C_1 &= C_2 = \frac{A}{2}
            \end{flalign*}
        \end{itemize}
    \end{itemize}
    \item The value function is given by
    \[
    V(A) = u(C_1,C_2) = 2 \ln(\frac{A}{2}) = 2 [\ln (A) - \ln(2)]
    \]
    \item Hence, we see that
    \[
    V'(A) = \frac{d}{dA} 2 [\ln (A) - \ln(2)] = \frac{2}{A} =\lambda
    \]
    \item The derivative of the value function is always equal to the Lagrange multiplier: $V'(A) = \lambda$
    \begin{itemize}
        \item In economics, $\lambda$ is the marginal value function (shadow price): $\Delta V \approx \lambda \Delta A$
    \end{itemize}
    \item General budget with prices $P_1,P_2$
    \begin{itemize}
        \item $\max_{C_1,C_2} u(C_1,C_2) = \max_{C_1,C_2} [\ln(C_1) + \ln(C_2)]$ such that $g(C_1,C_2)=A$ where $g(C_1,C_2) = P_1C_1+P_2C_2$
        \item We optimize the Lagrangian $L = u(C_1,C_2) + \lambda(A - g(C_1,C_2))$ and so
        \begin{flalign*}
            \left[ \begin{array}{cc}
                C_1^{-1} - \lambda P_1 & C_2^{-1} - \lambda P_2
                \end{array} \right] &= \left[ \begin{array}{cc}
                    0 & 0
                \end{array} \right] \\
            \nabla u(C_1,C_2) &= \lambda \left[ \begin{array}{cc}
            P_1 & P_2 \end{array} \right]  \\
            \Rightarrow (\forall i \in \{1,2\}) \ C_i &= \frac{1}{\lambda P_i}
        \end{flalign*}
        \item $\left[ \begin{array}{cc}
        P_1 & P_2 \end{array} \right]$ is the price vector
        \begin{flalign*}
            P_1C_1 = P_2C_2 &= \lambda^{-1} \\
            A &= 2\lambda^{-1} \\
            \Rightarrow \lambda &= \frac{2}{A} \\
            \Rightarrow (\forall i \in \{1,2\}) \ C_i &= \frac{1}{2}\frac{A}{P_i} \\
            V(A) &= \ln(\frac{1}{2}\frac{A}{P_1}) + \ln(\frac{1}{2} \frac{A}{P_2}) \\
            V'(A) &= \frac{2}{A} = \lambda
        \end{flalign*}
    \end{itemize}
    \item Optimization problem finding min, max, saddle point with eigenvalue of Hessian
    \item Linear regression problem
    \item Monotone Convergence Theorem and Bolzano-Weierstrauss Theorem
    \item Real analysis question on stuff after midterm and another question on anything
    \item Linear algebra question (vector spaces of polynomials)
    \item Constrained optimization question
\subsection{Slide 7 (Constrained Optimization with Cobb–Douglas Production)}
    \item A firm produces output using two inputs, $A$ and $B$. Production is given by
    \[
    h(a,b) = a^\alpha b^\beta, \qquad \alpha > 0, \beta > 0, \alpha + \beta = 1
    \]
    \item Inputs have prices $p_a$ and $p_b$, and the budget constraint is $I = p_a a + p_b b$
    \item Lagrangian Setup
    \begin{itemize}
        \item We maximize $h(a,n) = a^\alpha b^\beta$ subject to $p_aa + p_bb = I$
        \begin{flalign*}
            \mathcal{L}(a,b,\lambda) &= a^\alpha b^\beta - \lambda(p_a a + p_b b - I) \\
            \frac{\partial \mathcal{L}}{\partial a} &= \alpha a^{\alpha - 1} b^\beta - \lambda p_a = \frac{\alpha h}{a} - \lambda p_a = 0 \\
            \frac{\partial \mathcal{L}}{\partial b} &= \beta a^{\alpha} b^{\beta-1} - \lambda p_b = \frac{\beta h}{b} - \lambda p_a = 0 \\
            \frac{\partial \mathcal{L}}{\partial \lambda} &= I-p_aa-p_bb = 0 \\
            \nabla \mathcal{L} &= \left[ \begin{array}{ccc}
                \alpha a^{\alpha - 1} b^\beta - \lambda p_a & \beta a^{\alpha} b^{\beta-1} - \lambda p_b & I-p_aa-p_bb
            \end{array} \right] \\
            \frac{\alpha h}{a} &= \lambda p_a \\
            \frac{\beta h}{b} &= \lambda p_b \\
            \Rightarrow \frac{\frac{\alpha h}{a}}{\frac{\beta h}{b}} &= \frac{\lambda p_a}{\lambda p_b} \\
            \Rightarrow \frac{\alpha}{\beta} \frac{b}{a} &= \frac{p_a}{p_b} \\
            \Rightarrow b &= \frac{\beta p_a}{\alpha p_b} a \\
            p_a a + p_b \left( \frac{\beta p_a}{\alpha p_b} a \right) &= I \\
            p_a a \left( 1 + \frac{\beta}{\alpha} \right) &= I \\
            p_a a \left(\frac{\overset{1}{\overbrace{\alpha + \beta}}}{\alpha} \right) &= I \\
            \Rightarrow a &= \frac{\alpha I}{p_a}, b = \frac{\beta I}{p_b} \\
            (a^*,b^*) &= \left( \frac{\alpha I}{p_a}, \frac{\beta I}{p_b} \right)
        \end{flalign*}
    \end{itemize}
    \item Show that
    \[
    \frac{p_a a^*}{p_b b^*} = \frac{\alpha}{\beta}.
    \]
    \begin{itemize}
        \item Using the optimal quantities:
        \[
        (a^*,b^*) = \left( \frac{\alpha I}{p_a}, \frac{\beta I}{p_b} \right)
        \]
        we compute the expenditure ratio:
        \[
        \frac{p_aa^*}{p_bb^*} = \frac{p_a\frac{\alpha I}{p_a}}{p_b \frac{\beta I}{p_b}} = \frac{\alpha}{\beta}
        \]
    \end{itemize}
    \item Suppose the price of good $B$ increases to $kp_b$ with $k>1$. How does the optimal quantity $b$ change?
    \begin{itemize}
        \item Original optimal quantity: $b_1 = \frac{\beta I}{p_b}$
        \item New price: $kp_b$
        \item New optimal quantity: $b_2 = \frac{\beta I}{kp_b} = \frac{1}{k} \frac{\beta I}{p_b} = \frac{b_1}{k}$
        \item $b^*$ falls by a factor of $\frac{1}{k}$
        \item Thus quantity demanded decreases proportionally with the price increase
    \end{itemize}
\subsection{Slide 8 (Linear Algebra: Bases and Linear Maps)}
    \item Question 1: Let $V$ be a finite-dimensional vector space over $\mathbb{R}$ and let
    \[
    \mathcal{S} = \{v_1, \dots, v_n\} \subseteq V
    \]
    be a finite non-empty subset.
    \begin{itemize}
        \item (a) What does it mean for $\mathcal{S}$ to be a linearly independent set?
        \begin{itemize}
            \item $(\forall v_i \in \mathcal{S}) (\exists \lambda_i \in \mathbb{R}) \ \sum_{i=1}^n \lambda_iv_i = \bm{0} \Longleftrightarrow (\forall i \in \{1,\dots,n\}) \lambda_i = 0$
        \end{itemize}
        \item (b) What does it mean for $\mathcal{S}$ to be a spanning set for $V$?
        \begin{itemize}
            \item $(v_i : v_i \in \mathcal{S})(\forall v \in V)(\exists \lambda_i \in \mathbb{R}) \ \sum_{i=1}^n \lambda_iv_i = v$
        \end{itemize}
        \item (c) What does it mean for $\mathcal{S}$ to be a basis for $V$?
        \begin{itemize}
            \item $\mathcal{S}$ is a spanning set and a linearly independent set
            \item Equivalently, $\mathcal{S}$ is a minimal spanning set of $V$, or a maximal linearly independent subset of $V$
        \end{itemize}
    \end{itemize}
    \item Question 2: Let
    \[
    \bm{a} = \left[ \begin{array}{c}
        1 \\ 2 \\ 0
    \end{array} \right] \in \mathbb{R}^3
    \]
    and consider the map
    \[
    \phi: \mathbb{R}^3 \rightarrow \mathbb{R}^3, \qquad \phi(\bm{v}) = 5\bm{v} - (\bm{a} \cdot \bm{v})\bm{a}
    \]
    where $\cdot$ is the standard dot product
    \begin{itemize}
        \item (a) Show that $\phi$ is a linear map
        \begin{itemize}
            \item Additivity: For all $u,v \in \mathbb{R}^3$:
            \begin{flalign*}
                \phi(u+v) &= 5(u+v) - (a \cdot (u+v))a \\
                &= 5u+5v-[a\cdot u + a \cdot v]a \\
                &= (5u - (a \cdot u)a) + (5v - (a \cdot v)a) = \phi(u) + \phi(v)
            \end{flalign*}
            \item Homogeneity: For all $\lambda \in \mathbb{R}, v \in \mathbb{R}^3$:
            \begin{flalign*}
                \phi(\lambda v) &= 5(\lambda v) - (a \cdot (\lambda v))a \\
                &= 5\lambda v - \lambda (a\cdot v)a \\
                &= \lambda [5v - (a \cdot v)a] = \lambda \phi(v)
            \end{flalign*}
            \item Therefore, $\phi$ is linear
        \end{itemize}
        \item (b) Find the matrix of $\phi$ with respect to the standard basis of $\mathbb{R}^3$
        \begin{itemize}
            \item Let $v = (x,y,z)^T$.
            \begin{flalign*}
                a \cdot v &= \left[ \begin{array}{c}
                1 \\ 2 \\ 0 
                \end{array} \right]^T \left[ \begin{array}{ccc}
                       x \\ y \\ z
                    \end{array}\right]  = 2y \\
                (a \cdot v) a &= (x+2y)\left[ \begin{array}{c}
                    1 \\ 2 \\ 0 
                \end{array} \right] =  \left[ \begin{array}{c}
                    x+2y \\ 2x+4y \\ 0 
                \end{array} \right] \\
                \phi(v) &= 5 \left[ \begin{array}{c}
                    x \\ y \\ z 
                \end{array} \right] - \left[ \begin{array}{c}
                    x+2y \\ 2x+4y \\ 0 
                \end{array} \right] = \left[ \begin{array}{c}
                    4x-2y \\ -2x+y \\ 5z 
                \end{array} \right]
            \end{flalign*}
        \end{itemize}
        \item (c) Hence or otherwise, derive bases for $\text{ker}(\phi)$ and $\text{im}(\phi)$
        \begin{itemize}
            \item To find $\text{ker}(\phi)$, we solve $A \bm{v} = \bm{0}$ for $\bm{v} = (x,y,z)^T$
            \begin{flalign*}
                ker(\phi) &= \{ v \in \mathbb{R}^3: \phi(v) = \bm{0} \in \mathbb{R}^3 \} \\
                Av &= 0
            \end{flalign*}
            \[
            \left\{ \begin{array}{r}
                4x-2y = 0 \\
                -2x+y = 0 \\
                5z =0
            \end{array} \right.
            \]
            \begin{itemize}
                \item From the third equation: $z=0$
                \item The first two equations are equivalent and given $y=2x$
                \item Thus $v = \lambda(1,2,0)^T$ for any $x \in \mathbb{R}$, so:
                \[
                ker(\phi) = span\left\{ \left[ \begin{array}{c}
                    1 \\ 2 \\ 0
                \end{array} \right] \right\} = span\{a\}
                \]
            \end{itemize}
            \item $im(\phi) = \{ \phi(v) : v \in \mathbb{R}^3 \}$
            \begin{itemize}
                \item Brute force:
                \begin{flalign*}
                    \left[ \begin{array}{ccc}
                        4 & -2 & 0 \\
                        -2 & 1 & 0 \\
                        0 & 0 & 5
                    \end{array} \right] \left[ \begin{array}{c}
                        x \\ y \\ z
                    \end{array} \right] &= \left[ \begin{array}{c}
                        4x-2y \\ -2x+y \\ 5z
                    \end{array} \right] \\
                    &= x \left[ \begin{array}{c}
                        4 \\ -2 \\ 0
                    \end{array} \right] + y \left[ \begin{array}{c}
                        -2 \\ 1 \\ 0
                    \end{array} \right] + z \left[ \begin{array}{c}
                        0 \\ 0 \\ 1
                    \end{array} \right] \\
                    &= -2x \left[ \begin{array}{c}
                        -2 \\ 1 \\ 0
                    \end{array} \right] + y \left[ \begin{array}{c}
                        -2 \\ 1 \\ 0
                    \end{array} \right] + z \left[ \begin{array}{c}
                        0 \\ 0 \\ 1
                    \end{array} \right] \\
                    &= (y-2x) \left[ \begin{array}{c}
                        -2 \\ 1 \\ 0
                    \end{array} \right] + z \left[ \begin{array}{c}
                        0 \\ 0 \\ 1
                    \end{array} \right] \\
                    \Rightarrow B_I &= \left\{ \left[ \begin{array}{c}
                        -2 \\ 1 \\ 0
                    \end{array} \right], \left[ \begin{array}{c}
                        0 \\ 0 \\ 1
                    \end{array} \right] \right\}
                \end{flalign*}
                \item Orthogonal decomposition:
                \item Since $\phi(v) = 5v - (a \cdot v)a$ and $\vert \vert a \vert \vert^2 = 5$, we can decompose:
                \[
                v = v_\perp + v_{\vert \vert}, \text{ where } v_{\vert \vert} = \frac{(a \cdot v)}{\vert \vert a \vert \vert^2} a = \frac{(a \cdot v)}{5} a
                \]
                \item Then: $\phi(v) = 5 v_\perp + 5 v_{\vert \vert} - 5v_{\vert \vert} = 5 v_\perp$
                \item So $\phi$ scales the component orthogonal to $a$ by 5:
                \[
                im(\phi) = \{ w \in \mathbb{R}^3 : a \cdot w = 0 \}
                \]
                \item Solving $w_1+2w_2=0$ gives $w_1=-2w_2$. A basis is:
                \[
                im(\phi) = span \left\{ (-2,1,0)^\top, (0,0,1)^\top \right\}
                \]
            \end{itemize}
        \end{itemize}
    \end{itemize}
    \item Question 3: Let $\mathbb{R}[x]_4$ denote the real vector space of polynomials in $x$ of degree at most 4. Consider the linear map
    \[
    \phi : \mathbb{R}[x]_4 \rightarrow \mathbb{R}^3, \qquad \phi(f) = \left[ \begin{array}{c}
        f(0) \\ f'(0) \\ f''(0)
    \end{array} \right]
    \]
    \begin{itemize}
        \item (a) Show that $\phi$ is surjective
        \begin{itemize}
            \item For $f(x) = a_0 + a_1x + a_2x^2 + a_3x^3 + a_4x^4 \in \mathbb{R}[x]_4$
            \begin{flalign*}
                f(x) &= a_0 + a_1x + a_2x^2 + a_3x^3 + a_4x^4 \\
                f'(x) &= a_1 + 2a_2x + 3a_3x^2 + 4a_4x^3 \\
                f''(x) &= 2a_2 + 6a_3x + 12a_4x^2 \\
            \end{flalign*}
            \item At $x=0: f(0) = a_0, f'(0)=a_1, f''(0)=2a_2$
            \item Thus $\phi(f) = \left[ \begin{array}{c}
                a_0 \\ a_1 \\ 2a_2
            \end{array} \right]$
            \item For any $(y_1,y_2,y_3)^T \in \mathbb{R}^3$, choose:
            \[
            a_0 = y_1, \quad a_1=y_2, \quad a_2 = \frac{y_3}{2}, \quad a_3,a_4 \text{ arbitrary}
            \]
            \item Then $\phi(f) = \left[ \begin{array}{c}
                y_1 \\ y_2 \\ y_3
            \end{array} \right]$
            \item Therefore $\phi$ is surjective
        \end{itemize}
        \item (b) Find a basis for $ker(\phi)$
        \begin{itemize}
            \item We find all $f$ such that $\phi(f) = (0,0,0)^\top$.
            \item From $\phi(f) = (a_0,a_1,2a_2)^\top = (0,0,0)^\top$:
            \[
            a_0 = 0, a_1 = 0, a_2 = 0.
            \]
            \item Thus any $f \in ker(\phi)$ has the form $f(x) = a_3x^3 + a_4x^4$.
            \item Therefore:
            \[
            ker(\phi) = span\{x^3,x^4\}
            \]
            \item Since $dim\mathbb{R}[x]_4 = 5$ and $rank(\phi)-3$, we have $dim \ ker(\phi)=2$, confirming our result.
        \end{itemize}
    \end{itemize}
\end{itemize}

\section{Tutorial 10 (11/12/2025)}
\begin{itemize}
    \item Exercise 7.4. Find the equations of the tangent plane and the normal line to the graph of the specified function at the specified point.
    \begin{itemize}
        \item (a) $f(x,y) = x^2-y^2$ at $(-2,1)$
        \begin{itemize}
            \item Let $z=f(x,y)$. Then, the graph of $f$ can be written as $F(x,y,z)=0$, where $F(x,y,z) = f(x,y)-z$. Note that $f(-2,1)=3$.
            \item The gradient of $F$ is $\nabla F(x,y,z) = (2x,-2y,-1)$, which is $(-4,-2,-1)$ at $(x,y,z)=(-2,1,3)$. This is normal to the graph of $f$.
            \item So, the normal line is
            \[
            r = (-2,1,3) + \lambda(-4,-2,-1).
            \]
            \item The tangent plane is then
            \[
            4x+2y+z = (-2,1,3)^\top \cdot (4,2,1)^\top = -3.
            \]
        \end{itemize}
        \item $f(x,y)=e^{xy}$ at $(2,0)$.
        \begin{itemize}
            \item Again, define $F(x,y,z) = f(x,y) - z$. The graph is then $F(x,y,z)=0$. Note that $f(2,0)=1$.
            \item This time, $\nabla F = (ye^{xy},xe^{xy},-1)$ is the normal direction. So, at $(2,0,1)$, the normal direction is $(0,2,-1)$.
            \item So, the normal line is
            \[
            r = (2,0,1) + \lambda(0,2,-1).
            \]
            \item The tangent plane is then
            \[
            2y-z = (2,0,1)^\top \cdot (0,2,-1)^\top = -1.
            \]
        \end{itemize}
    \end{itemize}
\end{itemize}

\end{document}
